% ============================================================
%  Harald Høffding, "Spinoza's Ethica. Analyse og Karakteristik"
%  "Spinoza's Ethica: Analysis and Characterization"
%  D. Kgl. Danske Vidensk. Selsk. Skrifter, 7. Række, historisk og
%  filosofisk Afd. III, 3.  København: Høst & Søn, 1918.  97 pp.
%
%  TRANSLATION (English), from transcription.tex (Danish) in this directory,
%  the text of record (collated 2026-10-04; see its header).  Translated
%  2026-10-04 from that file as it stood after the odd spots found in
%  translating were read at the image and fixed in it.  Assembled from
%  .parts/tr/pNN.tex (gitignored working files) by .parts/tr/build.sh; this file
%  is the result and is edited directly from now on.
%
%  BOOK-SPECIFIC CONVENTIONS (beyond TRANSLATION-PLAYBOOK.md §2):
%   * Høffding's terms: Substans = substance; Attribut = attribute; Modus =
%     mode; Emne = thing (det enkelte Emne = the individual thing; Høffding
%     glosses Modus as "Enkeltvæsen, Emne"); Enkeltvæsen = individual being;
%     Grundvæsen = fundamental being; Grundegenskab = fundamental property;
%     Følelse = feeling (Spinoza's affectus); Drift = desire (cupiditas); Trang
%     = urge (appetitus; Selvopholdelsestrang = urge of self-preservation);
%     Forestilling = idea; Erkendelse = knowledge; Forstaaelse = understanding;
%     Virkekraft = power of action (potentia agendi); Iboen = immanence, iboende
%     Aarsag = immanent cause; Lovmæssighed = conformity to law; Sjæl = soul;
%     Rækkedannelse = series formation; Værdi = value; Glæde/Sorg = joy/sorrow
%     (Ulyst = pain); godt og ondt = good and evil; Aandskraft = strength of
%     mind (fortitudo), Livsmod = courage (animositas), Højmod = generosity
%     (generositas), Højsind = magnanimity; Erkendelsestrang = urge for
%     knowledge; Motivforskydning = displacement of motive; Tilskyndelse =
%     impulse; aandelig = spiritual.
%   * "Ethica" stays as the title of the book, roman, as printed.
%   * The paragraph labels as printed: a. b. c. in the Introduction and the
%     Concluding Characterization, α. β. γ. δ. in the five Books.
%   * Names in their English forms (Platon -> Plato, Aristoteles -> Aristotle,
%     Baco -> Bacon, Galilei -> Galileo, Thomas af Aqvino -> Thomas Aquinas,
%     Hugo de St. Victore -> Hugh of Saint Victor), in small capitals.
%   * Names in small capitals (\textsc) as in the Danish.  Titles in the notes
%     (italic in the Danish) and Latin terms (roman in the Danish) -> \textit.
%   * Quotations in Latin, German, French, English, Italian and Greek stay in
%     the original.  Spinoza quoted by Høffding in Danish is translated from
%     his Danish.  Titles cited in the notes stay as cited (Høffding's own
%     works in Danish); "Smlgn." / "Smlg." / "Se" -> "Cf." / "See"; citation
%     numbers as printed.
%   * Footnotes are numbered through the book (the print numbers them per page).
%   * Misprints: translate the sense, with a comment at the site (the Danish
%     marks them PRINTED AS IS).
%   * Page markers \opage{N} placed by hand at the English word where printed
%     page N begins (the book's own pagination, as in the transcription).
% ============================================================
\documentclass[12pt]{article}
\usepackage[utf8]{inputenc}
\usepackage[T1]{fontenc}
\usepackage[english]{babel}
\usepackage[protrusion=true,expansion=true]{microtype}
\usepackage[margin=1.3in]{geometry}
\usepackage{setspace}
\usepackage{enumitem}
\usepackage{libertinus}
\usepackage{libertinust1math}
\usepackage{textalpha}
\usepackage{fancyhdr}
\usepackage[colorlinks=true,linkcolor=black,urlcolor=black,
  pdftitle={Spinoza's Ethica: Analysis and Characterization},
  pdfauthor={Harald Høffding}]{hyperref}
\setlength{\headheight}{15pt}
\pagestyle{fancy}
\fancyhf{}
\fancyhead[LE,RO]{\thepage}
\fancyhead[C]{\textit{Harald Høffding: Spinoza's Ethica (1918)}}
\setstretch{1.3}
\usepackage{marginnote}
\usepackage{graphicx}
\newcommand{\opage}[1]{\marginnote{\footnotesize\textit{[#1]}}}
\newcommand{\tocline}[2]{\noindent #1\dotfill\ #2\par}

\title{Spinoza's Ethica\\[0.4em]\large Analysis and Characterization}
\author{Harald Høffding}
\date{Translated from the Danish\\[0.6em]
  \small \textit{Spinoza's Ethica. Analyse og Karakteristik.} D.\ Kgl.\ Danske
  Vidensk.\ Selsk.\ Skrifter, 7.\ Række, historisk og filosofisk Afd.\ III, 3.
  Copenhagen: Andr.\ Fred.\ Høst \& Søn; Bianco Lunos Bogtrykkeri, 1918}

\begin{document}
\maketitle
\thispagestyle{empty}
\clearpage

% ===== pp. 3--8 =====
\opage{3}\section*{Introduction.}

1. A philosophical system is an attempt to lead what is given in experience,
the things before us, back to one fundamental thing, a primal phenomenon, and
thereby to gather all knowledge into one fundamental thought. The validity of
this fundamental thought can be sought to be shown partly through an analysis
of the things at the philosopher's disposal, partly through an attempt to derive
these things from the fundamental thing. By the first procedure the fundamental
thing becomes a consequence of the individual things; by the second these become
consequences of it. In no system are attempts to apply both procedures lacking.
But even if such attempts could succeed, the question would still remain whether
the circle of things the thinker commands is not limited; and since new things
can always turn up, this question must be answered yes. No system can therefore
become anything but a hypothesis built on the analysis of a more or less
comprehensive circle of things. As for the second procedure, difficulties of
principle will also arise, since it already proves impossible on purely logical
grounds to derive a manifold of things from a single fundamental thing; every
inference needs, after all, more than one premise.

This consideration may well be said to have penetrated the philosophical world
to such a degree that the time of systems is now past. In a certain sense there
will of course always be ``system'' in serious thought, since every single
thought must be able to ``stand together'' with the other thoughts that are
held fast. But ``system'' in the sense of a world-explanation or world-view that
is not merely rounded off but closed and definitive belongs to the past.

The study of the philosophical systems of the past has, however, more than a
purely historical interest. Such edifices of thought were natural in themselves
so long as no thorough critical examination of human knowledge had been
undertaken. And on closer consideration it often turns out that problems we are
inclined to regard as wholly modern lie hidden as elements in the spontaneous
crystallizations that took place in the thinkers of the past. If now the things
with which thought is occupied keep their significance, --- if the fundamental
thoughts with which it works are closely bound up with the nature of human
thought, --- and if the work is done in earnest and with a sense for the depth of
the problems, --- then such a \opage{4}system can not merely stand as the
expression of the thought of a bygone time, but can also serve as a model of how
one must think under other historical presuppositions.

Of few systems does this hold to the degree that it holds of \textsc{Spinoza}'s.

Nor did \textsc{Spinoza}, for that matter, think he had set up a complete system,
closed for all time, as a frequent but mistaken view of him supposes. In a letter
to \textsc{Oldenburg} (Ep. 30 in \textsc{van Vloten}'s edition) he says that he does
not know how things really hang together, and how every part fits into the whole
to which it belongs; --- for that, one would need to know the whole of nature and
all its parts! In a letter to \textsc{Boxel} (Ep. 56) he says that we do not know
the greater part of nature's fundamental properties (attributes); we know only the
two fundamental properties that make themselves known in our own being, thought
and extension (Ep. 64; cf. Ethica II, Ax. 5). And when a young friend of his had
gone over to Catholicism and now, in a letter in which he tried to convert him,
asked mockingly whether he thought he had found the best philosophy of all,
\textsc{Spinoza} answered (Ep. 76) that he by no means thought he had arrived at the
best philosophy, but that he was certain he knew the true philosophy. Not only with
regard to nature's fundamental properties (attributes), but also with regard to the
existence of the individual, determinate things (\textit{existentia modorum}) we
stand on the ground of mere experience (Ep. 10); and as for the understanding of
these individual things given in experience, we stand before an infinite task,
since the individual thing (the individual being) is understood through its
connection with other individual things, so that the series of causes goes on to
infinity (Ethica I, 28). Knowledge from experience can thus never be closed.

How far the statements brought forward here can be brought into strictly necessary
connection with the whole of \textsc{Spinoza}'s other views will appear when we go
through the several parts of the principal work. But before we go on to such an
inquiry, we must try to find traces of the work of thought from which the system
arose. If the finished system can be compared with a crystal, it becomes
intelligible only when one can trace the formation of the crystal, that is, find
the presuppositions that were decisive for the building of the system.

The students of Spinoza have, especially in the most recent times, laid the main
weight on finding historical predecessors of \textsc{Spinoza}'s peculiar line of
thought. For all the learning displayed in this, one cannot say that it has got
further than mere possibilities. \textsc{Spinoza}'s first biographers
(\textsc{Lucas}, \textsc{Colerus}) pointed to \textsc{Descartes}. Modern
scholars have examined his relation to Jewish medieval theology (\textsc{Joel}),
to the new scholasticism (\textsc{Freudenthal}), to \textsc{Giordano Bruno}
(\textsc{Sigwart}, \textsc{Avenarius}), to the freer movements in the Dutch
intellectual life of the time (\textsc{Meinsma}), to the theologians of the past
in general
(\textsc{Dunin-Borkowski})\footnote{It has been \textsc{Spinoza}'s fate that his
personality and his teaching have repeatedly been defended by men who stood far
from, or were even hostile to, his whole direction. Thus the pastor
\textsc{Colerus}, who was himself much offended by \textsc{Spinoza}'s teaching,
cleared his memory of various slanders, and has therefore been called a modern
Balaam. Such a Balaam was \textsc{Jacobi} too, whose theology of feeling made him
look on \textsc{Spinoza}'s philosophy with horror. And in our own day the Jesuit
\textsc{Dunin-Borkowski} has studied \textsc{Spinoza}'s development not only with
great learning but also with great sympathy. --- The sympathy has its ground in
his seeing \textsc{Spinoza} as a counterweight to the scepticism and dilettante
free-thinking of that age and of ours. But compared with what
\textsc{Dunin-Borkowski} calls the world-philosophy, namely the medieval
scholasticism with \textsc{Thomas Aquinas} at its head, \textsc{Spinoza}'s
philosophy nevertheless stands as a great illusion. (\textsc{Stanislaus v.
Dunin-Borkowski}: \textit{Der junge Spinoza}. Münster 1910, p. 351---485).}. It
has long since been established that \opage{5}\textsc{Spinoza} did not spin his
wisdom out of himself, and that on the whole he was not the hermit his first
disciples found it most advisable to make him appear. He was a well-read man in
many fields and was able to take from many sides what he could use. But the
decisive thing was a current of thought that arose in him spontaneously and, at
first, without clear awareness, and that determined the choice of what he took
up from the thoughts of his time, and also decided what place they got in the
edifice of his thought. At bottom, a spiritual need that arose in him early was
the source of the current of thought and determined its direction as well.
% "hos kam" in the Danish is read as "hos ham" (him)

2. The work of thought that began early in \textsc{Spinoza} started within Jewish
theology and was taken up again on the basis of the Western natural science and
philosophy into which he had, especially after his expulsion from the synagogue,
lived himself. In both fields he went the way of analysis and consequence and
sought to hold fast to and carry through presuppositions of which his
predecessors had not been clearly aware. And throughout all this work he was led
by an urge to understand himself and his connection with the great, lawful and
infinite whole of existence, an urge that in him had at once an intellectual and
a religious character. These three elements or tendencies I will characterize
briefly, before going on to show, through an analysis of the principal work
(Ethica), how they are refracted in his thought. In the definitions and axioms
with which the first book of the Ethica begins, thoughts have crystallized that
can be traced back to the three sources indicated. The peculiarity of the system
and the personality of its author show themselves in the way in which the three
motives work together and strive to come into their own.

\textit{a.} \textsc{Spinoza}'s youthful work, the ``Short Treatise'', has a more
theological and mystical stamp than the principal work, the Ethica. There
\textsc{Spinoza} starts from the concept of God as it had developed in the higher
popular religions, and he draws all the consequences of God's being an
unconditioned and infinite being, --- a being without presuppositions and without
limits. It is by this way that he arrives at the strict concept of substance put
forward in the first definitions of the Ethica, --- the concept of that which
exists through itself and is understood through itself, that is, whose concept
does not presuppose any other concept from which it would have to be formed. Here
an urge shows itself to reach something last, something without presuppositions,
that does not in turn send thought beyond itself. \textsc{Spinoza} meets the
objection that every definition after all consists precisely in bringing a
concept under a higher concept (K. Tr. I, 7, 9---10). That does not hold, he
maintains, of the fundamental concepts that condition all knowledge;
\opage{6}if they too presupposed other concepts, nothing could ever be fully known!
A peculiar union of mysticism and theory of knowledge appears here. The urge to
find concluding concepts is not purely intellectual; it is one with the urge to
come to rest. If we are not to be thrown out into endless series of thoughts, we
must hold to the thought of a single substance, a single fundamental being (K.
Tr. I, 2, 10). Our thought finds rest neither by stopping at the knowledge of the
spiritual nor at the knowledge of the corporeal, but must go back to that which
conditions the existence of these two domains and makes the understanding of them
possible: God or substance (K. Tr. II, 22, 5).

In Jewish, in Mohammedan and in Christian theology alike, the concept of God has a
theoretical significance: it supplies the principle of understanding, in any case
of the concluding understanding of everything. And it is about this side of the
concept of God that \textsc{Spinoza}'s thinking concentrates more and more. His
relation to theology may be described by saying that he takes the question
seriously of what character the concluding thought would have to have. It cannot
be unintelligible, and if it is not to lead beyond itself, it must be understood
through itself. To make clear to ourselves what \textsc{Spinoza}'s substance must
mean, we must therefore ask what it is that is understood through itself and can
make everything else intelligible. It can only be the necessary connection between
subject and predicate, or between premises and conclusion, the connection that
underlies every transition of thought. Further back than to the principle of all
thinking, and thereby also of all the rationality we can find in existence, we
cannot go. We can give no reason why the conclusion follows from its premises. It
is by an analytic intuition\footnote{On this concept see my paper on intuition in
Vid. Selsk. Oversigter 1914, p. 77---80.} that we bring out the principle of all
rationality. The principle of all understanding stands in \textsc{Spinoza}'s
metaphysics as substance as well, the being that underlies everything. His
metaphysics is a projection of his theory of knowledge. He begins theologically by
saying that God is what makes everything intelligible to us, and he ends by saying
that what makes everything intelligible to us is God. Whereas in the Short Treatise
he goes from the concept of God to the concept of substance, in the first
definitions of the first book of the Ethica he goes, conversely, from the concept
of substance to the concept of God. This reversal is prepared in the first
appendix to the Short Treatise, and it appears decisively in the first letters
(see especially Ep. 2).

The concept of God that \textsc{Spinoza} took with him from his theological period
and always held fast to in his principal work was of course very different from
the concept of God of the popular religions, and he himself calls attention to
the difference (see e.g. Ep. 19. 21. 73). He constructs his concept of God from
the presuppositions of all rational understanding of existence, but must then
(continuing what Jewish philosophers of religion in the Middle Ages, especially
\textsc{Maimonides} and \textsc{Ibn Gebirol}, had begun) seek to purge all
anthropomorphisms, all concrete psychological properties, from the content of
this concept.

In the same years in which \textsc{Spinoza}'s break with the popular religion in
which he had grown up was being prepared by the independence of his thought,
another great \opage{7}mind, \textsc{Blaise Pascal}, returned through a religious
crisis to the concept of God of his popular religion, with an express affirmation
of all anthropomorphisms. ``God of Abraham, of Isaac and of Jacob, God of Jesus
Christ, not of the philosophers and the men of science!'' he cries in the
passionate utterance that he wrote down on 25 November 1654 and thereafter always
carried on him, sewn into his clothing.

\textsc{Spinoza}'s God was precisely ``the God of the philosophers and the men of
science'', in so far as his concept of God is a metaphysical expression of the
last presupposition of all thought and science. Now since this presupposition does
not stand in an external relation to the individual things, but makes itself known
in our understanding of them and in every fixing of the conditions of their
existence, \textsc{Spinoza}'s God must be an indwelling (immanent) God, not an
external (transcendent) God. \textsc{Spinoza} describes his opposition to the
usual concept of God thus (Ep. 73): ``I hold an opinion of God and nature wholly
different from the one modern Christians embrace. For I maintain that God is the
indwelling (immanent), not the external (transeunt) (= transcendent) cause of all
things.'' In the Short Treatise (first dialogue) he had illustrated this point by
a comparison with the understanding, which is active in all individual thoughts
but is not the external cause of them\footnote{\textsc{Hume}, who cannot have
known the Short Treatise, and who probably knows \textsc{Spinoza} only from
\textsc{Bayle}'s Dictionnaire, uses \textsc{Spinoza}'s teaching to frighten his
spiritualist opponents. When these maintained that our individual sensations and
ideas were all members of one and the same consciousness, and did not, as
\textsc{Hume} thought, stand in a purely external, mechanical relation of
association to one another, then they were using, he says, the same line of
thought (the concept of immanence) that led \textsc{Spinoza} to his impious
opinions! (\textit{Treatise on Human Nature}. I, 4, 5).}. --- \textsc{Spinoza}'s
inner or indwelling cause corresponds to what \textsc{Bruno} calls ``principle'',
whereas \textsc{Bruno} by ``cause'' means what \textsc{Spinoza} calls external
cause\footnote{\textsc{Bruno}: \textit{De la causa, principio et uno}. Opere ed.
Lagarde, p. 230, 272.}. ---

From God's being the principle of all understanding it follows not only that he
must be the ``indwelling'', not the ``external cause'', since he relates to all
things as the ground relates to the consequence. It follows as well that no limit
can be set to God's being. All the fundamental properties of existence must be
God's properties. It is through this line of thought that for \textsc{Spinoza}
the concept of God is not only one with the concept of substance, as we have
already seen, but also one with the concept of nature, so that he can use the
three expressions God, substance and nature as synonymous. In the first dialogue,
which seems to be the oldest part of the Short Treatise, it is asked whether ``a
being that is most perfect, since it cannot be limited by any other being'' can
be thought, and the answer is that nature is such a being, since everything is
comprised in it. And it is shown besides that this infinity is closely bound up
with the concept of the indwelling cause. If the things (the things in the world)
were due to an external cause distinct from them, this cause would be limited.
--- There is thus an inner connection between the principle of understanding,
immanence and infinity, and it is \textsc{Spinoza}'s consistent thinking through
of the traditional concept of God that leads him to all three concepts.
% the Danish has „iboende“, ikke ydre Aarsag“ (the opening „ before ydre lost; see the odd-spots list)

\opage{8}\textsc{Avenarius}\footnote{\textsc{Rich. Avenarius}: \textit{Die beiden
ersten Phasen des Spinozischen Pantheismus}. (1868).} and
\textsc{Sigwart}\footnote{\textsc{Chr. Sigwart}: \textit{Spinoza's kurzer
Tractat\textsuperscript{2}}. Prolegomena. (1869).} suppose that the concept of
nature became a member of \textsc{Spinoza}'s thought through the influence of the
thinkers of the Renaissance, perhaps of \textsc{Bruno}. This is most probable, and
on closer consideration it will lead us over to the second great motive of thought
in \textsc{Spinoza}. % the Danish misprints Indfiydelse (influence)
But as the concept of nature turns up in the first dialogue, in the answer to the
question whether there is a being that is not limited by anything else, it is
presupposed that an urge has already stirred beforehand to close thought in a
concept that in turn leads beyond itself no longer. That is \textsc{Spinoza}'s
fundamental urge, and it developed while he thought through the traditional
concept of God. It is from the thought of God as the principle of all
understanding that he came to the thought of God's immanence and of the infinity
of God or nature.

The line of thought that leads \textsc{Spinoza} to maintain the validity of the
concept of infinity, both as regards substance (as absolutely infinite) and as
regards its attributes (as each infinite in its kind), may recall the line of
thought that led \textsc{Bruno} to reject the validity of the old, limited picture
of the world. The first two propositions of the second book of the Ethica indicate
the way \textsc{Spinoza} went, the way that led him to the second and sixth
definitions of the first book. We can, it says (Eth. II, 1. Schol.), always add
thought to thought without there being any limit here, and likewise we can (Eth.
II. 2) add extension to extension without limit. The world of thought and the world
of extension are each infinite within their kind (\textit{in suo genere}), since
within them one can always advance further by the same law. And it is the same
with the concept that underlies all concepts, namely the concept of something that
can appear in infinitely many forms or kinds, since the series of these different
kinds (\textit{genera}) or attributes can always be continued, just like the
series of the different things within the same kind. Hence the definition of God
at the beginning of the Ethica (Eth. I, Def. 6): ``By God I understand an
absolutely infinite being, that is, a substance consisting of infinitely many
attributes, each of which expresses something eternal and infinite.'' --- In Ep.
36 this whole line of thought is clearly indicated.

In \textsc{Bruno} too what we may call the law of continuous series formation is
applied. % the Danish misprints BBUNO
Sensation, imagination and thought alike show us, he says, that we can always add
member to member. The sensible horizon can always be widened, and we can add image
to image, magnitude to magnitude, number to number, thought to thought to infinity.
Therefore he asks whether we can think the world limited: at the very limit we try
to set to the world, a new horizon would form around us, if we imagine ourselves
transported there\footnote{See on this \textit{Den nyere Filosofis
Historie\textsuperscript{2}}, I, p. 118---123. --- This thought comes out
especially in the book \textit{De Immenso} (Francofurti 1591). This work also
presents the relation between freedom and necessity in the same way as
\textsc{Spinoza}. (Eth. I, Def. 7).}. The logical genesis of the concept of
infinity proves to be the same in the two thinkers, and \textsc{Bruno} is the
forerunner here.

% ===== pp. 9--13 =====
The difference \opage{9}between their use of this concept is bound up with the
difference between their problems. In \textsc{Bruno} the widening of the picture
of the world is the main thing, --- in \textsc{Spinoza} it is the carrying through
of conformity to law within the picture of the world. But there is here a
possibility of an influence of \textsc{Bruno} on \textsc{Spinoza} that concerns a
more essential point than those to which one has usually kept in discussing the
relation between the two philosophers.

In both of them, of course, the question remains by what right we ascribe real
validity to the psychological or logical possibility\footnote{It is not clear
what principle of series formation \textsc{Spinoza} has in mind when he says that
merely by keeping to the attribute of thought we arrive at the thought of an
infinite being (Eth. II, 1. Schol.). It could be a purely psychological series of
sensations (Eth. II, 18. Schol.), or (more probably) a series of grounds and
consequences (Eth. I, 28), or (improbably) a series of members each of which
relates to the preceding as subject to object (Eth. II, 21. Schol.). --- Cf. on
various kinds of series formation \textit{Den menneskelige Tanke}, p. 164---172,
and specially on series formation by virtue of the principle of repetition
\textit{Totalitet som Kategori} (Vid. Selsk. Skr. 1917), p. 34---37.} of continued
series formation. This question first comes forward in the critical philosophy.
If we compare \textsc{Bruno} and \textsc{Spinoza} in this respect, \textsc{Spinoza}
stands more dogmatically than his predecessor. Whereas for \textsc{Spinoza} the fact
that we can construct a concept of infinity is proof enough that there is an
infinite being, \textsc{Bruno} distinguishes between construction in thought and
reality, and specially between mathematics and physics\footnote{E.g. \textit{La
cena de le ceneri} (Opere ital. ed. Lagarde, p. 184): Altro è giocare con la
geometria, altro è verificare con la natura.}.

\textit{b.} After his departure from the synagogue \textsc{Spinoza} became
acquainted with the science, philosophy and literature of the West. And when now,
having set up the concept of God as the principle of all understanding, he asks in
turn what understanding really means, he maintains that we learn this by keeping
to the fundamental thoughts (\textit{notiones communes}) that prove to apply
everywhere in nature, in the part as in the whole, that is, to the universal laws
of everything existing (Tract. theol. polit. Ch. 6---7). Only by this way can we
come to know the divine. ``As for me,'' \textsc{Spinoza} says in a letter (Ep. 21),
``I have not learned, and have not been able to learn, any of God's eternal
attributes from Holy Scripture.'' --- Theology had taught him that God is the
principle of intelligibility, and he now turns to the presuppositions of science
to learn what understanding is.

Not enough attention has been paid, I think, to what a comparison between the
Short Treatise and the Ethica clearly shows, that principles of natural science
form the background of \textsc{Spinoza}'s philosophy in its definitive form.
% the Danish misprints Sammenligniug (comparison)
When he demands (Eth. I, 28) that every phenomenon (thing) be explained by another,
this by a third, and so on, this contains the requirement set up by
\textsc{Kepler} (which \textsc{Kepler}\footnote{\textit{Den nyere Filosofis
Historie\textsuperscript{2}}, I, p. 165. --- There was a work of \textsc{Kepler}'s
in \textsc{Spinoza}'s library, but it was a chronological work, not one of his
purely scientific ones. Of mathematical and scientific works \textsc{Spinoza} owned
books by \textsc{Vieta}, \textsc{Huyghens}, \textsc{Riolan}, \textsc{Steno} and
\textsc{Boyle}, besides \textsc{Descartes}.} himself applied in criticizing the
speculations of his youth), that the causes that can be recognized in science must
themselves be capable of being shown in \opage{10}experience. For him this
principle coincides with the principle of inertia set up by \textsc{Galileo} and
\textsc{Descartes}, which he sets up in Eth. II, Lemma 3 as valid for material
things and also applies (Eth. III, 6; cf. II, 1. Schol.) to spiritual things. To
this is joined the proposition of the constant relation between rest and motion,
another expression of \textsc{Descartes}'s proposition of the conservation of
motion and, like it, an imperfect forerunner of the modern proposition of the
conservation of energy (Eth. II, Lemma 7, Schol.; cf. Ep. 32. 64). --- In the
``geometrical'' presentation that \textsc{Spinoza} unfortunately gave of his
teaching, the propositions mentioned are derived from the concept of attribute,
which (Eth. I, def. 4) is defined as a property of substance that is understood
through itself and therefore is not understood by means of any other property. It
follows, among other things, that an extended phenomenon must be explained by
another extended phenomenon. But in reality \textsc{Spinoza} arrived at his strict
concept of attribute through analysis of the fundamental concepts of natural
science. What is put first in the synthetic presentation came last in the
analysis. That extension cannot be limited by thought, or thought by extension,
that one must not, therefore, in one's explanation of nature jump over from one
attribute to another (Eth. I, def. 2; I, 10 and II, 6), --- this is in the end a
principle of method found by analysis of the scientific propositions mentioned,
although \textsc{Spinoza} puts it forward as an eternal truth.

One of \textsc{Spinoza}'s correspondents objected to \textsc{Spinoza}'s definition
of attribute that it was not, after all, beyond doubt that thought could not be
limited by body or body by thought, since thought might be a bodily motion
(\textit{motus corporeus}) (Ep. 3). \textsc{Spinoza} answers (Ep. 4) that
extension, considered as extension, is not thought, and that this is all he needs
in order to hold fast his proposition that thought and extension cannot limit each
other. --- The principle of series formation mentioned above clearly underlies
this. \textsc{Spinoza} demands series of homogeneous members. And this brings us
over to the most fundamental proposition in \textsc{Spinoza}'s system, the
proposition of causality, and to the way in which \textsc{Spinoza} conceives it.

The causal relation is intelligible to us, according to \textsc{Spinoza}, only when
its members, what we call ``cause'' and what we call ``effect'', are of the same
nature, so that a relation of identity holds between them. In the treatise on ``the
improvement of the understanding'' he says that the knowledge of the effect is
nothing but a perfect understanding of the cause, and presupposes it. ``If two
things have no property in common, they cannot be understood through each other,
or the concept of the one does not contain that of the other'' --- so say the
fourth and fifth axioms of the first book of the Ethica. And in Ep. 4 it is stated
that if two such things were to stand in a causal relation to each other, there
would be something in the effect that is not in the cause, and this something
would then have its ground in nothing.

\textsc{Spinoza} thus demands equivalence between cause and effect, just as in
inferences that can be converted there is equivalence between premises and
conclusion. Within the world of phenomena (things) he wants to find a connection as
strict and logical as that between the members of a series of logical inferences
that can \opage{11}be converted. One is not to begin with abstractions, but with
thoughts that make the origin and succession of things, spiritual and material
alike, intelligible. The more we understand individual things, the more we
understand God (Eth. V, 24), for what makes them intelligible to us is precisely
God. What exists through itself and is understood through itself is not itself an
individual thing, an individual being. It is the principle of grounding, or of
conformity to law, which is the last presupposition of all science and which
proves itself in the whole as in all the individual parts of existence. For
\textsc{Spinoza} the knowledge of God and natural science were so far from excluding
each other that, on the contrary, they coincided. To think and to inquire was for
him a religious act.

In a famous work that appeared some ten years after \textsc{Spinoza}'s Ethica, a
similar line of thought is in part found. In his ``Principia'' (and later in the
``Optics'') \textsc{Newton} conceived space as God's sense organ
(\textit{sensorium}), the organ of omnipresence. Mathematical inquiry then also
became knowledge of God. Both \textsc{Newton}'s and \textsc{Spinoza}'s thought, if
followed historically, would lead back to \textsc{Plato}. In other respects
\textsc{Spinoza}'s and \textsc{Newton}'s philosophy of religion are very different.
\textsc{Newton} wants to ground the assumption of God's existence essentially in a
teleological way, from the purposiveness in the structure of the world. But for
\textsc{Spinoza}, to think of God as acting according to purposes was one with
thinking of him as a finite and imperfect being, which does not yet possess all
that accords with its nature. The sharpest statements found in \textsc{Spinoza}
come precisely in his criticism of teleology.

Whereas in the foregoing weight has been laid on \textsc{Spinoza}'s philosophy,
especially his concluding conception of the concepts of substance and attribute,
being due to an analysis of the causal explanation of empirical science, an acute
student of Spinoza has on the contrary objected that \textsc{Spinoza}'s concept of
substance contained no causal determinations, and that no such determinations can
be derived from his definition of substance either (Eth. I, def.
3)\footnote{\textsc{Gabriel Huan}: \textit{Le dieu de Spinoza}. Paris 1914, p. 44
f.}. This remark, correct in itself, contains a criticism of \textsc{Spinoza}, not
a criticism of my view of him. Since \textsc{Spinoza} in fact conceives causality
as a purely logical relation, he needed to take no other than purely logical
determinations into his definition of substance. \textsc{Spinoza}'s philosophy is an
attempt to think the world by the formal categories, to which he led back all real
categories (such as causality) and all ideal categories (the concepts of
value)\footnote{On the different groups of categories see \textit{Den menneskelige
Tanke}.}. One can criticize him partly by showing that he did not in reality carry
through the reduction of all categories to the formal ones, partly by establishing
the impossibility of such a reduction. But as his work actually stands, one will
not be able to deny that it is characterized by an attempt, magnificent in itself,
in the direction named. In the following analysis of the several parts of the work
there will be occasion to throw light on this main point.

\opage{12}\textit{c.} By \textsc{Spinoza}'s own declaration at the beginning of
the unfinished treatise ``on the improvement of the understanding'', his personal
experience of life was of decisive significance for him in the development of his
main thoughts. When he had to account to himself for what significance
philosophical thought had for him, or rather, how philosophical thought had become
for him that in which he at last found the value of life, the explanation lay for
him in this: that the only thing that is in our power and makes us independent of
external circumstances is the work of thought itself, and the knowledge that can be
won through it of the connection of our mind with the whole of nature.
% the Danish misprints Tanksarbejdet (the work of thought)
Other goods --- pleasures, honour, riches --- rest on something that does not spring
from our innermost being. But through the work of thought a joy of knowledge and a
love of the divine are reached that lift one above all sorrow and unrest. Here he
saw an aim for his life, an aim that at the same time made possible, and brought
with it, an endeavour that other men may reach the same aim. For, as he later said
(Eth. IV, App. cap. 9), the best way in which a man can show what ability and
spirit he commands is so to influence others that they come to live under the very
rule of true knowledge.

It was not to flee real life that \textsc{Spinoza} chose this way. He was no
pessimist. He would not have sympathized with the young \textsc{Schopenhauer}'s
words to \textsc{Wieland}: ``Life is a precarious business; I mean to spend mine
thinking about it.'' The life of thought was for \textsc{Spinoza} the highest life,
not because it stood in opposition to the life of reality, but precisely because in
its own way it was able to give expression to all the sides of life, just as it was
itself a side or form of life.

Whether \textsc{Spinoza}'s noble and earnest conduct of life was consistent from his
own standpoint was a question much discussed in and after \textsc{Spinoza}'s own
time, whereas nowadays --- with the necessary knowledge of his line of thought ---
one will hardly see any problem in it. There is another question that has more
interest for us, namely, to what degree the after-effect of the positively religious
standpoint had a lasting influence on his mind.

It has certainly not been mere accommodation when he used the words God, nature and
substance as synonymous. In the fundamental concept of his philosophy he saw a
religious and a rational content expressed at once. No opposition made itself felt
for him here. In particular it has not been mere accommodation when he often,
especially in his letters, uses Christian expressions as symbols for thoughts to
which he himself gave a wholly different content. It would certainly have been best
had he followed the rule he himself lays down in one of his letters (Ep. 23), that a
philosopher should preferably avoid using theological expressions; but
\textsc{Wundt} is decidedly right in the remark he makes in his characterization of
\textsc{Spinoza}, that the latter's view was wholly different from that of the
Enlightenment: ``Whereas the Enlightenment deliberately covers over the mystical
depths of religious thought, it is precisely the mystical content of the concept of
God and of religious feeling that \textsc{Spinoza} seeks to translate into the form
of rational knowledge''\footnote{\textit{Ethik\textsuperscript{4}}, II, p. 113.}.
Through the working thought itself, whose \opage{13}significance for the genesis of
the system I seek to stress in this treatise, \textsc{Spinoza} thought he heard the
voice of the divine. In the first four books of the Ethica the mystical character of
the system recedes in comparison with the Short Treatise, but in the last part of
the fifth book it asserts itself very strongly. The mystical element here comes to
push the rational element aside.

In the most recent times \textsc{Dunin-Borkowski} has maintained that
\textsc{Spinoza} was himself caught in an illusion with regard to his ``moral
energy''. ``\textsc{Spinoza}'', he says\footnote{\textit{Der junge Spinoza}, p.
145.}, ``believed that the light that streamed in over his life came from his
self-made philosophical firmament. But this light was in fact due to other
stars\dots{} [those] that had shone upon him in his youth!'' This remark can be
supported by a psychological fact of great significance. The effects on feeling of
an experience or a line of thought can persist long after the original motive has
fallen away; they can also enter as elements into feelings that arise under the
influence of new experiences or new thoughts. In my Psychology I have called this
phenomenon the expansion of feeling. But the learned Jesuit, so sympathetic towards
\textsc{Spinoza}, lays too little weight on the influence of \textsc{Spinoza}'s great
and deep-going work of thought in the last twenty years of his life, a work of
thought in which his whole personality took part, and which therefore led to a
concentration, a total desire, in relation to which other incentives had to pale.
The proposition that a feeling can only be subdued or displaced by another feeling of
opposite direction and greater strength (Eth. IV, 7) \textsc{Spinoza} could learn
from his own experience, as his whole spiritual energy gathered about the urge to
know.

When a seafarer comes up from the southern hemisphere to the northern, it is new
stars he catches sight of and must take his bearings by. The memory of the southern
sky does not therefore fall away, but it no longer has decisive significance.

The three motives of thought that can be traced in \textsc{Spinoza}, and whose
refraction characterizes his system, in our day stand out more or less sharply
against one another and give rise to different problems. This difference between
problems hardly exists for \textsc{Spinoza}. He sought thoughts by whose help the
different problems could find their solution at once. It is the differentiation of
the problems that above all separates us from \textsc{Spinoza}. But since an inner
connection must always be presupposed between the different problems that occupy
human thought, \textsc{Spinoza}'s great attempt of thought will always keep its
significance as orienting, although we could not now raise edifices of thought in the
same style as his.

% ===== pp. 14--20 =====
\opage{14}\section*{First Book.}
\subsection*{a. Definitions and Axioms.}

The definitions and axioms with which the Ethica begins are, as we have seen in
the foregoing, not set up arbitrarily, but won by analysis partly of inherited
religious ideas, partly of the procedure the new science of the time applied.
Definitions, \textsc{Spinoza} says in a letter (Ep. 9), are arbitrary only when they
are set up merely to be tested, not when they are set up because experience leads
us to them. Experience, it is true, he says in a following letter (Ep. 10), cannot
teach us the essence of things, but it can confine our thought to certain essential
sides of things (\textit{mentem nostram determinare, ut circa certas tantum rerum
essentias cogitet}).

On closer consideration we find in the first definitions and axioms of the Ethica a
formulation of the presuppositions on which, according to \textsc{Spinoza}, all
science rests. As we have seen in the Introduction, these presuppositions are not
regarded by \textsc{Spinoza} as purely formal; they are paraphrases of the principle
of sufficient reason, of the strict logical relation between ground and consequence,
but the norm of truth they express is thought of as an existing being.

The first definition starts from there being something whose existence cannot be
denied without contradicting oneself. When this something is called cause of itself
(\textit{causa sui}), one must remember that in the treatise on the improvement of
the understanding he declares this expression to be purely popular. For it says
there: ``If something is in itself, or, as is popularly said (\textit{ut vulgo
dicitur}), is cause of itself (\textit{causa sui}), then it must be understood
through its own essence alone.'' \textsc{Descartes} had already used this
expression, and it had been objected against him (in the first and the fourth set
of objections to his Meditations) that for a being to be able to be the cause of
itself, it would have to be before it existed. It is perhaps for this reason that
\textsc{Spinoza} criticizes the expression in his treatise on the theory of
knowledge. The fact is, however, that as \textsc{Spinoza} uses the word cause, the
temporal relation between cause and effect plays in the end no part. Hence ground
(\textit{ratio}) and cause (\textit{causa}) become for him one and the same. Hence
the first definition coincides with the third, which defines substance as ``that
which is in itself and is understood through itself, that is, that whose concept
does not need the concept of another thing from which it would have to be formed''.
\opage{15}What kind of concept can it now be that should be able thus to close our
knowledge without pointing beyond itself? As I have sought to show in the
Introduction, for \textsc{Spinoza} it is the principle of all conformity to law in
the world --- the principle that makes it possible to understand the succession of
things as a series of grounds and consequences. The thought of the relation between
ground and consequence becomes the first (or last) thought, and its content is, for
\textsc{Spinoza}, the being that underlies everything, while it has itself no ground
outside itself.

In the second and sixth definitions a distinction is made between something that is
finite in its kind, something that is infinite in its kind, and something that is
absolutely infinite. With this the relation between individual being (mode),
fundamental property (attribute) and fundamental being (substance) is stated.

As we have seen in the Introduction, the concept of attribute is grounded in there
being series formations that can be continued to infinity by the same rule with which
they began. This is shown in the first and second propositions of the second book in
the case of thought and extension. No limit can be set either to the world of
thought or to that of extension, and neither can they limit each other; for in that
case the law of the series formation would be changed. While no limit can be set
within an attribute, the attribute itself is nevertheless limited in so far as it
presupposes a definite rule of series formation. Therefore it can express only a
single side or property of existence. And then a new series can be formed of the
different attributes of existence. This series too can have no limit (although, as
will appear later, we know only two members of it), and the concept of God receives
its definition precisely in the concept of the substance consisting of infinitely
many attributes (Def. 6).

Through the series in which the attributes show themselves, substance expresses
itself. For the individual members (modes) of these series are held together by the
law of grounding. But there is (according to Def. 3 and 5) this difference between
substance and mode, that the former has its ground in itself, the latter its ground in
something other than itself. The common presupposition is that there must be a ground,
wherever it may lie. The principle of sufficient reason coincides for \textsc{Spinoza}
with the principle of causality, as we have already seen, and as is expressly
enjoined in the fourth and fifth axioms: ``The knowledge of the effect depends on the
knowledge of the cause and presupposes it. Things that have nothing in common with
each other cannot be understood through each other either, or the concept of the one
does not presuppose the concept of the other.'' The relation between cause and effect
is thus a pure relation of concepts, in which the difference of time plays no part.

\textsc{Spinoza} accordingly had really no reason to set up the proposition of
causality as a special proposition. We find, however, that in more special cases,
where the discussion seemed to demand it, he expressly formulates it. Thus Eth. I, 8.
Schol. 2: ``For every existing thing there is necessarily a cause on account of which
it exists.'' Here, then, necessary validity is ascribed to the proposition of
causality, and \textsc{Spinoza} expresses himself in the same way in Ep. 34. Where this
\opage{16}necessity lies, \textsc{Spinoza} does not say, and that he sees no problem
in it is understood when one holds fast that cause and ground are for him one, as he
expressly says in a demonstration of the eleventh proposition of the first book:
``For the existence or non-existence of every thing a cause or ground (\textit{causa
seu ratio}) must be given.''

But now the difference made between substance and modes nevertheless brings with it a
twofold causal relation, since there is a difference between having one's ground in
oneself and having it in something else. Eth. I, 28 (this important proposition,
already mentioned in the Introduction, and to which we shall come back later) demands
that each individual mode be explained by another mode, this by a third, and so on to
infinity. But at the same time each individual mode is to be understood as a peculiar
determination (\textit{affectio}) of substance (Def. 5), and it must therefore find
its explanation in substance --- besides finding its explanation in other modes. This
twofold causal relation --- one might call it the ideal and the elementary --- would
really have to be reduced for \textsc{Spinoza} to a single one, the ideal. For
substance is, through its attributes, the principle of all conformity to law in the
world, and so also of the laws that hold between the different modes; and on the
other hand these modes are understood through their mutual lawful content, that is,
by being regarded as moments in substance (\textit{affectiones substantiæ}).
Besides, according to Ax. 4 and 5, there must be no difference between cause and
effect. The elementary concept of cause, which was to hold for the relation between
modes different from one another, cannot therefore be a true causal relation. And in
the end the external relation between modes, and with it the elementary causal
relation, holds only on the standpoint of sensuous ideas (of the imagination). And the
last problem then becomes how the different standpoints --- that from which there
exists only the purely logical (or mathematical) relation between ground and
consequence in mutual equivalence, and that for which there exist different things in
an external causal order --- can be possible. The reduction of these two standpoints
has not been carried through to this day\footnote{Cf. my criticism of the Marburg
school in \textit{Totalitet som Kategori} (Vid. Selsk. Skr. 1917, p 46---48.} --- % the Danish has no ) and no stop after p; kept as printed
and certainly cannot be carried through.

--- A special form of the proposition of causality, the principle of inertia, is
contained in the second definition. For when it is said here that extension cannot be
limited by thought, this is grounded, as we have seen in the Introduction, in our
being able to think any line or surface or any volume lengthened or enlarged, that is,
to form members in a continuous series without leaving the attribute of extension. But
this proposition also contains for \textsc{Spinoza} the methodical principle that
every extended thing must be explained by another extended thing, that is, every
change in a body by changes in the state or position of another body. \textsc{Spinoza}
works by this principle, as we shall see in what follows, when faced with the special
questions. And when he demands (Eth. I, 10, to which proposition, besides Ax. 4, Eth.
II, 6 refers back) that every attribute (and so everything that falls under this
attribute) shall (\textit{debet}) be explained by itself, this follows, no doubt, in
the end from every attribute, according to Def. 4, expressing the essence of substance
and so, like \opage{17}substance, having to have its ground in itself. But this
definition itself is at bottom a methodical principle, grounded in there having to be
homogeneity between cause and effect if the causal relation is to be intelligible. ---
There is an interesting difference here between the Short Treatise and the Ethica.
There the principle of inertia is already set up (II, 19, 8), but no consequences are
yet drawn from it with regard to method (or, in \textsc{Spinoza}'s terminology, with
regard to the relation between the attributes), and therefore no new psychophysical
hypothesis is set up, as in the Ethica. \textsc{Descartes}, in his formulation of the
principle of inertia (Principia Philosophiæ II, 37---41), had expressly made a
reservation with regard to angels and souls. And on the strength of such a reservation
the Short Treatise confidently lets the soul influence the direction of physiological
processes, and on the other hand lets the body influence the soul and thereby be known
by it (K. Tr. II, 19, 9---14). Here, then, in the interval between the Short Treatise
and the Ethica, a change took place in \textsc{Spinoza}'s view, probably under the
influence of the methodical principle (\textit{vera causa}) mentioned in the
Introduction.

Now it is characteristic of \textsc{Spinoza} that these different principles (which,
for that matter, are not different for him) are not conceived as forms or conditions
of human understanding, determined by the nature of human thought and by the character
of the tasks set it, but are regarded as expressions of the essence of things, in the
end of the essence of the eternal ground of existence. When a thought is true, says
Ax. 6, it must agree with its object\footnote{This is the correct translation of Ax.
6 (\textit{idea vera debet cum suo ideato convenire}). In the very law of thought
itself the norm of truth is given, which also decides what is false (\textit{veritas
norma sui et falsi}, Eth. II, 43 Schol.). In the treatise on the improvement of the
understanding \textsc{Spinoza} develops the view that truth does not rest on anything
distinct from knowledge; every false idea will be dissolved through its consequences.
We shall come back to this question later. --- \textsc{Freudenthal} (Festschrift for
\textsc{Zeller} p. 128) wrongly regards Ax. 6 as a definition of truth and finds in it
a remnant of scholasticism.}. And what is the content of strict thought --- so it says
in Eth. I, 30 Dem.\footnote{The word ``objective'' in \textsc{Spinoza}'s usage means
``subjectively''. Cf. below at II, 7.} as a closer explanation of Ax. 6 --- must
necessarily be given in nature.

That in our search for the last presuppositions of knowledge we cannot go further back
than to certain principles is for \textsc{Spinoza} proof that these principles express
the nature of an eternal and necessary being that underlies everything (\textit{res
omnium prima}, as it is called in the treatise on the improvement of the
understanding), and that does not itself have its ground in anything different from
it. It is \textsc{Parmenides}' old teaching that still asserts itself: ``Thought and
the object of thought are one and the same'' (ταὐτὸν δ᾽ἐστιν νοεῖν τε καὶ οὓνεκεν ἐστι
νοήμα).

If we want to know more about substance, we are referred to the attributes, each of
which expresses the essence of substance from a single side (\textit{in suo genere}).
An attribute is, according to Def. 4, that which the understanding perceives of
substance as making up (constituting) its essence. From this definition
\textsc{Hegel}\footnote{\textit{Vorlesungen über die Geschichte der
Philosophie\textsuperscript{2}}, Berlin 1844, III, p. 344: ``Wie kommt ausser Gott der
Verstand herbeigelaufen, der die beiden Attribute des Denkens und der Ausdehnung auf
die absolute Substanz anwendet?''} concluded, and after him J. E. \opage{18}\textsc{Erdmann},
that it is only for our knowledge that there are such attributes. But ``as''
(\textit{tanquam}) does not mean ``as if'' here. \textsc{Spinoza} uses ``tanquam'' in
several places of valid perception (see e.g. Eth. II, 45 Schol.). The critical
philosophy's distinction between things in themselves and our perception of them must
not without more ado be read into \textsc{Spinoza}. For he teaches expressly that the
attribute, like substance, has its ground in itself. If \textsc{Hegel} were right, it
would have to be said that the attribute had its ground in the human understanding.
\textsc{Spinoza}'s characteristic teaching amounts precisely to this, that the last
foundation of knowledge is also the last foundation of existence. % the Danish misprints SRINOZA

In the relation \textsc{Spinoza} has stated between the concepts substance
(fundamental being), attribute (fundamental property) and mode (individual being,
thing), his whole philosophy is really contained. The special series of propositions
he sets up are only further developments of what is contained in the definitions. They
do not come about, as a really ``geometrical'' method would require, by combining
construction from the definitions set up, but are explications, paraphrases, not
proper demonstrations. The ``geometrical'' method is merely a shell that must be
broken if one wants to get hold of \textsc{Spinoza}'s real, working thought. And it has
the great drawback, which \textsc{Spinoza} himself sometimes feels, that the different
doctrines are put on the same line despite their mutual relations of super- and
subordination.

\subsection*{Excursus on Definitions, Axioms and Postulates in \textsc{Spinoza}.}

{\small
With regard to \textsc{Spinoza}'s ``definitions, axioms and postulates'' the following
is to be noted.

1) The definitions are set up as nominal explanations of fundamental concepts, but on
closer examination prove to be due to an analysis of \textsc{Spinoza}'s concept of
science. The results of this analysis are dogmatically fixed under the (unformulated)
presupposition that the foundation of thought is also that of existence. % the Danish misprints Tænkningeus
\par
2) The axioms have a different character in the different books. In the first book
they state principles regarded as self-evident, but which, like the definitions, prove
to rest on analysis of the concept of science. --- In the second, third and fourth
books the axioms are empirical propositions that stand for \textsc{Spinoza} as beyond
doubt. That ``man thinks'' (II, Ax. 2), and that ``we do not feel or perceive any
individual things other than bodies and thoughts'' (II, Ax. 5), stand fast as facts of
experience. Likewise, that every individual thing will be able to find its master (IV.
Ax.). After II, 13 --- in connection with borrowed propositions (\textit{lemmata})
from natural science --- axioms are set up that rest on experience (e.g. that bodies
are either at rest or in motion, and so on). These axioms are the scientific
counterparts of the axioms drawn from psychological experience at the beginning of the
book. --- Of the two axioms in the fifth book, one is an empirical-psychological
proposition, the other a proposition of the theory of knowledge derived from a
proposition in the third book. % the Danish misprints empirsk-psykologisk
\par
3) The expression postulate \textsc{Spinoza} uses only in the second and third books,
and he makes no difference in principle between axiom and postulate (cf. III, Ax. 1:
\textit{postulatum seu axioma}). The series of postulates set up in the third book
after Proposition 13 are only summaries of what was contained in the preceding
borrowed propositions (\textit{lemmata}) and axioms. These postulates are referred
back to when it says in Proposition 17 Schol.: ``All the postulates I have set up
contain hardly anything that does not stand fast by experience.''\par}

\opage{19}\subsection*{b. Propositions 1---14.}

\begin{center}\textit{Purging of the Concept of Substance and Demonstration of its
Identity with the Concept of God.}\end{center}

\textit{α.} That there must be substance in the world, that is, something that exists
in itself without presupposing anything else, is for \textsc{Spinoza} a matter of
course, because the series of our thoughts (of grounds and causes) must find its close
in something that does not in turn demand continued thought. The urge to a close shows
itself, as already mentioned in the Introduction, in another way in the Ethica than in
the Short Treatise. In the youthful work it was above all the urge to rest that was
stressed. ``If we'', it says here (I, 2), ``wanted to seek the cause of substance,
which is the principle of all things, we should have to seek in turn the cause of this
cause, and so on to infinity, so that --- since we must necessarily stop somewhere and
come to rest, which we must --- it is necessary to come to rest in this one
substance.'' This line of thought \textsc{Spinoza} later expressly abandoned. The
grounding of the existence of substance or God, he says in a letter (Ep. 12), does not
lie in our otherwise getting an infinite series of causes, but in there otherwise
being no cause necessary in itself of the things that are not
necessary\footnote{\textsc{Aristotle} already has a similar line of thought: ``If
there is no first, there is no cause at all.'' (\textit{Metaphys.} p. 994).}. That the
series of modes (individual beings, things) leads us from A to B and from B to C and so
on to infinity, --- in this \textsc{Spinoza} now finds no difficulty. On the contrary,
he is --- like \textsc{Bruno}\footnote{See \textit{Den nyere Filosofis
Historie\textsuperscript{2}}. I, p. 123 f.} --- convinced that from substance (nature,
the divine) an infinity of consequences or effects must proceed. So long as the work of
thought is occupied with modes, no close can therefore be expected. If a close is to be
reached, it must be found in the principle that underlies the lawful connection between
all modes, as well as these modes themselves according to their innermost nature; and
it is precisely this principle that \textsc{Spinoza} calls substance, nature or God.
Not the urge to rest, but absolutely rational satisfaction now motivates the close.

\textsc{Spinoza} thus does not think of the series of causes as closed by a first
member, from which all other members spring, but finds the close he seeks in the inner
law or force that makes itself known in the causal series that empirical science
demonstrates. Either nothing exists, it says in Eth. I, 11 Dem., or else an absolutely
necessary being exists. And on the other hand it also holds that the more we understand
individual beings, the more we understand God (Eth. V, 24).

\textit{β.} When the existence of substance stands for \textsc{Spinoza} as an eternal
truth, it is thus because it expresses the fundamental truth on which all other truth
depends, --- namely, that in existence there reigns a lawful order which finds its
expression in the concept of cause, which for \textsc{Spinoza} coincides with the
concept of ground.

As already remarked, the proposition of causality is expressly set up in the second
scholium to the eighth proposition. But this proposition is clearly presupposed in the
first eight propositions. For when \textsc{Spinoza} wants to establish that, according
to the definition of substance, there \opage{20}can be only a single substance, he
does it in the following way: If there were several substances, they would have to have
different attributes, for through attributes the essence of substances is expressed;
but substances with different attributes would not be able to be causes or effects of
one another, since they would have nothing in common; their concepts could not be led
back to one and the same concept.

There is a presupposition for the validity of this line of thought that
\textsc{Spinoza} does not indicate, namely, that different substances cannot be thought
without standing in a relation of interaction with one another. This is not necessary,
since every substance is to have its ground or cause in itself. It could after all be
thought --- and \textsc{Leibniz} in his doctrine of monads did think this possibility
--- that what springs from the one substance does not cross at all with, or stand in a
relation of interaction to, what springs from another substance. The different
substances could stand in the same relation to one another as, in \textsc{Spinoza},
the different attributes stand to one another.

If one defines the proposition of causality so that it states that everything we are
to be able to know must stand in relation to something else as ground to consequence
(or cause to effect) or conversely, then it would be correct that if there were several
``substances'', they would have to stand in interaction with one another. But then the
concept of a cause or ground of itself would have fallen away. This concept in reality
denotes a blind alley and is a purely negative concept, which states only that one has
not found the ground or cause of something and therefore cannot bring it into
connection with what one otherwise knows or thinks one knows. The first presuppositions
of our conception of the world can only be justified through their consequences, not
themselves derived in turn. The concepts they contain are therefore limiting concepts.
--- This consideration \textsc{Spinoza} would not be able to accept. It is first
carried through by the critical philosophy.

It is of interest to compare \textsc{Spinoza}'s concept of substance with the concepts
of substance of scholasticism and of \textsc{Descartes} on the one side, and with
\textsc{Kant}'s ``Ding an sich'' on the other.

% ===== pp. 20--26 =====
The older concept of substance (that of scholasticism and of \textsc{Descartes})
places substance, as existing in itself, in sharp opposition to the properties, and
it thereby gets a mystical character, since we know things only through their
properties and the connection of these with one another. In \textsc{Spinoza}, on
the other hand, there is an inner relation between substance and its properties
(attributes), since each of these expresses the essence of substance. What holds of
substance, that it has its ground in itself, holds also of each single attribute.
And if the interpretation given above of \textsc{Spinoza}'s substance as the very
principle of grounding or of knowledge is correct, the concept of attribute contains
the scientific principle that the domains of nature (the spiritual as well as the
corporeal) are each to be explained by their own laws. Nature is to be explained by
nature itself. --- Such an inner relation between theology and science, between the
knowledge of God and the knowledge of nature, scholasticism could not recognize.
According to it God cannot even stand in a real relation to the world he has
\opage{21}created. The world stands in relation to God, namely in a relation of
dependence; but God does not stand in relation to the world. The creation has no
significance for God's own being\footnote{Cf. my \textit{Religionsfilosofi} § 22 and
the passages of \textsc{Hugh of Saint Victor} and \textsc{Thomas Aquinas} cited in the
note. Cf. on \textsc{Augustine} the end of § 78.}. \textsc{Dunin-Borkowski} accordingly
finds the decisive difference between \textsc{Spinoza} and scholasticism in this, that
according to the former there is a relation of reciprocity between the infinite and
the finite (God and the world), whereas according to the latter the world does indeed
stand in a relation to God, but only by way of analogy (which here really means:
figuratively) can one speak of God's standing in relation to the
world\footnote{\textit{Der junge Spinoza}, p. 354: ``Dem Verhältnis der Abhängigkeit
des Geschöpfes vom Schöpfer entspricht in der unendlichen Ursache kein wirkliches
Gegenverhältnis''. \textsc{Dunin-Borkowski} ascribes to this difference ``eine
ungeheure Tragweite''.}.

\textsc{Kant}'s ``Ding an sich'' is a residual concept; it expresses that something
always remains over when scientific work comes to an end. It is at the same time a
limiting concept and a problematic concept, and --- what \textsc{Kant} does not always
hold fast --- doubt can always be raised whether the limit is fixed in the right way.
As we have seen in the Introduction, \textsc{Spinoza} admits that a remainder always
remains; only he maintains, and rightly, that this does not rob the knowledge that has
really been won of its significance. The criterion of truth is for him (as we shall
see in a later connection) the agreement of thought with itself within the world of
experience.

\textit{γ.} The concept of attribute denotes a property in which the essence of
substance, as existing in itself and intelligible through itself, shows itself.
\textsc{Spinoza} enjoins (10 Schol.) what \textsc{Arnauld} and \textsc{Regius} had
already maintained against \textsc{Descartes}, that one has no right, from two
attributes being irreducible to each other, to conclude that different beings or
substances lie behind them. On the contrary, the more reality or perfection a being
possesses, the more different attributes it must possess, and the concept of God is,
according to \textsc{Spinoza}, the concept of a substance that has infinitely many
attributes, each of which expresses eternity and infinity (Def. 6 and Proposition
11). The proof of God's existence coincides with the proof of the existence of
substance, except that with the concept of God the thought of reality or perfection
is of decisive significance. It stands for \textsc{Spinoza} as absurd that the
perfect should not exist, --- that only finite, limited beings with few properties
should exist, but not an infinite being with infinitely many properties.
``Perfection'', he says (11 Schol.), ``does not take away the existence of a thing,
but on the contrary grounds it, whereas imperfection takes it away. Therefore we can
be certain of the existence of no thing as of God's, the existence of the absolutely
infinite and perfect being.'' In the sixth definition of the second book reality and
perfection are expressly declared one and the same, and in the appendix to the first
book \textsc{Spinoza} rejects the standard of perfection taken from the relation of
things to human needs and human wishes: ``The perfection of things is to be judged
solely by their nature and their power.''

\opage{22}There is nevertheless a valuation hidden in \textsc{Spinoza}'s concept of
perfection, even if it is made one with the concept of reality. Precisely in this
identification there is expressed a deep admiration for the fullness of existence as
that which oversteps all limits, as regards richness of quality, time and space alike.
The same enthusiasm stirs here for the richness and unboundedness of nature that in
\textsc{Montaigne} lay behind his apparent scepticism as well as behind his sense for
the manifoldness of individualities, --- that made \textsc{Bruno} burst the bounds of
the old picture of the world, --- and that also inspired the great Italian artists.
\textsc{Leonardo da Vinci} was, as \textsc{Julius Lange} has
said\footnote{\textit{Udvalgte Skrifter}, II, p. 239. --- For \textsc{Montaigne} and
\textsc{Bruno} reference is made to \textit{Den nyere Filosofis
Historie\textsuperscript{2}}. I, p. 26; 30 and p. 111; 117---123.}, ``in love with
the manifoldness of the world''. There is here also an aesthetic element in
\textsc{Spinoza}'s view of the world. He admires the power and fullness of nature,
which has been able to let both the great and the small, both the good and the evil,
both the useful and the harmful come into being (see the closing words of the appendix
to the first book). Accordingly he does not merely make reality one with perfection,
but declares that perfection and beauty are closely related concepts (Ep. 54).

\textsc{Spinoza} really rejects the concept of value only when one wants to use it to
explain individual, special phenomena of nature. On the other hand he applies it to
the innermost ground and highest law of nature when he identifies God and substance
and declares reality and perfection one. The urge to value has not vanished in
\textsc{Spinoza}, but it has once and for all come to rest in the fundamental thought
of the great conformity to law of existence as determining all that, for finite
beings, comes to stand as good or evil. For \textsc{Spinoza} personally, as the
autobiographical hints at the beginning of the treatise on ``the improvement of the
understanding'' show, the knowledge of truth had come to stand as the highest human
good, and at the end of the fifth book ``the intellectual love of God'' is grounded in
its being the divine order of nature that makes such knowledge possible. Here there
was for \textsc{Spinoza} a value that had to have its ground deep in the nature of
existence.

\textit{δ.} In the assumption of infinitely many attributes the relation between
\textsc{Spinoza}'s religious-speculative line of thought and the knowledge from
experience on which he always lays weight shows itself in a peculiar way. For
\textsc{Spinoza} admits, as we have already seen in the Introduction, that he does not
know all the divine attributes, not even the greater part of them. In the fifth axiom
of the second book it is laid down that we do not feel or perceive any things (modes)
other than those that fall under the attributes of thought or extension, that is,
souls and bodies, and in Ep. 64 it is explained that this comes from human nature
neither presupposing nor offering any attributes other than these two. When
\textsc{Spinoza} nevertheless (Ep. 55) claims an exact knowledge of God, it must be
because he is convinced that the logical relation between ground and consequence is the
same within all attributes, which, after all, according to the interpretation given in
the foregoing, lies in their all being attributes \opage{23}of the same substance.
Only an infinite understanding (\textit{intellectus infinitus}) would be able to
comprise thoughts expressing the infinitely many attributes. That \textsc{Spinoza}
assumes such an \textit{intellectus infinitus} we shall come to see later.

The very idea of an infinite number of attributes is of Neoplatonic origin. For
\textsc{Plotinus} the principle of existence was an absolute unity that enclosed in
itself a fullness that could not be exhausted in any existence or grasped by any
thought. In medieval scholasticism this thought always lay in the background, but it
was refracted by the rationalist character of scholastic theology. After
\textsc{Descartes} had sharply marked off the life of the soul and matter as the two
domains of our inquiry, the question naturally came forward by what right we assume
that these two domains completely express the content of existence. In a work
approximately contemporary with \textsc{Spinoza}'s Ethica, \textsc{Malebranche} comes
to this question. ``One must not claim'', he says (Recherche de la vérité. III, 2, 9),
``that there are only minds and bodies, beings that can think and beings that are
extended. That could be a mistake. For although they suffice to explain nature, and
although one can therefore conclude, without fear of being mistaken, that the natural
things we can know depend on extension and on thought, there may yet be other beings
of which we have no idea and whose effects we do not see. It is then a hasty judgement
that men pass when they set it up as an indubitable principle that every being must be
either a soul or a body.'' --- The difference between \textsc{Malebranche} and
\textsc{Spinoza} on this point is only that where \textsc{Malebranche}, as a
Cartesian, speaks of beings or substances, \textsc{Spinoza} speaks of attributes of one
and the same substance.

There is not, as has sometimes been thought, any self-contradiction in the assumption
of infinitely many attributes. But a positive proof of this assumption
\textsc{Spinoza} would not be able to give. His most acute correspondent,
\textsc{Tschirnhausen}, accordingly urged him in vain to give one. No more has he been
able to give an explanation of how the manifoldness of the attributes can be derived
from the unity of substance. This difficulty is, however, the same whether only two or
infinitely many attributes are assumed. \textsc{Tschirnhausen} touches on yet another
difficulty, namely how it is known that all attributes stand on the same level. These
questions of \textsc{Tschirnhausen}'s (Ep. 65---66; 70) still stand. In an earlier
letter (Ep. 9) to a younger friend \textsc{Spinoza} had declared that he had two proofs
of the manifoldness of the attributes --- the one, that the more reality a being has,
the more attributes must be ascribed to it; --- the other, that the more attributes a
being has, the more we are compelled to take it to be real. --- Against both these
``proofs'', however, stands the fact, acknowledged by \textsc{Spinoza} himself, that we
do after all know only two attributes, and that the surest knowledge of reality we have
is built up on the basis of them. % the Danish misprints Kendsgerntng (fact)
--- In the Short Treatise \textsc{Spinoza} appealed to a presentiment we have that there
exist other sides of existence than those we know; and he had not given up the hope of
coming to know them (K. Tr. I, 1, 7. Note and VII, 1, 1. Note). In the letters and in
the Ethica this thought does not appear.

\opage{24}\textsc{Spinoza} could quite simply have kept to the spiritual and the
material standing to each other not in a contradictory but only in a contrary
relation, since we have no guarantee that all the sides of existence should appear
within our experience. Just as little as our experience can be said to be exhaustive
with regard to the number of things (modes), spiritual and corporeal, just so little
may it be said to be exhaustive with regard to the number of the main kinds of
existence it reveals to us. --- This consideration is, as I have sought to show
elsewhere\footnote{\textit{Den menneskelige Tanke}, p. 309 f; 334---338.}, not
without significance for our attitude to the problem of the relation between mind and
matter.

As for the equal standing of the infinitely many attributes, it rests, according to
\textsc{Spinoza}, on substance being expressed in all of them according to its essence.
If now the essence of substance consists in the very logical relation between ground
and consequence, the equal standing can hold only in so far as all attributes have a
rational character, contain the possibility of applying logic. But there is in itself
nothing to prevent their being otherwise, with regard to character and value, highly
different --- unless (as \textsc{Spinoza} does tend to) all character is reduced to
logic and all value to ``reality''. The thought of infinitely many attributes opens a
far larger vent for possibilities than \textsc{Spinoza} was aware of. The fact that we
know only two attributes sets a narrower limit to a scientific view of the world than
he saw.

\subsection*{c. Propositions 15---28.}
\begin{center}\textit{Closer Determinations of the Concept of God in its Relation to
the Concept of the World.}\end{center}
\begin{center}\textit{(Unity and Manifoldness).}\end{center}

\textit{α.} If God is the substance consisting of infinitely many attributes, and if
extension is one of these, then extension, the fundamental form in which everything
material appears, is a form of God's being. When this astonishes or offends many, it is,
\textsc{Spinoza} thinks, because by matter they understand a chaos of separated
elements. But it is only for the imagination, sensuous representation
(\textit{imaginatio}), that separated elements exist, just as it is only the
imagination that lets extension stand as limited. Extension is continuous, and
therefore unlimited and indivisible\footnote{\textsc{Huan} (\textit{Le Dieu de
Spinoza}, p. 14) reverses the relation, in that he thinks \textsc{Spinoza} relies on
indivisibility to prove that extension is a divine attribute. But \textsc{Spinoza}'s
polemic against the divisibility of matter is only to clear away an obstacle, not to be
a proof. Indivisibility follows, like infinity, from continuity. Cf. what was remarked
above on series formation.}. That matter, whose fundamental form is extension in space,
is a divine attribute does not therefore mean that God is one with the sum of the
material things the senses show us (15 Schol.). It is extension as the object of
geometry that \textsc{Spinoza} gives a place in principle in his philosophy, guided by
the decisive significance geometry has for all natural science. Besides, the
impossibility asserts itself \opage{25}of deriving the material from the spiritual
unless one assumes an incomprehensible act of creation. With equal necessity, according
to \textsc{Spinoza}, God reveals himself in material and in spiritual nature.
\textsc{Spinoza}'s attribute of extension corresponds to what \textsc{Descartes} in the
synopsis of his Meditations (\textit{Synopsis Meditationum}) calls ``body in general''
(\textit{corpus generaliter sumptum}) as opposed to the individual bodies. As
\textsc{Dunin-Borkowski} shows\footnote{\textit{Der junge Spinoza}, p. 357---360;
573.}, attempts were made among the scholastics from the end of the Middle Ages to form
a concept of existence in space that could fit both God and finite beings; but with
them there was only a question of analogy between divine and material space, not of
their identity. --- The relation between \textsc{Spinoza} and \textsc{Newton} with
regard to the concept of extension has been spoken of in the Introduction.

As was to be expected, \textsc{Spinoza}'s conception of the attribute of extension in
its relation to the things that sensation shows us is analogous to his conception of
the attribute of thought in its relation to the things that psychological observation
shows us. Just as he warns against making matter as a divine attribute one with
sensuous matter, so he also warns against conceiving the attribute of thought as one
with the states and activities of the soul that psychological observation teaches us to
know. Understanding and will cannot belong to God's being in the sense in which they can
be used of men. Human understanding shows itself in thinking elaboration and
observation, and so presupposes a given material; but for God there can be nothing
distinct from him and his thought that would only now have to be worked up into the form
of thought: for then God would indeed be dependent and limited! And human willing always
strives towards an aim not yet reached; such striving can be thought only in finite
beings. For God everything stands in eternal completion, without succession in time
(\textit{simul natura}). Only the name could divine and human understanding and will
have in common --- like the Dog Star and the dog, the barking animal on earth (17
Schol.). --- It is of interest to see that \textsc{Kant} in all three of his principal
works comes to the same result. For him too only a word remains when one wants to apply
psychological expressions to God\footnote{\textit{Den nyere Filosofis
Historie\textsuperscript{2}}, II, p. 95.}. The founding critical philosophy here
comes, then, to the same conclusion as the most thoroughgoing dogmatic philosophy.

When \textsc{Spinoza} nevertheless conceives ``thought'' as a divine attribute, he
undertakes a reduction and idealization of spiritual things similar to the one he
undertakes of material things in the attribute ``extension''. In both places it must be
understood thus, that the attribute means the conformity to law that holds within one
of the essential sides from which existence can be seen, --- a conformity to law that is
the condition of the existence of the individual things. Thought as a divine attribute
means that the ideal of knowledge exists, is not first to be attained, but is an eternal
reality in which it is the task of the man who seeks truth to gain a share. When
\textsc{Spinoza} in his theologico-political treatise (Ch. 14) says that the God of the
positive \opage{26}religions has significance as a practical model (\textit{exemplar
veræ vitæ}), the same really holds, although \textsc{Spinoza} would not acknowledge
it, of the God of his philosophy. Under the attribute of thought \textsc{Spinoza}'s God
is the eternal ideal of truth, and also, by having the ground in himself, the eternal
ideal of freedom.

\textit{β.} That all individual beings (modes, \textit{res particulares}) that appear as
things in our experience spring from God's being as closer determinations of it
(\textit{affectiones substantiæ}), \textsc{Spinoza} grounds in their presupposing
substance and the attributes (15---16; 25). It is, as we have seen earlier, the
scientific investigation of spiritual and material phenomena that leads to the knowledge
of the conformity to law of nature and through it to the knowledge of God or substance.
But --- it is now maintained in Proposition 28 --- what is finite and has a limited
existence cannot be derived from the absolute nature of any divine attribute, for from
this only what is itself infinite can follow. A finite thing finds its explanation only
in another finite thing, this in a third, and so on to infinity. --- The principle that
\textsc{Spinoza} maintains here coincides, as already stated in the Introduction, with
what one might call the principle of the valid cause, \textsc{Kepler}'s \textit{vera
causa}. What appears in experience is to be explained by a cause that can itself be
shown in experience\footnote{In Eth. II, 17 Schol. \textsc{Spinoza} uses the expression
\textit{vera causa} in this way. Otherwise he uses the expression \textit{causa
proxima}.}. \textsc{Galileo} applied this principle in his dialogues when he maintained
that in the middle of a special investigation (e.g. of the relation of the earth to the
sun) one could not appeal to an intervention by God. And his grounds are the same as
\textsc{Spinoza}'s: from an infinite cause follow infinite effects; God is the cause of
everything, but in the individual case the task is to find the special causal relation.

The opposition, not to say contradiction, that undeniably exists between Proposition 15
and Proposition 28 can be removed within \textsc{Spinoza}'s presuppositions if one takes
as the basis the conception of \textsc{Spinoza}'s concept of substance put forward in the
foregoing. When a thing given in experience can be explained by another, and an infinite
series of finite causes and effects can thus be formed, then precisely in the law that
binds the members of such a series together a principle expresses itself, an order of
things, that is not itself a member of the series and does not itself appear as a thing
for observation. And it is such a principle that \textsc{Spinoza} has expressed in his
concept of substance and in the concept of attribute partly identical with it. On this
conception the two concepts of cause with which \textsc{Spinoza} operates can be
joined. That a thing is explained by its place in the series of things means the same as
that it is explained by substance (or by the attribute). Thereby a proposition such as
Eth. V, 24 is understood: ``The more we understand individual things, the more we
understand God.''

It seems to me that the agreement (or at any rate freedom from contradiction) that
becomes possible in this way between the two important Propositions 15 and 28 speaks in
favour of the conception of \textsc{Spinoza}'s concept of substance that I have put
forward.

% ===== pp. 26--32 =====
But a difficulty still remains, which follows from the
identifica\opage{27}tion of ground and cause in \textsc{Spinoza}'s system. The relation
between ground and consequence is timeless (``eternal''); the relation between cause
and effect presupposes (in so far as it is different from the relation between ground
and consequence) a temporal relation. There is a temporal relation, not merely a
relation of grounding, between the members of the causal series. Thereby a difference
of principle nevertheless remains between modes (things) as members of the causal series
that can be built on the ground of experience, and the relation of partial identity in
which the individual mode stands to substance. Besides, the individual modes each have
their qualitative peculiarity. The problem in \textsc{Spinoza}, as for modern standpoints
in the theory of knowledge, is the relation between the timeless and qualityless world
of logic and mathematics and the world of experience, which is the world of qualitative
changes. This world presupposes the former if it is to be understood, but cannot be
derived from it. \textsc{Tschirnhausen} already pointed out that motion and figure cannot
be derived from the general attribute of extension, and that the difference and
manifoldness of things cannot in general be derived from substance or the attribute in
their absolute unity (Ep. 59. 80. 82).

That \textsc{Spinoza} noticed the problem is seen from his inserting an intermediate
concept between substance (and the attributes) on the one side and the concrete things
(modes) on the other. This intermediate concept is the concept of the \textit{modus
infinitus} (23). Infinite modes spring (see 28 Schol.) immediately or mediately from the
essence of substance or of the attributes and are eternal like them. As examples he
names under the attribute of thought the infinite understanding, and under the attribute
of extension the constant relation between rest and motion (Eth. I, 21. 30. V, 23. 40
Schol.; cf. Ep. 32 and 64). These concepts already appear in the Short Treatise, and
\textsc{Spinoza} there (I, 2), as in Ep. 73, calls their objects ``sons of God'', to
signify that they stand in a closer relation to substance than the finite modes.

If we try to translate these half-mythological expressions into the language of the
theory of knowledge, \textsc{Spinoza}'s view can be said to be that the persistence of
the life of thought and of motion in the world under all vicissitudes is a consequence of
the principle of grounding or of causality. Just as \textsc{Spinoza} identifies grounding
and causality, so he identifies the proposition of causality and the proposition of the
conservation of motion (the imperfect forerunner of the conservation of energy). More
recent philosophers too have made this last identification. But just as there is a leap
from grounding to causality, so there is also a leap from the proposition that
everything has its cause to the proposition that cause and effect are equivalent, and
there is again a leap from the general proposition of equivalence to the special, actual
transformations that in every individual case take place in a definite direction. Here
the critical philosophy must distinguish what \textsc{Spinoza} identifies\footnote{Cf.
on these matters \textit{Den menneskelige Tanke}, p. 211---215; 281---287.}.

That the intermediate members the infinite modes denote do not solve the difficulty lies
in this, that the finite modes can no more be derived from them than they can be derived
directly from substance. It is a presupposition in \textsc{Spinoza} that there are,
after all, finite things. Therefore, as we shall see later, experience is the first step
on \opage{28}the scale of knowledge, and therefore it is only experience that teaches us
that modes exist at all (Ep. 10). \textsc{Spinoza}'s philosophy is an attempt to climb up
by analysis of the data of experience to attributes and substance. The infinite modes are
rungs on this ladder. \textsc{Spinoza}, without being clearly aware of it, pulls the
ladder up after him, so that a deductive descent does not become possible.

The opposition between an elementary and an ideal causal relation, which always asserts
itself in scientific work, has in \textsc{Spinoza} crystallized into an opposition
between a world of qualitatively different things (modes), which arise and succeed one
another in time, and a system of eternal thoughts, within which continuity reigns (within
each individual attribute) and in the end identity (in substance). The former world is a
revelation of the latter system appearing to sensation, memory and imagination. Why and
how, besides substance, attributes exist, besides the attributes eternal modes, and
besides these eternal modes modes differing in quality and determined in time, --- or why,
besides pure thinking (\textit{intellectus}), sensuous representation
(\textit{imaginatio}) also exists, --- these are questions \textsc{Spinoza} does not
raise, and which, for that matter, neither he nor anyone else would be able to answer.
All science, philosophy too, works with certain presuppositions that can progressively
stand their test in the purgatory of experience and criticism, but which cannot be shown
to be the only possible ones. \textsc{Spinoza}'s successor in the history of philosophy,
\textsc{Leibniz}, expressed this in the proposition that the success of a hypothesis
does not prove its truth.

\subsection*{d. Propositions 29---36.}
\begin{center}\textit{Closer Examination of the Relation between the Fundamental Being
and the Individual Beings.}\end{center}

\textit{α.} Besides the concepts substance and God, which (when substance is thought
with infinitely many attributes) are identical, there is yet a third concept that
becomes identical with them both, so that \textsc{Spinoza} says ``God or nature''
(\textit{deus seu natura}, e.g. Eth. IV, Preface), just as he says God or substance.
This concept is suited to suggest a connection between the two worlds which, as we have
just seen, threaten to come into sharp opposition to each other.

After it has been said in 29 that in the nature of things (\textit{in rerum natura})
there is nothing contingent, but that everything is determined by the necessity of the
divine nature (\textit{ex necessitate divinæ naturæ}) to exist and act in a definite
way, it says in the scholium (29 Schol.): ``Before I go further, I will explain, or
rather recall, what I understand by nature producing (\textit{natura naturans}), and
what I understand by nature produced (\textit{natura naturata}). I think it is
sufficiently clear from the foregoing that by nature producing one must understand that
which is in itself and is understood through itself, or, what is the same, such
attributes (\textit{attributa}) of substance as express eternity and infinity, that is,
God in so far as he is regarded as a free cause. By nature produced, on the other hand, I
understand everything that follows from the necessity of God's nature or (what is
\opage{29}the same) of a divine attribute, that is, all the individual beings determined
by the divine attributes (\textit{dei attributorum modi}), in so far as they are regarded
as things that are in God and without God can neither be nor be conceived.'' --- In other
words, to \textit{natura naturans} belong the attributes, to \textit{natura naturata}
the modes.

As we have seen earlier, \textsc{Spinoza} set up the concept of infinite modes and stated
that such can follow either immediately from the attributes, or mediately, namely through
the infinite modes that have arisen immediately (21---23; 28 Schol.). Both these kinds of
infinite modes belong, like the finite modes, to \textit{natura naturata}.
\textsc{Spinoza}'s most acute correspondent, \textsc{Tschirnhausen}, wanted a closer
explanation here too and asked for examples of such infinite modes as sprang immediately,
and such as sprang mediately, from God's nature. In his answer (Ep. 64) \textsc{Spinoza}
names as examples of the first kind: within the attribute of thought the infinite
understanding, and within the attribute of extension motion and rest (by which, no doubt,
as already remarked above, is meant the relation between motion and rest --- \textit{ratio
motus ad qvietem}, as it is put in Ep. 32, that is, the conservation of motion). As an
example of the mediately produced infinite modes he names only ``the character of the
whole universe'' (\textit{facies totius universi}, ``the physiognomy of the All''),
which, despite all changes in the particular, nevertheless remains one and the same, and
he refers to the second book, the scholium to the seventh borrowed proposition (Eth. II.
Lemma 7. Schol.), where it says that ``the whole of nature is one individual, whose parts,
that is, all bodies, vary in infinitely many ways without any change in the whole
individual''. A similar statement is found in Ep. 32, where it says ``that every body is
surrounded by other bodies, and that all bodies determine one another to be and to act
in a certain definite way, yet so that in the whole universe the same relation between
rest and motion is preserved''. Lowest in the series come then the finite modes, the
phenomena of soul and body that are objects of observation.

We thus get the following scheme for \textsc{Spinoza}'s most important concepts:
% the print sets the scheme as a diagram with braces; given here as a table
\begin{center}\small
\begin{tabular}{lccl}
 & \multicolumn{2}{c}{\textit{Substantia.}} & \\
 & Cogitatio. & Extensio. & \}\ Natura naturans. \\[0.5em]
 & Intellectus infinitus. & Motus et qvies. & \\
Modi infiniti. & \multicolumn{2}{c}{(immediate producta).} & \\
 & \multicolumn{2}{c}{Facies totius universi.} & \}\ Natura naturata. \\
 & \multicolumn{2}{c}{(mediate producta).} & \\[0.5em]
Modi finiti. & \multicolumn{2}{l}{\{\ Res singulares (mentes et corpora).} & \\
\end{tabular}
\end{center}

We see, then, that \textsc{Spinoza}, as corresponding to ``the conservation of motion''
in material nature, assumes a spiritual persistence, an infinite thought that remains the
same under the alternation of the special spiritual phenomena. From hints given
\opage{30}in the fifth book one sees that he thinks of the \textit{intellectus infinitus}
as containing the eternal that appears in the individual spiritual beings in a form
limited in time and as the correlate of a limited material existence. On the other hand,
several students of Spinoza have rightly missed a concept that, within the attribute of
thought and subordinate to the \textit{intellectus infinitus}, could correspond to ``the
physiognomy of the universe'', which under the attribute of extension is subordinated to
``the conservation of motion''. Attempts have been made in various ways to construct this
missing member of the system. This point is, however, hardly of great significance,
although it shows how a system reaches its full articulation only gradually and under the
influence of questions and doubts. \textsc{Spinoza}'s own answer to the above-mentioned
question of \textsc{Tschirnhausen}'s is very short and seems written somewhat hastily;
he does not even take time to criticize the view of a correspondent according to whom
(Ep. 63) the infinite modes that proceed immediately from substance were supposed to be
the attributes! This correspondent (\textsc{Schuller}) had not made himself familiar with
the concept of attribute as that which expresses the very essence of substance. And it is
said expressly (23 Dem.) that an infinite mode is an immediate or mediate consequence of
the divine nature of an attribute.

The most obvious thing would be to think of a kind of world-soul as the correlate of ``the
physiognomy of the world'' or ``the world-individual''. \textsc{Bruno} sets up such a
concept\footnote{\textit{Den nyere Filosofis Historie\textsuperscript{2}}, I, p.
127---129.} as the expression of the force that makes interaction in space possible and
also stirs innermost in the individual elements. When \textsc{Bruno} says of ``the
world-soul'': \textit{amina tota in toto et qvalibet totius parte}! (the whole soul in % the Danish misprints amina (anima); kept as printed
the whole and in every part of the whole!) --- one is reminded by it of the expression
that occurs more than once in \textsc{Spinoza} (Ethica II, 37---38; 44): ``what occurs
in the same way in the part as in the whole''. By this expression \textsc{Spinoza} has in
mind the general laws, whereas \textsc{Bruno} has in mind a force that carries these laws
through in the whole and in the part. This concept was too mystical and poetical for
\textsc{Spinoza}, and he found nothing that could take its place.

For \textsc{Spinoza} the lawful order of nature, which is conditioned by ``the
conservation of motion'' and the corresponding relations within other attributes, was not
an abstraction (a \textit{universale}), but a characteristic side of existence, an
individual peculiarity of the unique world we know. The \textit{modi infiniti} are both
individual and universal. They probably coincide with what \textsc{Spinoza} in the
treatise on ``the improvement of the understanding'' calls ``the fixed and eternal things,
in which the laws of nature are inscribed as in law books''. Of them it is said that
``although they are individual (\textit{singularia}), they can yet stand for us as
universal (\textit{tanqvam universalia}) on account of their omnipresence and
comprehensive working; they are as it were general definitions of the changeable things
and the nearest grounds of everything''. These fixed and eternal things thus correspond to
what in the Ethica are called \textit{modi infiniti} of the second class (\textit{mediate
producta}). In any case this concept bears witness that \textsc{Spinoza} saw in the order
of nature that science discloses a great concrete fact, a unique thing. The general laws
serve only to characterize this unique thing. And as \opage{31}will appear later, the
knowledge of general laws is not for \textsc{Spinoza} the highest kind of
knowledge\footnote{This conception of ``the fixed and eternal things'' is, it seems to
me, clearly indicated in \textit{Den nyere Filosofis Historie\textsuperscript{2}}, I,
p. 302, for which reason \textsc{Huan} (p. 251) wrongly names me among those who
interpret the expression as a designation of general laws. As against \textsc{Huan} I
hold decidedly that \textit{res fixæ et æternæ} are \textit{modi infiniti mediate
producti}. --- What may perhaps be illuminating here is the demonstration of the
significance for the theory of knowledge of typical ideas of individuals:
\textit{Totalitet som Kategori} § 12---13.}.
% the Danish "fortolker Udtrykket for som Betegnelse" (an extra "for", as printed)

The question how we rise from the finite modes to ``the fixed and eternal things'', and in
general to the infinite modes, \textsc{Spinoza} did not examine more closely. He announces
an investigation of ``the procedure of the modern empiricists'', but his treatise on the
theory of knowledge was not completed, and so this question was not treated by him.

\textit{β.} When \textsc{Spinoza} in the definition of \textit{natura naturans} calls God
a free cause (29 Schol.), the definition of freedom given from the first (Def. 7), as
determination by the nature of one's own being, must be held fast. In Proposition 33,
Schol. 2 \textsc{Spinoza} goes more closely into this question in order to reject the
frequent conception of freedom as indifference, an equal possibility of two opposites.
Freedom in this last sense would, he thinks, conflict with God's perfection. One likes to
think of God first as existing and then as taking his decisions. But for an eternal being
there is no ``before'' or ``after''. God's decision is one with his eternal being and can
therefore neither begin nor be changed. There is in God no possibility distinct from
reality. Yet \textsc{Spinoza} finds the assumption of an absolute indifference in God less
unreasonable than the assumption that God should act for the sake of the good or with the
good in view (\textit{sub ratione boni}): that would be to let God strive after a model,
an aim, that he has not yet reached, --- indeed, it would really make him dependent on a
fate (\textit{fatum}), --- he who according to the nature of his being is the first cause
of everything!

\textit{γ.} The last proposition of the first book (36) asserts an often misjudged
element in \textsc{Spinoza}'s system. It runs: ``Nothing exists from whose nature some
effect does not follow,'' and it is grounded in everything that exists expressing God's
nature or essence in a certain definite way (cf. 25---26).

This proposition is of special significance for understanding the part the urge of
self-preservation in finite beings plays as the foundation of \textsc{Spinoza}'s
psychology and ethics (in the third and fourth books; see, however, already II, 45
Schol.). It is as an element in substance that the individual mode can be said to be cause
of itself, to have its ground in itself, although this independence is refracted by the
action of other modes, each of which can also be said to have its ground in itself. Every
member of the causal series is both cause and effect, although one is inclined to keep to
the dependence of the following members on the preceding ones. The refraction of activity
and passivity will prove to play a great part in \textsc{Spinoza}'s psychology and
ethics.

\opage{32}\subsection*{e. Appendix to the First Book.}
\begin{center}\textit{Theology and Science.}\end{center}

\textit{α.} When \textsc{Spinoza} asked himself what prevented the recognition of the
great connection between the things in the world (\textit{rerum concatenatio}) from
becoming general, he found the explanation in men being led in their contemplation of
nature by a practical urge, not by an urge for science. They are therefore inclined to
believe that everything is arranged for their own good, that which is the object of their
own constant striving. From the first they are ignorant of the causes of things, also of
the cause of the urge they have to seek their own good. They believe themselves free when
they are led by this urge, and they believe that the means of satisfying this urge, in so
far as they cannot produce them themselves, are produced by beings resembling themselves.
When misfortunes befall them, they derive them from the wrath of the gods over men's sins.
And if unsolved riddles still remain, they comfort themselves that the wisdom of the gods
far surpasses the understanding of men. Against this whole artificial system
(\textit{fabrica}) the cause of truth would have been hopeless had not mathematics, which
asks not about purposes but only about the essence and properties of things, given men
another norm of truth than the one their blind urge leads them to follow.

The teleological-theological view turns things upside down. It makes the first the last,
the cause the effect, the perfect the most imperfect. The eternal order of nature is not
the effect of men's urges and wishes, but contains the explanation of men's being able to
feel urges and cherish wishes at all. And the perfect is one with the eternal order of
nature, which rests on no presupposition; the imperfect is that which presupposes many
conditions. The theological view makes God imperfect by letting him strive after something
he still lacks. And it introduces a wholly new way of thinking: instead of proving
something by establishing the impossibility of the opposite (\textit{reductio ad
impossibile}) it appeals to ignorance (\textit{reductio ad ignorantiam}). And God's
incomprehensible will becomes the last refuge of ignorance (\textit{asylum
ignorantiæ}). Whoever seeks natural causes of what awakens wonder and awe is regarded as
an impious heretic by the persons in whom the great mass sees interpreters of the essence
of nature and of the will of the gods. And these persons see well enough that if ignorance
were removed, the thoughtless awe (\textit{stupor}) would also vanish, without which their
authority cannot be kept up\footnote{In Ep. 73 \textsc{Spinoza} distinguishes between
religion and superstition thus, that the former is grounded in wisdom, the latter in
ignorance.}.

One cannot at all, \textsc{Spinoza} continues, use concepts such as good and evil, order
and confusion, beautiful and ugly to explain nature, since such concepts are due to human
needs and do not come from any connection that can be shown in nature. And when one
complains of the many imperfections in nature, one forgets that it is not human benefit
and harm that can supply the standard \opage{33}of perfection and imperfection. The
standard must be sought in the very nature and power of things (\textit{ex sola rerum
natura et potentia})\footnote{From the identity of reality and perfection it follows that
what has no perfection at all cannot exist. Therefore the Devil cannot exist! (\textit{Den
korte Traktat}, II, 25).}. If one then asks why God did not create men so that they were
led by reason alone, the answer must be that the laws of nature are so comprehensive that
everything an infinite understanding could think, from the highest degree of perfection to
the lowest, could find its place.

This whole line of thought is highly characteristic of \textsc{Spinoza}. He stands here as
one who asserts the right of the strict scientific method, and also as one who is convinced
that by carrying this method through fully nothing is lost of what really has value. Only,
men must not measure the value of existence by the needs of their own little world. Within
the human world the concepts of value whose cosmic validity he rejects have their
justification, as he tries to show in the fourth book.

It is, as already mentioned, a concept of value after all that \textsc{Spinoza} uses when
he identifies perfection and reality. When in the fourth book he goes over to a specially
human theory of value, there therefore arises not only a refraction between cosmology and
ethics, but also a refraction between two measures of value. To this we shall come back in
its place. Here it need only be remarked that whereas \textsc{Spinoza} from the
cosmological point of view determines the perfect as that which has no presuppositions,
from the specially human point of view (which he himself adopts in the fourth book) it
must conversely be said that the perfect is precisely that which rests on the most
presuppositions. Life has more presuppositions than the lifeless, the life of the soul than
merely organic life, the ethical, intellectual and aesthetic life than the elementary life
of the soul. The doctrine of categories shows that the concepts of value are the most
special of all our fundamental concepts\footnote{See on this \textit{Den menneskelige
Tanke}, p. 237---241.}.

\textit{β.} In the dispute about \textsc{Spinoza}'s philosophy called forth by
\textsc{Jacobi}'s ``Briefe über die Lehre des \textsc{Spinoza}'' (1785), especially by
the oral utterances of \textsc{Lessing} that \textsc{Jacobi} quoted, the question came
forward in particular whether everything could be derived from the thought of a highest
purpose. In a conversation with \textsc{Jacobi}, \textsc{Lessing} had called himself a
Spinozist and declared it a prejudice that thought should be the first and highest;
extension and motion are, just as much as thought, grounded in a higher force that is not
exhausted in these forms. He had spoken against ``unsere elende Art, nach Absichten zu
handeln'' being held the best method. Besides \textsc{Spinoza}, \textsc{Lessing} may here
have been influenced by \textsc{Hume}, whose ``Dialogues on Natural Religion'' had
appeared shortly before. The conversation between \textsc{Lessing} and \textsc{Jacobi}
took place in the summer of 1780. When it became known, it aroused great attention and
offence, since the popular theology of the Enlightenment (which, it had been believed,
\textsc{Lessing} too held) was overthrown if one could not infer God's existence from the
purposiveness of nature\footnote{On \textsc{Lessing}'s place in the history of philosophy
see \textit{Den nyere Filosofis Historie\textsuperscript{2}}, II, p. 16---25).}.

\opage{34}According to \textsc{Jacobi}, \textsc{Spinoza} was an atheist because his
philosophy denied God's personality. \textsc{Spinoza} himself had rejected this
designation, which he thought was due to the misunderstanding that he made God dependent
on fate. On the contrary, he says (Ep. 43), I hold that everything follows with inevitable
necessity from God's nature. For \textsc{Spinoza}, after all, freedom and necessity do not
exclude each other. And he points out besides that in his view the ethical ideals are
independent of whether we derive them from God as judge or from the eternal necessity of
the divine nature.

It is not easy to decide how \textsc{Spinoza}'s standpoint is to be characterized in the
philosophy of religion. It is generally thought that every respectable philosophy must
represent some -ism or other, and it is thought besides that one is once and for all in
possession of a catalogue of all possible -isms, in which every philosophy must be able to
find its place. --- As for what is characteristic of \textsc{Spinoza}, two points above all
are of significance. He does not hold the usual theological idea of God and the world as
two different beings, but God is for him the innermost being of the world, which makes
itself known in its laws. And he denies that the concept of personality can be applied to
God, since personality (understanding and will), as we have seen above, presupposes in his
view something given that is to be understood, and a resistance that is to be overcome for
the sake of a purpose. In the concept of pantheism, which \textsc{Toland} introduced at the
beginning of the eighteenth century, and which is usually taken to express an inner
relation between God and the world and to deny God's personality, it has often been
thought that one had a rubric into which \textsc{Spinoza} could fit. But according to
\textsc{Spinoza}, after all, thought (\textit{cogitatio}), although everything cannot be
derived from it, is an eternal attribute of substance, even though the special forms we
call understanding and will do not belong to ``nature producing''. In the concept of this
attribute there lies an idealization similar to the one that philosophical theists (such
as \textsc{Lotze}) undertake when they sublimate the concept of personality so that it is
supposed to fit a being that is to have its ground in itself. The boundary between
pantheism and theism is therefore not easy to draw\footnote{See more closely on the
concepts of pantheism and theism my \textit{Religionsfilosofi} § 22 and § 95. --- On the
purely psychological way, adopted there, of considering what underlies these concepts, I
have, since I wrote my \textit{Religionsfilosofi}, come to lay ever greater weight.}.

By this consideration, which has already been put forward in ``Den nyere Filosofis
Histori'' (2nd ed. I, p. 316), I have only wanted to show how difficult it is to bring a
deep-going work of thought like \textsc{Spinoza}'s under rubrics fixed in advance.
\textsc{Huan} therefore wrongly names me among those who interpret \textsc{Spinoza}'s
concept of God in a theistic direction\footnote{\textit{Le dieu de Spinoza}, p. 221.}.
\textsc{Spinoza}'s thought has on the point spoken of here, as on many other points, such
an individual peculiarity that it is not easy to classify, but, as \textsc{Huan} aptly
says on the last page of his book, can really only be called --- Spinozism!
% "Histori" as in the Danish

The boundary between theology and philosophy \textsc{Spinoza} finds in this, that the
former uses psycho\opage{35}logical expressions of God, whereas the latter denies such
expressions any justification. When he warns philosophers against using theological
expressions (Ep. 23), it is this he has in mind. But behind his own presentation,
especially in the concluding remark to the first book, there lies a still more fundamental
problem, namely the justification of laying concepts of value at the foundation at all in
closing our view of the world. \textsc{Spinoza} himself has kept more of the theological
concept of God than he is aware of. His substance, according to the definition, was to
contain only the principle of intelligibility (of rationality); but it encloses ---
through the definition of perfection as reality --- a concept of value as
well\footnote{Cf. my \textit{Religionsfilosofi} § 24 and § 80---89.}. The ground
\textsc{Spinoza} gives for this (in the fifth book) is, to be sure, only that it is
substance that is the cause of the joy of knowledge. But then the question turns up of the
relation between intellectual value and other values.

\opage{36}\section*{Second Book.}
\subsection*{a. Definitions and Axioms.}

As an introduction to the second book \textsc{Spinoza} gives some initial determinations
and principles that are of interest for the theory of knowledge and for psychology. The
content of the second book may be said to be the psychology of knowledge and the
investigation of the validity of knowledge (theory of knowledge).

\textit{α.} Of special importance is Def. 3, according to which an idea (\textit{idea}) is
a concept that the soul (\textit{mens}) forms because it is a thinking being. And this
definition is explained more closely, the weight being laid on every idea being a product
of spiritual activity (\textit{actio mentis}); therefore the word concept
(\textit{conceptus}) fits better than the word perception (\textit{perceptio}), which might
seem to suggest that consciousness is passive towards its object. --- \textsc{Spinoza} here
corrects his own earlier view, since in the Short Treatise he let consciousness stand
passive towards its objects and let the one attribute act on the other.

Def. 4 gives a further development of these remarks, in that it determines an adequate
idea (\textit{idea adæqvata}) as one that, considered in and for itself and apart from
the object, has all the inner properties of a true idea. The outer properties concern the
agreement with the object (\textit{ideatum}). What the inner properties are can be seen
from the treatise on the improvement of the understanding, where it is said that the form
of a true idea rests on its own power and nature and shows itself in the connection
between subject and predicate, ground and consequence. This inner connection can lead us
out of falsity and doubt and is the real norm of truth. No outer sign of truth can be
given; the outer properties of the idea, its agreement with the object, are therefore
derived in relation to the inner ones --- as can already be seen from Eth. I, Ax. 6, when
this axiom is rightly understood. (See above).

When \textsc{Spinoza} was asked by \textsc{Tschirnhausen} what difference there was between
an adequate idea (\textit{idea adæqvata}) and a true idea (\textit{idea vera}), he
answered (Ep. 60) that at bottom there was no difference, except that the word ``true''
refers to the agreement of the idea with its object (\textit{ideatum}), whereas the word
``adequate'' refers to the nature of the idea in and for itself.

From the object of the idea (\textit{objectum = ideatum = res percepta}) \textsc{Spinoza}
distinguishes (as is seen from Propositions 5, 7, 13 and 15) the content of the idea
(\textit{esse formale ideæ}, \opage{37}cf. I, 30 Dem.: \textit{id, quod in intellectu
continetur}). The objects of the concepts triangle (I, 11 Dem.) or circle (II, 7 Schol.)
are the triangles or circles really existing (\textit{in natura}); the object of the idea
in which the soul consists is the body corresponding to the soul. But the content of these
concepts is given with their scientific definitions. The objects (triangles, circles,
bodies) fall under the attribute of extension, but the content of the concepts belongs to
the attribute of thought.

What is striking in this arrangement is, first, that the activity of thought itself, which
\textsc{Spinoza} stresses so strongly in Def. 3 and 4, as well as later on (43 Schol.; 49
Schol.), is not named here. It plainly falls under the attribute of thought, just as much
as the conceptual content it leads to fixing, and there might be reason to ask in its case % the Danish has a stop for a comma (spørge. hvad)
too what corresponds to it under the attribute of extension. It will appear that this point
is bound up with an unclarity in \textsc{Spinoza}'s teaching on the relation between the
attributes. --- Secondly, it is striking that triangles and circles stand on a line with
human bodies and are thus regarded as real objects of nature. Thus (II, 7 Schol.) the
circle existing in nature is spoken of. This point is bound up with \textsc{Spinoza}'s
dogmatic theory of knowledge, according to which a reality corresponds to every
scientifically formed concept.

\textit{β.} The concept of spiritual activity (\textit{actio mentis}) set up by Def. 3
presupposes a definition of activity that \textsc{Spinoza}, strangely enough, gives only in
the second definition of the third book, but which ought to have had its place here. It is
not the only occasion there is for agreeing with \textsc{Leibniz} when he wrote in his own
copy of the Ethica that demonstration was not \textsc{Spinoza}'s strong side
(\textit{\textsc{Spinoza} non est magnus demonstrator}).

The definition of activity (III, Def. 2) runs thus: ``I say that we are active
(\textit{agere}) when something happens in us or outside us of which we are the full cause
(\textit{adæqvata causa}), that is, when something follows from our nature, in us or
outside us, that can be clearly and distinctly understood through it alone. On the other
hand I say that we are passive (\textit{pati}) when something happens in us, or follows
from our nature, of which we are only partly the cause (\textit{causa
partialis}).''

In this definition it is implied that we are never wholly passive, whereas we might very
well be absolutely active, that is, the full cause of something that happens in us or
outside us. Now to be the full cause is, according to the seventh definition of the first
book, one with being free. Full activity or freedom becomes for \textsc{Spinoza} an ideal
(\textit{exemplar}) towards which ethical striving is directed (IV, Preface), and in a
series of propositions (IV, 67---73) he describes the free man as the summit of human
development. But in experience, as we shall see \textsc{Spinoza} himself admit, every
individual being (mode) is both active and passive. Only substance is absolutely active.

We can, however, still use \textsc{Spinoza}'s definition of activity if we alter it thus:
We are the more active, the more the cause of our states and actions lies in ourselves,
--- the more passive, the more their cause lies outside us.

% ===== pp. 38--44 =====
\opage{38}\textit{γ.} For \textsc{Spinoza}'s psychology Ax. 3 is of interest. It states
that such modes of the soul (\textit{modi cogitandi}) as love, desire
(\textit{cupiditas}) and so on presuppose ideas of what is loved, desired and so on, but
that there can be ideas without other modes of the soul being present.

\textsc{Spinoza} is of course right that we cannot love or desire without an idea of what
we love or desire. But he overlooks that there can be elementary feelings, such as
anxiety, well-being, low spirits, and besides these urge and striving, without there being
any idea present of the object of (the cause of) these states. He himself acknowledges this
in the scholium to the ninth proposition of the third book, where he distinguishes between
the blind urge of self-preservation (\textit{appetitus}) and desire (\textit{cupiditas}),
the latter of which presupposes an idea of an object. Indeed, he even adds that we do not
need something or will it because it is good, but conversely hold it good because we need
it and will it. A refraction appears here between intellectualism and voluntarism, which is
characteristic of \textsc{Spinoza}'s psychology and to which we shall come back later.
Here it need only be added that it is in the highest degree doubtful whether
\textsc{Spinoza} is right that there can be ideas without any degree whatever of feeling or
urge.

The fifth axiom, that we know no things other than bodies and modes of thought, we have
already discussed above.

\subsection*{b. Propositions 1---10.}
\begin{center}\textit{The Relation between the Attributes.}\end{center}

\textit{α.} To grasp rightly \textsc{Spinoza}'s teaching on the relation between the
attributes, one must hold fast that in \textsc{Spinoza}, even more than in other
philosophers before the eighteenth century, the definite distinction between different
philosophical problems that is now gradually making its way is lacking. Religion, view of
the world (cosmology), theory of knowledge, ethics, psychology --- all the endeavours and
tasks denoted by these words are so entwined in \textsc{Spinoza} that it can be difficult
to disentangle them, and that an attempt to do so easily leads to reading too modern a
line of thought into this thinker of the Renaissance. As for the relation between the
attributes, it is above all the difference between theory of knowledge and psychology (or
psychophysics) that we have to deal with.

What \textsc{Spinoza} puts forward in the first six propositions of the second book about
the mutual relation of the attributes is a consequence of the second and fourth
definitions of the first book, a consequence already drawn in the tenth proposition of the
first book. If every attribute is to express the essence of substance, and if the essence
of substance consists in having its ground (or cause) in itself, then every attribute must
also contain the ground (or cause) of all the things that fall under it, and no explanation
can be fetched from other attributes.

But the question then lies near, what led \textsc{Spinoza} to set up a definition of the
concept of attribute that entails such consequences. Of the \opage{39}motives of thought
that asserted themselves on this point we have already spoken in the Introduction (see
above). The one motive is scientific and is due to the conception of the method and
principles of the new science that \textsc{Spinoza} had won in the interval between the
writing of the Short Treatise and the working out of the Ethica, and that demanded that
every material thing be explained by another material thing which, like it, could be shown
in experience. The other motive belongs to the theory of knowledge. The criterion of truth
proves in the end to consist in the mutual agreement of thoughts and therefore demands that
every thought be derived from another thought.

These two motives of thought have in \textsc{Spinoza} mutually lit up each other. For the
continuity in principle that natural science found or demanded in the material domain, and
that had already led \textsc{Descartes} to conceive everything material as a single
substance, \textsc{Spinoza} naturally sought a counterpart in the spiritual domain. But
from II, 1. Schol. and 6. Cor. one sees that the domain of thought itself offered him a
lawful connection that could serve as a model for research in the material domain. In 1.
Schol. it is shown how series of thoughts can be built that can be continued to infinity by
the same law, and in Proposition 2 this is then transferred by analogy to the attribute of
extension. In 6. Cor. it says that in the same way and with the same necessity as thoughts
follow from the attribute of thought, the objects (\textit{res ideatæ}) follow and are
inferred from their different attributes. Here, then, one infers from the inner connection
within the attribute of thought to a corresponding connection within other attributes. The
only one of these other attributes that experience makes known to us (Ax. 5) is extension.
The transfer of the concept of series formation by immanent laws from the attribute of
thought to the attribute of extension meant for \textsc{Spinoza}, who, following
\textsc{Descartes}\footnote{\textit{Principia Philosophiæ}, II, 10---11. (Space and body
differ only as genus and species, which shows itself clearly in this, that even if we think
of all sensible properties as diminishing in degree, so that they vanish for us, space still
remains).} in this, found the essence of bodies exhausted in purely geometrical properties
(see Def. 1), that this concept must hold for all bodies as well. That all bodies are at
rest or in faster or slower motion he maintains later in the book (in the two axioms after
Proposition 13) as a fact of experience, whose relation to the definition of corporeality
as extension he does not examine; in his last years, however, he expressed doubt about the
correctness of the Cartesian definition (see Ep. 81 and 83). --- That there is such a
corresponding connection within every attribute is, on the interpretation of
\textsc{Spinoza}'s fundamental concepts set out in the foregoing, what \textsc{Spinoza}
denoted by saying that every attribute in its own way expresses the essence of substance.

A third motive too has certainly been at work, namely the decided rejection of all
teleology that \textsc{Spinoza} pronounces, not only in the appendix to the first book, but
in several places in the course of the book. His criticism of the theological concept of
God and his conception of scientific method worked in the same direction here. All
teleology presupposes that a thought or plan is the cause of changes in nature; but that
\opage{40}would presuppose, partly that there should be no inner lawful connection in
nature, partly that there should be a difference in God's being between thought and power.

In \textsc{Spinoza} himself, however, no sharp distinction certainly made itself felt
between these different motives of thought, which we have here disentangled one by one.

\textit{β.} In Proposition 7, which states that the order and connection of ideas is the
same as the order and connection of things, the hypothesis of identity is stated as holding
for the attribute of thought in relation to all other attributes. For by things must here be
understood all things within all these other attributes. What is meant by ideas is seen from
the demonstration, which refers to the fourth axiom of the first book: ``The concept
(\textit{idea}) of the effect depends on the concept (\textit{idea}) of the cause and
presupposes it.'' It is clear that the ideas spoken of here are scientifically formed
thoughts with their clear and definite content --- what \textsc{Spinoza} (see above p. 36)
calls \textit{esse formale idearum}\footnote{\textit{Formalis} means in \textsc{Spinoza}
what we call objective. And just as we speak not only of objective objects as opposed to
mere thoughts, but also of objective thoughts as opposed to associations of ideas, fancies
and feelings, so \textsc{Spinoza} too speaks both of \textit{essentia formalis attributorum
dei} (Eth. II, 40 Schol. 2), that is, of objects, and of \textit{esse formale idearum}
(Eth. II, 5), that is, of the content of thought. --- With regard to the terminology it
must further be noted that ``objective'' (7 Cor.) means what we would call
``subjectively'', since it is used as the opposite of \textit{formaliter}. See especially
I, 30 Dem., where ``objective'' stands in opposition to ``in Natura''.}. It is to this
clarified content of thought that ``things'' correspond, the mutual dependence of things
corresponding to the dependence of the consequence (the effect) on the ground (the cause).
One sees clearly here how substance finds expression in the different attributes: the
relation between ground and consequence is presupposed in them all, is what is identical in
them all (however many they may be).

In the scholium to the seventh proposition the hypothesis of identity is applied specially
to the relation between the two attributes we know from experience. ``The thinking and the
extended substance'', it says, ``are one and the same substance, which is now regarded under
this attribute, now under that. And so also a mode of extension and the idea of this mode
are one and the same thing, only expressed in two ways\dots{} So long as we regard things
as modes of thought, we must explain the whole order of nature, or the connection of
causes, by the attribute of thought alone, and so long as we regard things as modes of
extension, the whole order of nature must be explained by the attribute of extension alone.
And the same holds, in my view, of the other attributes. God is the cause of things as they
are in themselves\footnote{The expression ``things as they are in themselves'' occurs in
\textsc{Spinoza} in different connections and therefore in different senses. 1) In II, 7.
Schol. it denotes things in so far as they are regarded under all possible attributes
(which human thought cannot do). 2) In II, 44 Dem. (with reference to I, 29) it denotes
things in their necessity as opposed to the way in which they stand for mere observation.
3) In the remarkable letter (Ep. 6) in which \textsc{Spinoza} discusses some experiments of
\textsc{Boyle}'s, concepts such as motion, rest and their laws (these ``chaste concepts'',
\textit{notiones castæ}) are said to express nature as it is in itself, as opposed to ideas
(such as colour, cold, heat, solidity, fluidity and so on), which express only their relation
to human sensation. In this sense the word is also used by \textsc{Locke}, who (Essay II,
8, 23) says of ``the primary qualities'' (figure, number, motion and so on) that they are in
things whether we perceive them or not, so that they give us ``the idea of the thing, as it
is in itself''. This sense comes from \textsc{Galileo}, \textsc{Hobbes} and
\textsc{Descartes}. 4) Finally the expression ``\textit{res in se consideratæ}'' is used in
the preface to the fourth book (see also IV, 73 Schol.) to denote the opposite of the
concepts of value one is inclined to apply to things.}, because he consists of infinitely
many attributes.''

\opage{41}As the proposition is set up here, it is of the kind that belongs to the theory
of knowledge and concerns the relation between scientifically formed thoughts (objective
concepts) and other sides of existence. This is seen not only from the reference to the
\textsc{Kabbala}'s teaching of God, whose understanding and what he understands are one and
the same; it is seen also from \textsc{Spinoza}'s example: the relation between the circle
existing in nature and the concept of the circle.

But as is seen in what follows (Propositions 11---13 and III, 2. Schol.), \textsc{Spinoza}
regards the relation between soul and body as a special case of the relation between
knowledge (that is: scientific concepts) and its objects, and so passes, without noticing
it, from theory of knowledge to psychology (or psychophysics). The relation between the
mathematical concept of a circle and the circle existing in nature is after all a different
one from the relation between the concept as a psychical product and the physiological
process to which it is in some way or other bound, and which must fall under the attribute of
extension just as much as ``the existing circle''. --- The confusion becomes clearer still
when we ask how it stands with the activity of thought that leads to the forming of
scientific (objective, ``formal'') concepts. Does not something within the attribute of
extension correspond to this too? And the matter becomes clearest of all when we ask how it
stands with wrong formations of concepts and self-contradictory ideas. One can suppose that
certain physiological (cerebral) processes correspond to these too. But \textsc{Spinoza}
quite certainly does not suppose that anything whatever that ``exists in nature'' corresponds
to the idea ``square circle''.

It is \textsc{Spinoza}'s dogmatism that led him to confuse two problems. Like all dogmatic
philosophers and natural scientists he is certain that absolute realities must correspond
to our concepts formed by strict method (see above p. 16 f). And his strict conception of the
inner connection of material nature made it impossible for him to assume an action of the
one attribute on the other; states of the soul therefore became for him another expression
of certain material states, and conversely. These two relations he has identified. And so
now his dogmatic theory of knowledge, now his realist psychophysics has the word. --- We
shall have occasion several times more to speak of this question.

\textit{γ.} In Propositions 8---10 it is enjoined that the concept of substance cannot be
used of a man, as scholasticism and in part even \textsc{Descartes} had taught on account of
a vaguer concept of substance. Man is an individual being, a mode, and is as such a member
of a series of individual beings. He has his cause in God, not in so far as God is regarded
as an infinite being, but in so far as God's being shows itself in the lawful connection of
the individual beings. In the divine attributes there are contained possibilities of
infinitely many individual beings, just as in a circle there is the possibility of an
infinite number of chords cutting one another, whose \opage{42}segments form equal
rectangles, whereas in each individual case only a finite number can be given. In this
comparison, which is found in the scholium to the eighth proposition, the circle corresponds
to substance, the chords to the modes.

\subsection*{c. Propositions 11---15.}
\begin{center}\textit{On the Relation between Soul and Body.}\end{center}

\textit{α.} The general relation between thought and extension, which in the foregoing has
been determined as a consequence of the relation between attributes in general, is now
applied specially to the relation between soul and body. These, then, relate as the idea
(\textit{idea}) relates to its object (\textit{res}). The body (\textit{corpus}) is the
individual thing (\textit{res singularis}) that corresponds to consciousness
(\textit{mens}).

When it says in Proposition 13: ``The object of the idea that constitutes the human
consciousness (\textit{mens}) is a body, a certain mode of extension, and nothing else'',
the meaning must be that man's peculiarity consists in a mode of extension corresponding to
his thought, and not a mode of one of the other attributes. In the demonstration given of
this proposition reference is made to Ax. 5, which lays down that we feel or perceive no
individual things other than bodies and modes of thought. --- Now \textsc{Spinoza} does
assume, after all, that there are infinitely many attributes, all standing in a proportional
relation to one another. What holds of human consciousness in relation to the attribute of
extension must therefore consistently hold of it in relation to any other attribute
whatever, although here the one member of the relation is $x$. (Cf. above p. 22---24).

What holds of man must (13 Schol.) hold also of all individual beings within our experience,
by virtue of the general relation between the attribute of thought and that of extension.
All individual beings must be animate (\textit{animata}), although in different degrees.
The difference between the different beings with souls corresponds to the difference
between the different bodies, which above all offer highly different relations between
activity and passivity. Already as regards man himself there are great fundamental
differences in this respect. The more active and independent the body is, that is, the more
its functions depend on itself, the more perfect will the knowledge of the soul also be.
Obscurity and clarity in the life of thought correspond to passivity and activity --- in the
case of soul and body alike. By the obscure ideas \textsc{Spinoza} has in mind, as far as one
can see, sensations of life and elementary feelings of pleasure and pain and mere
associations of ideas, perhaps also the obscure urge that is not aware of its aim. According
to Def. 3 the word \textit{idea} really denotes only actively formed thoughts; but in Ax. 3
a distinction is made between such thoughts and other modes of thought (\textit{amor,
cupiditas} and other \textit{affectus animi}). As already remarked (see above p. 38), a
difficulty appears here for \textsc{Spinoza}'s psychology, since these passive elements
are always to presuppose the active ones, but not conversely.

\textit{β.} After Proposition 13 \textsc{Spinoza} inserts a series of borrowed propositions
(\textit{lemmata}) \opage{43}from physics, naturally as he understood the physics of his
time. After first laying down that no body is a substance, he sets up the principle of
inertia (Lemma 3), which has already been mentioned above (p. 9 f). Next it is maintained
that from motion follows motion, from rest only rest. --- Here one clearly sees the main
difficulty in the Cartesian physics, which was first removed by \textsc{Leibniz}, namely how
a thing at rest can come into motion, or a thing in motion come to rest, when rest and motion
are two wholly different states. \textsc{Hobbes}, for that matter, following in
\textsc{Galileo}'s track, had already indicated the possibility of a solution through the
concept of the tendency to motion (\textit{conatus}), which does not vanish because it is
checked by an opposite tendency to motion\footnote{\textit{Den nyere Filosofis
Historie\textsuperscript{2}}, I, p. 268. --- \textsc{Gassendi} too corrected
\textsc{Descartes} in this way (ib. p. 254).}.

Then it is presented purely dogmatically how the simple elements of which bodies consist
bring about, by their different relations of position, the difference between the states of
aggregation, between solid, soft and fluid bodies. Several bodies can in turn be joined into
a whole through common working, a whole that can persist even if the individual parts
change. From such wholes more composite individual wholes can arise in turn, and in the end
the whole of nature makes up one great individual, whose parts, the individual bodies, are
changed in many ways without change in nature as a whole. --- What \textsc{Spinoza} here
calls ``the whole of nature'' is plainly the same as what in other places (see above p. 29)
is called ``the character of the whole universe'' (the physiognomy of the world), and which
through the constant relation between rest and motion has its ground in the attribute of
extension.

That the individual changes in the world (\textit{variatio rerum}) cannot be derived from
the attribute of extension, and that \textsc{Descartes} wrongly defines matter as mere
extension, \textsc{Spinoza} admits in the last letter (Ep. 83) we have from his hand, and he
announces a closer discussion that he never came to undertake.

Finally it is laid down (in the form of postulates) that the human body is composed of
smaller wholes that are different from one another and each composite, --- that the action
of other bodies can leave traces (\textit{vestigia}), which consist in one of the fluid
parts of the body affected changing a soft part by altering its position, --- and that our
body constantly needs to be regenerated by the action of other bodies. And from this
composite character of the human body it is inferred (Proposition 15), by virtue of the
proportionality between soul and body, that the idea that constitutes the essence of the
human consciousness is composed of very many ideas.

\textit{γ.} \textsc{Spinoza} applies his theory of the relation between soul and body in
the special parts of the Ethica in such a way that he now infers from the knowledge we have
(or that he thinks we have) of the body what goes on in the soul, now from what is known to
go on in the soul infers what goes on in the body. Both are justified, since according to
\textsc{Spinoza} the relation between soul and body is a purely logical or mathematical
relation of functions, without either of the members being declared once \opage{44}and for
all to be the independent variable. All attributes are simultaneous expressions of the
essence of substance (\textit{simul natura}. I, 17 Schol., III, 2 Schol.). It has wrongly
been found inconsistent that \textsc{Spinoza} goes now the one way, now the other. He goes
from the corporeal to the spiritual in Propositions 16---19; the opposite way he goes e.g.
in III, 12 Dem. and especially V, 1\footnote{It is of interest to compare
\textsc{Avenarius}' teaching of what he calls the two vital series (the physiological and
the psychological) with \textsc{Spinoza}'s teaching of the attributes. See \textit{Moderne
Filosofer}, p. 101 f. --- \textsc{Avenarius} declares the physiological vital series in
advance to be the independent variable.}.

Concerning the relation between the two series of things \textsc{Spinoza} uses either the
expression ``indicate'' or ``point to'' (\textit{indicare}) (II, 16 Cor. 2; III, 14 Dem.) or
the expression ``presuppose'' (\textit{involvere}) (II, 16; III, 29 Dem.). The first
expression describes the relation as symptomatic; the second, which is more correct,
describes it as a purely logical relation of inference, a relation between a ground of
knowledge and what is derived from it. Yet \textsc{Spinoza} also often uses the expression
that the soul ``perceives'' (\textit{percipit}) the body, an expression bound up with the
confusion mentioned above of theory of knowledge and psychophysics, since the relation
between soul and body is regarded as a relation between subject and object.

It is perhaps not superfluous to explain the difference between the two points of view by
an example. The physicist infers from a sensation of colour a certain degree of refraction
of white light. That is a relation belonging to the theory of knowledge. The physiologist
infers from the same sensation of colour certain processes in the sense organ and the brain.
This is a psychophysical relation. In every sensation both points of view assert themselves.
And likewise in every thought, every feeling, every act of will. When \textsc{Spinoza}
thought the concept of substance, a principle of the theory of knowledge was to be expressed
by it; but at the same time something went on in his brain (if only we knew what it was!). A
heroic resolution (what \textsc{Spinoza} calls \textit{mentis decretum} as a
\textit{species fortitudinis}, see III, 2 Schol. and 59 Cor.) concerns different
relations in and outside the resolving subject, and on this rests its ethical significance;
but from another point of view such a resolution is a physiological phenomenon (according to
\textsc{Loeb} a purely chemical phenomenon), which in the end must be derivable ``from the
laws of motion and rest''. Which point of view one adopts depends of course on the context
in the individual case.

If one wanted to say that thought (consciousness) could itself be a material process,
\textsc{Spinoza} would answer as he answered in a letter (Ep. 4): ``Even if I grant this,
which I by no means do, it cannot be denied that extension, considered as extension, is not
thought''\footnote{In the treatise on the improvement of the understanding
\textsc{Spinoza} gives among examples of unclear and false ideas the assumption that there
are bodies from whose mere composition thought arises (\textit{dari corpora, ex quorum sola
compositione fiat intellectus}).}. --- Apropos of this remark it may be added that the point
of view of the theory of knowledge must find application even if one keeps to the
physiological side of man's being. The problem then becomes why and how one can regard a
certain state of the brain as a sign of something that perhaps has a wholly different
physical character from it, perhaps \opage{45}relating to it as an optical process to a
chemical one. \textsc{Spinoza} himself gives hints in that direction in what immediately
follows.

\subsection*{d. Propositions 16---36.}
\begin{center}\textit{The Psychology of Ideas.}\end{center}

\textit{α.} \textsc{Spinoza} introduces the doctrine of ideas from the attribute of
extension. Our body stands in constant interaction with other bodies, and its state is
therefore at every moment determined not only by its own nature, but also by the character
of the other bodies that act on it. The soul's perception (\textit{perceptio}) of its own
body therefore also contains a perception of other bodies, but conditioned by the way in
which one's own body has been affected by them. It regards them, so to speak, from the
standpoint of one's own body. Its knowledge of other bodies therefore cannot become adequate
(\textit{adæqvata}) (25), but stands as a consequence without premises, --- and the same
holds of the soul's knowledge of its own body, since this is always determined by its
relation to other bodies (28). --- What \textsc{Spinoza} puts forward here is what is now
called the subjectivity of the sense qualities. He puts it forward in a less dogmatic way
than one would expect from his doctrine of the attributes.

\textit{β.} We often perceive several other bodies at once, and when later we perceive one
of them again, we come to remember the others. For here the traces (\textit{vestigia}) left
by the first perception are at work\footnote{These traces are material (III, Postulate 5).
Sometimes the material processes that correspond to memories and fancies are called images
(\textit{imagines}), since things stand as present. \textsc{Spinoza} apologizes for this
expression, because those material processes do not reproduce the look of things (II, 17
Schol.). In III, Postulate 2 a distinction is made between traces and images, and the latter
are regarded as a consequence of the former (though as material!). Cf. also V, 11---14. ---
But in II, 49 Schol. ``image'' is used psychologically.} (18). --- It is the so-called
association by contiguity that \textsc{Spinoza} speaks of here. Elsewhere (Tract. Theol.
polit. cap. 4) he gives association by similarity as an independent form of association,
and in his explanation of pity (III, 27) he in any case presupposes immediate recognition as
a condition. When he says in IV, 18 Schol. that a word can call up the idea of an object
although there is no similarity between them, this really implies the presupposition that
the association of ideas goes more easily when there is similarity between the members.

The association of ideas takes place spontaneously, and from the first we have no reason to
doubt the validity of the ideas that turn up in us. What we imagine we regard provisionally
as valid, until other ideas come that exclude its validity. In themselves our ideas contain
no error; we lack only an idea that could exclude the validity of the ideas before us (17).
Just as \textsc{Spinoza} has called the sensations consequences without premises, so on his
view he could also have used this designation of ideas that arise spontaneously.

\opage{46}This interesting way of looking at things was later carried further by English
psychology, which connected it partly with a spontaneous tendency to continue a series of
ideas until definite obstacles come (\textsc{Hume}), partly with the spontaneous urge to
movement that makes us act without criticism on ideas that turn up spontaneously, until this
acting brings disappointments or sufferings (\textsc{Bain}).

\textit{γ.} \textsc{Spinoza} then distinguishes very sharply between mere memory or
association of ideas and the connection of our ideas that takes place according to the rules
of the understanding (\textit{secundum ordinem intellectus}), and that leads to an
understanding of things from their first causes, independently of the changes of state that
differ for each individual (18 Schol.). True knowledge must comprise the laws that hold for
all bodies and for all souls. Human consciousness has only an imperfect knowledge of itself,
of its body and of external things, so long as man is determined to his ideas of things from
outside, namely by chance encounter. It is otherwise when the ideas are determined by an
examination of the agreements, differences or oppositions of things (29 Schol.).

In substance (God) there is a perfect knowledge (\textit{idea seu cognitio}), because it
acts at once under all attributes and in all modes, not merely under a single attribute or
in a single mode (19---32). Untruth is due to limitation and incompleteness and the confusion
that follows from them. An unclear sensation or idea (\textit{idea confusa}) is, as we have
seen, a consequence without premises (28 Dem.); if the premises can be supplied, the
unclarity ceases. But these imperfect and obscure ideas arise with the same necessity as the
perfect, clear and true ideas --- and for God, to whom this necessity stands clear, no obscure
thoughts therefore exist. The obscurity is present only in so far as the ideas are referred
to a single and limited consciousness (36).

A difficulty has been found in \textsc{Spinoza}'s ascribing to human consciousness an idea
of itself, so that it can be an object for itself, just as the body is an object for it
(21). \textsc{Spinoza} has himself given the necessary explanation, saying that this idea
of our consciousness (\textit{idea mentis}) --- which, since consciousness is itself an
idea, is an idea of an idea (\textit{idea ideæ}) --- is nothing but the form of
consciousness or of the idea, regarded as a mode of thought, apart from the relation to the
body (21 Schol.). This corresponds to what is said in the treatise on the improvement of
the understanding, that the method of thinking is a reflexive knowing (\textit{cognitio
reflexiva}), by which we become aware of the form and procedure of our thoughts. Only he who
has true thoughts can know what truth is. --- When \textsc{Spinoza} says that this idea of
ourselves relates to our consciousness (our soul) in the same way (\textit{eodem modo}) as
this relates to the body, ``the same way'' naturally means only analogy, not identity.
There is nothing to prevent there being, within one and the same attribute, relations
analogous to the relation between two attributes. --- On the psychophysical meaning of the
attributes, \opage{47}the task in the case of this self-consciousness, exactly as in the case
of all other phenomena of consciousness, will be to find what brain process goes on when it
appears.

\subsection*{Propositions 37---47.}
\begin{center}\textit{Theory of Knowledge.}\end{center}
% the print gives this section no letter (e.)

\textit{α.} It is the natural science of the time that gives \textsc{Spinoza} his
understanding of what science is at all. He starts from all bodies agreeing in certain
properties, above all in being subject to the laws of motion and rest and in the end all
falling under a single attribute, the attribute of extension. It is the general mathematical
and mechanical laws he has in mind. Their content is something that appears in equal degree
in existence as a whole and in its individual parts (37). --- The concept of the whole is in
\textsc{Spinoza} derived in relation to the concept of conformity to law. As we have seen
(above p. 29 and 43), he conceives the universe as an individual whole, and this conception
is grounded in there being in the universe a constant relation between rest and motion
despite all changes in the particular\footnote{The expression ``that which is equally in the
part and in the whole'' (\textit{id, quod æqve in parte et in toto est}) has its history.
\textsc{Spinoza} may have it from \textsc{Bruno}, who teaches that ``the world-soul'' (which
is the principle of nature's conformity to law and is also called \textit{forma
universalis}) is ``tutta in tutto et qualsivogla parte del tutto'' (\textit{Op. ital.} ed.
Lagarde p. 203, 242). (Likewise \textit{De Immenso} IV, 15). But he may also have it from the
scholastics. It goes back to \textsc{Augustine} and \textsc{Athanasius}, who surely have it
from \textsc{Plotinus}. See on this my \textit{Religionsfilosofi}, Note 15. A hint is found
already in \textsc{Plato} (\textit{Parmenides} p. 131 B). --- \textsc{Hume}
(\textit{Treatise} I, 4, 5) speaks of the expression as scholastic and finds it ``shocking
when crudely propos'd''. He is probably thinking of the mythological and dogmatic use of the
concept. But its history is interesting, because it shows that the thought of the conformity
to law of existence presses forward as soon as men begin to think. This thought runs through
different forms. In \textsc{Spinoza} it receives a decidedly rational form, although it is
put forward dogmatically.}.

The knowledge that concerns the general laws of existence is thus not conditioned by any
particular thing in its difference from other particular things. Or, as \textsc{Spinoza}
puts it in his metaphysical-theological language, God's being appears in the same way
whether he is regarded as constituting the essence of man or as constituting the essence of
another thing. % the Danish misprints fremtrædrr (appears)
For what is in question here is what man's body has in common with all other bodies, and
therefore man can here have an adequate knowledge (37---39). We stand here before what
\textsc{Spinoza} in the treatise on the improvement of the understanding called ``the fixed
and eternal things and the laws contained in them'', which correspond to what the Ethica
calls infinite modes (of the second degree); and the last foundation here is the attribute
of extension.

At this point theory of knowledge and psychology coincide. If the life of the soul had no
other content than that of scientific knowledge of law, there would therefore be no
opposition between the interpretation of the doctrine of attributes in terms of the theory
of knowledge and its psychophysical interpretation. But \textsc{Spinoza} himself recognizes
that there is other content of the soul than scientific knowledge of law. Partly he
maintains, as we shall see, that the knowledge of law is not the highest form of science,
because it does not \opage{48}take in the individual phenomena in their peculiarity. Partly
he maintains the reality of error, that is, the fact that there are ideas that conflict
with what scientific knowledge establishes. Partly, finally, he assumes feelings and desires
that have a certain independence in relation to mere thought. And to all these irrational
elements (since they all belong after all to the attribute of ``thought'') elements within
the attribute of extension must correspond. Here the psychophysical interpretation shows
itself in its opposition to the interpretation in terms of the theory of knowledge. --- But
we return to the theory of knowledge.

It is the general laws (\textit{notiones communes}) that form the foundation of all thinking
about nature. \textit{Fundamenta ratiocinii nostri} they are called (40 Schol. 1). They
disregard all differences of time and can therefore be said to show us things under the
aspect of eternity (\textit{sub qvadam æternitatis specie}, 44 Cor. 2). By eternity
\textsc{Spinoza}, after all, means not continual persistence, but an existence that follows
with necessity from the essence of things (I, Def. 8). And this necessity of things is
(according to I, 16) one with the necessity of God's eternal nature (44 Cor. 2).

The expression \textit{notiones communes} comes from the Stoics, and it originally denotes
concepts that are common to all men\footnote{Cf. \textsc{Paul Barth}: Die Stoa.
(Frommanns Klassiker der Philosophie). 1903. p. 72 f.}. \textsc{Spinoza} too assumes that
there are certain concepts that are common to all men and that are (in part) perceived by
all ``adequately, that is, clearly and distinctly''; but this (possible) community is
grounded in him by inference from these concepts concerning only what holds ``equally of the
whole and of every part'' (38 with Cor.).

On the way by which one reaches the knowledge of the general laws on which we must rely in
our understanding of existence \textsc{Spinoza} will not speak more closely, but refers to
the treatise on the improvement of the understanding, which unfortunately, as we know, he
did not finish. He has, however, at an earlier place in the Ethica (29 Schol.) given a small
hint, in that he sets over against the ``external'' perception, which is due only to the
encounter of chance circumstances, the knowledge that can be called ``inner'', because it
rests on comparison and thereby leads to an understanding of agreements, differences and
oppositions; by such a procedure one arrives at a clear and distinct perception of things.
(Cf. above p. 46). It is, after all, also on the possibility of ordering our thoughts in
series by definite laws that he grounds the concept of attribute. (See above p. 8 and p.
39). Of interest in this connection is also his exchange with the famous chemist
\textsc{Boyle} (with \textsc{Oldenburg} as intermediary). Against \textsc{Boyle} he
maintains that no experiment can establish the justification of the mechanical view of
nature, which leads all sensible properties back to motion and rest and their laws (Ep. 6).
It seems to be \textsc{Spinoza}'s meaning that every experiment can only supply an example
of these laws, but that these stand as the constant presupposition, which according to
\textsc{Spinoza} is an eternal and necessary truth. Despite the dogmatic line of thought,
his position here nevertheless offers an analogy to \textsc{Kant}'s. He stood towards the
natural science of his time in a way similar to that in which \textsc{Kant} stood towards
\textsc{Newton}'s Principia, regarding \opage{49}it as a spiritual fact whose foundation was
to be found by way of analysis. The analogy lies all the nearer as there is, on this point,
still a dogmatic element in the founder of criticism himself\footnote{Cf.
\textit{Totalitet som Kategori}, p. 36 f.}.

From the concepts of general, pervasive laws (\textit{notiones communes}) \textsc{Spinoza}
distinguishes such concepts or words as have general meaning only because the ability to
comprehend different things in a distinct way has been lacking. To this belong such words as
``thing'', ``being'', ``something'', and besides all concepts of kinds (\textit{notiones
universales}) such as ``man'', ``horse'', ``dog'' and so on (40 Schol. 2)\footnote{It is of
interest to compare \textsc{Spinoza}'s criticism of general concepts with \textsc{Leibniz}'.
For \textsc{Spinoza} the main point was that general concepts do not state real laws for the
relation between the different marks or properties. When \textsc{Leibniz} doubts the
validity of general concepts, it is because he is convinced that ``nature has differences and
likenesses that we do not know''. Therefore our determinations of the natural kinds must
always be provisional and conditioned by our knowledge at a given time. Besides, the
transition from kind to kind can be continuous, for which reason concepts of kinds can only be
set up purely by convention. (\textit{Nouveaux Essais}. III, 6, 14. 23. 28).}. This
distinction recalls how \textsc{Bruno} and \textsc{Bacon} interpreted \textsc{Plato}'s
``ideas'' as laws, not as mere general concepts, which is what they originally were.

Among the mere general concepts \textsc{Spinoza} also counts the concept of number. ``We
do not perceive things'', he says (Ep. 50), ``under the concept of number except after
having led them back to a common kind.'' I can say that I have two coins when I have a
daler and a skilling. And since concepts of kinds, as we have seen, are according to
\textsc{Spinoza} signs of an imperfect power of thought, he cannot give the concept of number
any eminent place in thought. It cannot be applied to substance. ``Measure
(\textit{mensura}), time and number are only modes of thinking or rather of imagining
(\textit{imaginandi modi})'' (Ep. 12). It is only arithmetic, not geometry, that
\textsc{Spinoza} has misgivings about. The difference he makes here springs, no doubt, from
continuity appearing more clearly in the geometrical than in the arithmetical domain; number
is the expression of discontinuity. But from \textsc{Descartes}' analytic geometry (which
was in his library) he could have learned that geometrical forms can be expressed by
arithmetical formulas (the circle, for example, which according to \textsc{Spinoza} ``exists
in nature'', by the equation $r^2 = y^2 + x^2$), so that number and extension can go
together, the formulas changing as the lines and their relations change, and conversely.
Perhaps, however, \textsc{Descartes}' analytic geometry influenced \textsc{Spinoza}'s
thought in another way, in that the relation between the different attributes is of just
the same kind as, according to \textsc{Descartes}, the relation between equation and
curve\footnote{See the interesting discussion in L. \textsc{Brunschwicg}: \textit{Les
Étapes de la philosophie mathématique}. p. 144 f.}. % the Danish misprints BBUNSCHWICG (Brunschwicg, as on p. 53)
But the qualitative discontinuity, which in general brings such great difficulties for
\textsc{Spinoza}'s system, he could not acknowledge as belonging to thought, and therefore
he refers it to mere ``imagination''.

It was first \textsc{Leibniz} who saw that the origin of number is due to the ability to
continue a series by the same law by which it began, and that this
con\opage{50}tinuation was always in our power. It is the positive significance of repetition
that appears here and marks the boundary between logic and mathematics. A new phase in the
history of the problem of knowledge begins with \textsc{Kant}'s taking up this point of view
as decisive, if rational empirical science is to be seen in its grounding and
justification\footnote{Cf. \textit{Totalitet som Kategori}, p. 34---37.}.

Another, and more justified, use of his critical attitude towards concepts of kinds
\textsc{Spinoza} makes in his criticism of the assumption of a ``free will'' in the sense of
the word in which everyone is supposed to have the same power of self-control. This
assumption is due, he thinks, to an abstract concept of men and overlooks the great
difference between men's abilities and between the ways by which individuals must work
their way forward to perfection. (Ep. 19 and 23).

\textit{β.} It is not \textsc{Spinoza}'s meaning that only material nature is subject to
general laws. In the teaching on the mutual relation of the attributes there lies hidden the
presupposition that the spiritual cannot be derived from the material and must therefore be
regarded as a peculiar side, a special attribute of substance. What is common to the
different attributes is, after all, precisely the relation between ground and consequence,
the relation that appears in different forms within the different attributes. All
individual things (no matter under which attribute they are referred) are subject to a world
order that determines their existence at a given time and place, and that has in the end its
ground in the necessity of the divine nature. (II, 45 Schol. and V, 29 Schol.). As for the
laws that hold under the attribute of thought, one could in \textsc{Spinoza} --- owing to the
confusion of the points of view of the theory of knowledge and of psychology shown above ---
just as well expect a reference to logical-mathematical as to empirical-psychological
connections. Whereas, where there is talk of the possibility of always continued series
formation under the attribute of thought (Eth. II, 1), one would most readily think of
logical-mathematical series of concepts (see above p. 9 Note), \textsc{Spinoza} himself,
where (Tr. theol. polit. cap. 4) he has to name examples of conformity to law in nature,
points to the law of the association of ideas as the correlate of the laws of motion in the
material domain. And in the third and fourth books, and the beginning of the fifth, he shows
how, through the association of ideas, displacements can take place in the domain of the
life of feeling and will, and he infers from these phenomena to correlated material
processes (V, 1). And on the other hand he finds in the spiritual domain a counterpart to
the principle of inertia holding in the material domain (III, 6) in the urge of
self-preservation, which plays so great a part in his ethics. It was in general his
conviction that ``whether one regards nature under the one or the other attribute, it is one
and the same order or connection of things one has before one'' (III, 2 with reference to
II, 7 Schol.). And on this foundation he builds, in the last three books of the Ethica, his
teaching on the highest development of human life.

\textit{γ.} As for the actual development of knowledge, \textsc{Spinoza} distinguishes
three stages (II, 40 Schol. 2). On the first stage knowledge is conditioned by individual,
scattered observations, which stand as conclusions without premises. (Cf. above p. 45).
\opage{51}In the treatise on the improvement of the understanding a closer characterization
of them is given. Their scattered and chance appearance (\textit{sensationes fortuitæ et
solutæ}) is stressed and connected with their not springing from the energy of
consciousness itself, but being due to external causes. Scatteredness and passivity are thus
closely bound up together. It does not matter, \textsc{Spinoza} adds, how one thinks of the
sensuous ideas and their origin\footnote{\textsc{Hume} too (\textit{Treatise} I, 3, 5)
leaves it open [for the theory of knowledge] how the sensations arise. Had \textsc{Kant}
taken a similarly cautious position, the concept of the Ding an sich would hardly have been
formed.}; the decisive thing is their vague and passive character. With these scattered and
passive observations \textsc{Spinoza} puts the ideas that are due to report and tradition,
and that are formed on the likeness of what we ourselves have seen and heard. This first
stage of knowledge \textsc{Spinoza} calls vague experience (\textit{experientia vaga}),
opinion (\textit{opinio}) or imagination (\textit{imaginatio}). % the Danish misprints expeenritia
The second stage is the scientific standpoint just described, which leads to the knowledge
of general laws and relies on the critical comparison of observations and reports (29
Schol.). Besides natural science and psychology, \textsc{Spinoza} could count historical
criticism here, as he himself applied it to the books of the Bible in his
theologico-political treatise. This second stage \textsc{Spinoza} calls reason
(\textit{ratio}). Already here we know things ``under the aspect of eternity'', because we
rely on what is independent of times, places and individual differences (44 Cor. 2). The
third stage is designated intuitive knowledge (\textit{scientia intuitiva}).

Intuitive knowledge stands in opposition both to the knowledge that builds on observation
or tradition and therefore concerns only individual things (\textit{singularia}), and to
rational knowledge, which demonstrates general laws. It unites the peculiarity of both in
itself, in that it is a knowledge of the essence of an individual thing through the
understanding of the lawful connection within which the thing stands as a peculiar member.
The example by which \textsc{Spinoza} illustrates this description of intuitive knowledge is
taken from simple arithmetical proportions. We see immediately that 6 is the number that is
to 3 as 2 is to 1. Underlying this, to be sure, is a knowledge of the nature of numbers and
of proportional relations; but this knowledge is not applied with full awareness. In this
example (II, 40 Schol. 2) \textsc{Spinoza} plainly regards the number 6 as an individual, a
\textit{res singularis}, which by virtue of the law of proportionality fits in a definite way
into a whole connection.

This third and highest kind of knowledge can be reached only by one who has gone through the
first two degrees. Insight into the conformity to law of existence is a condition of its
arising. In this the Ethica's concept of intuition differs from the more mystical concept of
intuition in the Short Treatise. In the fifth book of the Ethica intuition is said to consist
in gaining, from the very essence of an individual thing (\textit{ex ipsa essentia rei
cujuscunqve singularis}), insight into its dependence on the substance of existence, as
regards both essence and existence, and it is added that in the first book it had only been
shown how things in general were dependent on substance (V, 36 Schol.). \opage{52}The line of
thought in the first book thus belongs to the second kind of knowledge (\textit{ratio}),
which does show us things under the aspect of eternity, but in an abstract way (I, 25). The
individual being that \textsc{Spinoza} in the fifth book wants to show in its connection with
substance is the human mind, which through the knowledge of the laws of existence stands as a
member of the whole connection that has its last ground in substance through the infinite
understanding and the attribute of thought.

It is perhaps characteristic that some of the examples of intuitive knowledge that
\textsc{Spinoza} gives in the treatise on the improvement of the understanding are lacking in
the Ethica. Thus the example that when I know something, I know from it what it is to know.
This is an analysis, or, if one will, an immediate inference. In the treatise named, method is
defined as knowledge of (consciousness of) knowledge (\textit{cognitio reflexiva}), an idea of
an idea (\textit{idea ideæ}). (See also Ethica II, 20---21; cf. above p. 46). In the Ethica
too it is expressly maintained that whoever has a true idea must also know that he has a true
idea and cannot doubt its truth, since an idea is not something dumb, like a picture
(\textit{pictura}) one can stare at, but is an activity, is the very act of understanding.
Every true idea contains the norm of truth in itself (II, 43). When \textsc{Spinoza} in the
Ethica does not name this as an example of intuition, it is thus not because he does not
acknowledge its correctness. It is rather because he thinks that it is not the kind of
intuition he calls the third and highest degree of knowledge. When the immediate knowledge of
what goes on in an act of knowing is called intuition in the treatise on the improvement of
the understanding, it is what I have called analytic intuition, the consciousness that
expresses itself in every immediate inference, every simple transition from member to member
in a demonstration\footnote{\textit{Om Begrebet Intuition} (Vid. Selsk. Oversigt 1914),
p. 77---80).}. In the Ethica, on the other hand, \textsc{Spinoza} lays weight on an individual
thing (\textit{res singularis}) being known by intuitive knowledge in its connection with
existence as a whole. Here it is not a matter of simple analysis, immediate transition of
thought, but of an apprehension of the whole, a construction, which forms the close of a long
work of thought and builds on the knowledge of law won through this work of thought. This is
what I have called synthetic intuition\footnote{\textit{Om Begrebet Intuition}, p.
80---82.}. I suppose, as indicated above, that \textsc{Spinoza} in his example from numbers
regards a number (here the number six) as an individual, which by an intuition won on the
basis of the law of the formation of the series of numbers gets its place fixed in relation to
other numbers. Only so does the example fit what it is plainly his intention to express. In
the fifth book it is the thinking man himself, let us say \textsc{Spinoza} himself, whose
personality is seen as a member of the connection of the world, as an eternal element of mind
(\textit{æternus cogitandi modus}) (V, 40 Schol.).

Analytic intuition must already be at work in the transition from the first stage of knowledge
to the second. It is applied in the careful comparison that does not stop at scattered
observation, but through examination of the likenesses and differences of things finds the
foundation of our thinking, which \opage{53}consists in concepts of general relations (II,
28---29; 44 Cor.). Analytic intuition is thus a means, a preparation. Synthetic intuition, on
the other hand, which \textsc{Spinoza} sets up in the Ethica, is an apprehension of the whole
that stands as the last aim and great ideal of knowledge\footnote{\textsc{Joachim} (\textit{A
Study of the Ethics of Spinoza}, 1901, p. 183) is right only in this, that intuition is
already presupposed at the second stage of knowledge, if one has analytic intuition in mind.
\textsc{Huan} (\textit{Le dieu de Spinoza}, p. 16---19) has seen this clearly. --- What I
call analytic intuition is already described by \textsc{Descartes} (\textit{Regulæ}, Ch. 8),
and after him by \textsc{Locke} (\textit{Essay} IV, 2, 1 and 10, 3) and \textsc{Hume}
(\textit{Treatise} I, 3, 1---3). \textsc{Descartes} uses examples similar to
\textsc{Spinoza}'s in the treatise on the improvement of the understanding, and
\textsc{Spinoza} may well have known the Regulæ, which are known to have circulated in copies
in \textsc{Spinoza}'s circle (see on this \textsc{Brunschwicg}: \textit{Les Étapes de la
philosophie mathématique}, p. 141).}. What \textsc{Spinoza} says of intuitive knowledge in
the treatise on the improvement of the understanding must be understood of it: ``What I have
so far been able to understand by this kind of knowledge is very little (\textit{perpauca}).''
It is clear, after all, that a wholly transparent, rational connection between the individual
and the universal sets an infinite task for thought. \textsc{Leibniz}, who had so clear an
eye for the individual and nuanced and for the countless transitions, saw this clearly. It is
impossible to know the full individuality of a thing, since every individual being stands in
an infinity of interactions with everything else in the world, and since there is an infinity
of small differences not known to us (Nouveaux Essais, III, 3, 6). It is nevertheless of great
interest that precisely \textsc{Spinoza} set up synthetic intuition as an ideal, he who has so
often been regarded as the man of abstractions.

The mathematical example \textsc{Spinoza} uses to illustrate ``synthetic intuition'' shows that
he very well acknowledges the possibility of such intuition in special domains, especially in
the mathematical domain. He would thus not disagree with \textsc{Zeuthen}, who maintains that
synthetic intuition (a knowledge of the whole) is possible in mathematics, though he adds that
this comes from mathematical knowledge being itself symbolic. ``For it in reality comprises
only a limited number of particulars, which can be gathered into a whole, which one can then,
by the use of suitable symbols, let appear as a particular, either in a further mathematical
investigation or in one of the applications of mathematics''\footnote{\textsc{Zeuthen}:
\textit{Om Anvendelse af Regning og Ræsonnement} (Vid. Selsk. Overs. 1914, p. 275).}. As I
have suggested in my paper on intuition, a knowledge of the whole can also be reached in
special domains of experience, e.g. the psychological and the historical. But there it
receives a decidedly hypothetical character.

There is a difference between the way in which intuitive knowledge is spoken of here in the
second book and later in the fifth book. The fact is that in the fifth book the life of feeling
and will works together with the life of thought. The joy of knowledge, which in the end gave
life its value for \textsc{Spinoza}, reaches its summit when it appears that it is the
innermost ground of existence that makes knowledge possible, just as knowledge itself on its
side reaches its completion in the thought of this innermost ground and its appearance in
attributes and in modes. Thereby \textsc{Spinoza}'s synthetic intuition gets a religious stamp
and falls under what I have called practical intuition. It also gets an
ar\opage{54}tistic stamp, since it presupposes the ability to see the individual in its whole
concrete peculiarity and yet as though shot through with the great laws of existence.

What in the paper named I have called concrete intuition, under which fall the images of
intuition that form spontaneously in sensation, memory and imagination, corresponds to
\textsc{Spinoza}'s first stage of knowledge (\textit{imaginatio}). It is on the basis of it
that analytic intuition can lead to the knowledge of law\footnote{An interesting example of a
transitional form between concrete and analytic intuition is offered by the Indian
mathematician \textsc{Apastamba}'s presupposition that two triangles that are equal to one and
the same third are equal. It is intuition (the perception of surfaces) that rules; but it marks
a beginning of analysis that that presupposition is separated out from the whole visual
experience. \textsc{Zeuthen}: \textit{Hvorledes Matematiken i Tiden fra Platon til Euklid
blev rationel Videnskab} (Vid. Selsk. Skr. 1917), p. 55 f. --- In the relation between the
perception of surfaces and the perception of contours demonstrated by \textsc{Edgar Rubin}
(\textit{Synsoplevede Figurer}. 1915) there is another fine example of the transition of
concrete intuition to analysis.} and thereby possibly pave the way to synthetic intuition.
For \textsc{Spinoza} the general laws were not the last and highest task of science. The last
aim was the understanding of the particular and individual by means of the knowledge of
law\footnote{See on this point \textit{Den menneskelige Tanke}, p. 222---227. ---
\textit{Totalitet som Kategori}, p. 63---66.}. ---

With epigrammatic brevity \textsc{Spinoza} has stated the main stages of human knowledge,
beginning with the gathering of things, rising to science and closing with the intuitive
knowledge that stands on the border of science, religion and art, and that also --- which,
however, \textsc{Spinoza} would hardly admit --- must in each individual case bear a
different stamp according to the thinker's individual personality and experience of life. Nor
would he have admitted that a rounded picture of the whole, such as is to mark the summit of
knowledge, must scientifically be regarded as hypothetical.

\subsection*{f. Propositions 48---49.}
\begin{center}\textit{The Concept of Will.}\end{center}

\textit{α.} \textsc{Spinoza} closes his psychology of knowledge and his theory of knowledge
by enjoining once more that one must not confuse thought, image and word. A real thought,
consistently developed, is not something resting and passive, but the fruit of the mind's
activity and itself constantly at work, unfolding its consequences further. When it is thought
that one can believe or see something and yet not will it, one understands by belief or
insight a passive relation to an image or word (``a dumb picture on a tablet''), not a
spiritual activity. The activity by which a thought unfolds itself in its consequences is one
with a willing, and conversely: willing (\textit{volitio}) is consequence of thought. It is
impossible to have the idea of a triangle and not affirm (\textit{affirmare}) that the sum of
its angles is equal to two right angles. --- It is also enjoined that one can have full
certainty only of what is fully and clearly seen and proved. Certainty is not the same as an
absence of doubt, which may be due to the chance circumstance that no occasion for unrest in
the mind has arisen. % the Danish has "knn" for kan (can)

By the word will (\textit{voluntas}) \textsc{Spinoza} understands in this connection only the
ability \opage{55}to affirm or deny. Affirmation is quite simply a consequence of an idea one
has. \textsc{Spinoza} declares expressly that by will he does not understand desire
(\textit{cupiditas}), which expresses itself in an urge towards something or aversion to
something.

What \textsc{Spinoza} has at heart in this whole line of thought is to enjoin the active
character of the life of thought, its springing from man's innermost being, which in turn is
an individual mode of the substance at work in everything. It is in the concept of
self-activity that understanding and will are one for him. His standpoint here recalls that of
\textsc{Socrates}. \textsc{Socrates} too maintained, after all, that whoever knows cannot do
otherwise than act on his knowledge; if he does not, that is for \textsc{Socrates} a proof that % the Danish misprints kan for han (gør kan det ikke)
in reality he had no knowledge. % the Danish "gør kan det ikke" read as "gør han det ikke"

When \textsc{Descartes} taught that the will can extend further than the understanding,
\textsc{Spinoza} grants this only in so far as by understanding one means the ability to think
clearly and distinctly. He grants, of course, that there are obscure thoughts; but they rest
only on incompleteness. And he grants that there is error; but it is something negative,
resting on the absence of cases that conflict with what one has provisionally settled down
with.

When \textsc{Spinoza} declares willing and thinking one, he has in mind such characters as
those in whom deliberation is able to give a clear idea of the different possibilities, so
that the decision becomes a simple consequence of the relation between these possibilities,
--- and in whom, next, the thought on which the mind has concentrated in the decision is held
fast in its consequence until it is realized. But even such characters have gone through a
development that led through obscurities and errors, and if what is meant by the word will is
not something that first comes into being at the stage of completion, there must be other
forms of will than the one that is one with clear thinking, forms that may very well be called
active, since they can have their real cause in man's own being. Such forms of will are what is
denoted by words such as urge, striving, longing, desire, appetite. They can often stand in
opposition to the willing grounded on clear reflection --- as when \textsc{Goethe} says of his
friend:

\begin{quote}
Noch ist, bei tiefer Neigung für das Wahre,\\
ihm Irrthum eine Leidenschaft.
\end{quote}

In the admission that the will can under certain circumstances and in a certain sense extend
further than the understanding, there lies the hint of a fuller concept of will than the one
\textsc{Spinoza} maintains in this place. The one-sidedness of it is that obscurity and error
are conceived only as something negative, although they are real, positive states. Absence of
doubt means spontaneous belief, trust or assent, and does not exclude clarity and full
consistency within the circle of ideas at one's disposal. That \textsc{Spinoza} did not see
the positive character of such states is bound up with his not distinguishing between two
kinds of positive judgement, --- primary ones, which come about by immediate reflection
(analysis) on given things, and secondary ones, which reject a denial or a doubt concerning
such things. \opage{56}He misjudges above all the significance and independence of immediate
intuition, especially of imagination, when he claims (49 Schol.) that to imagine a winged horse
is nothing but to predicate wingedness of the horse\footnote{See on the matters touched on
here \textit{Den menneskelige Tanke}, p. 53---57; 78 f. --- \textit{Totalitet som
Kategori}, p. 8.}.

Finally, the identity of understanding and will can in itself just as well mean the reduction
of understanding to will as conversely. What \textsc{Spinoza} cares about is the assertion of
spiritual activity. To this side of the matter I shall come back later (at III, 9).

\textit{β.} Whereas \textsc{Spinoza} here has made willing one with thinking, in the third
book (9 with Schol.) he conceives willing as the expression of a striving for
self-preservation (\textit{conatus perseverandi in suo esse}), which, when it is felt, is
called urge (\textit{appetitus}), and when it is joined with the idea of its aim is called
desire (\textit{cupiditas}). And in the definitions that form the appendix to the third book,
desire (\textit{cupiditas}) is said to be the very essence of man in so far as it is determined
from an inner state to undertake a certain action. Finally the word desire is given so wide a
meaning that it can denote all human striving (\textit{conatus}), impulse (\textit{impetus}),
urge (\textit{appetitus}) and willing (\textit{volitio}).

It is plainly two different concepts of will that are set up at the end of the second book and
at the beginning of the third, that is, in close proximity to each other.

Attempts have been made in various ways to explain this difference. \textsc{Sigwart}
supposes (in his translation of the Short Treatise p. 208) that the urge of
self-preservation shows itself in the very affirmations and denials spoken of in II, 49. He
thus takes the presentation in the third book as the more fundamental. But he admits that
\textsc{Spinoza} himself does not say this in so many words, however near it must have lain to
him. \textsc{Baensch} goes (in his translation of the Ethica p. XVI) the opposite way, finding
in the third book precisely the summit of the intellectualism that already asserts itself in
the second book: feeling and urge, which according to II, Ax. 3 were merely to be dependent on
ideas, are made one with them in the third book, and the urge of self-preservation is then to
be nothing but the self-affirmation of the idea (\textit{idea}) that constitutes man's
essence. But \textsc{Baensch} must admit that at its summit the intellectualism then turns
over into its opposite, since it is maintained (III, 9 Schol.) that we do not strive after
something, do not will it, because it is good, but conversely regard it as good because we
will it. Here urge and willing are expressly set as more fundamental than knowledge.
\textsc{Tönnies} (Studien zur Entwickelungsgeschichte des \textsc{Spinoza}.
Vierteljahrsschrift für wissenschaftliche Philosophie VI) has sought to solve the difficulty
by the assumption that there was an interval of time between the second and third books in
their present form, and that the third book was only completed after \textsc{Hobbes}'
psychology had influenced \textsc{Spinoza}. What speaks for this explanation is that
\textsc{Hobbes} conceives the will precisely as urge or desire. In the Elements of Law (1640)
the last desire that expresses itself in deliberation is called (I, 12, 2) the will. Urge,
fear and hope cannot be willed (voluntary) \opage{57}because they are the will. The same
determinations of concepts are found in De Homine (1658) Ch. 11 and in Leviathan (1651).
Their historical presuppositions we shall come to speak of at the third book. --- There is
much that speaks for \textsc{Tönnies}' hypothesis. But it is and remains striking that
\textsc{Spinoza} in the final working together of his drafts did not notice the discrepancy.
The explanation could lie in his having spontaneously identified the urge and desire that
cannot be made one with pure thinking with the willing that is led by obscure or false
ideas. In his concession to \textsc{Descartes} he has, after all, as already remarked,
hinted at a widening of the concept of will he sets up in II, 49. But to this comes, as the
discussion of the last three books will show, that a refraction between an intellectualist
and a voluntarist conception runs through his whole psychological and ethical view.

\textit{γ.} In a concluding remark \textsc{Spinoza} first rejects the assumption of the
indifference of the will (``free will''), partly by maintaining that an idea is not a resting
image but an activity of thought that cannot be reconciled with indifference, partly by
directly rejecting the possibility of absolute states of equilibrium. A man who is equally
hungry and equally thirsty and therefore cannot decide either to eat or to drink,
\textsc{Spinoza} would --- alluding to \textsc{Buridan}'s famous ass\footnote{As for this
animal, \textsc{Schopenhauer} already searched in vain for it in the writings we have of
\textsc{Buridan} (a scholastic of the first part of the fourteenth century). Later
\textsc{Siebeck} (\textit{Beiträge zur Entwickelungsgeschichte der neueren Psychologie}.
Giessen 1891) hunted for it without success. \textsc{Siebeck} supposes that \textsc{Buridan}
used the example in oral teaching to show the difference between man and animal: man will be
able to find a difference where the animal finds none. --- \textsc{Spinoza}'s line of thought
regarding the ass goes precisely in this direction. --- As \textsc{Schopenhauer} shows
(\textit{Die beiden Grundprobleme der Ethik\textsuperscript{2}}, p. 58 f) the example is
older than \textsc{Buridan}. It occurs already in \textsc{Dante} and is hinted at in
\textsc{Aristotle}. What is new in \textsc{Buridan} is that it is an ass that is caught
between two equal possibilities.} --- rather call an ass than a man.

Finally \textsc{Spinoza} shows how the conviction of the necessity of all things teaches us to
strive to reach the highest possible understanding of existence, leads us to calm and
resignation towards the swings of fate, enjoins compassion and gives guidance for the right
government of the state.

\opage{58}\section*{Third Book.}
\subsection*{a. Preface, Definitions, Axioms and Postulates.}

In the preface the conviction is expressed in forceful words that definite laws reign in the
spiritual domain, just as in the material domain. Nature is always the same, its power and
efficacy everywhere alike; it is the same laws by which everything happens and all
transformations from one form to another take place. Therefore it is senseless to blame
nature, and it is also senseless to derive human weakness and inconstancy from sinfulness, or
to laugh or weep or be indignant over human nature as it is. What matters is to come to know
the nature of human feelings and, on the basis of this knowledge, to show what the mind itself
can do to harmonize (\textit{moderari}) them, even if it does not, as \textsc{Descartes}
thought, have absolute mastery over them. But apart from this, the understanding of human
feelings is as valuable as the understanding of the phenomena of material nature, and
therefore \textsc{Spinoza} will here apply the same method as elsewhere and regard human
actions and impulses as if it were a question of lines, surfaces or bodies.

The second definition speaks of the important determination, already mentioned, of the
concepts of activity and passivity, which, as will appear, is of decisive significance for
understanding the development of the feelings. It finds application at once in the very
definition of what feeling is. Taking his starting point here from the bodily side,
\textsc{Spinoza} defines feeling as a bodily state by which the body's power of action
(\textit{agendi potentia}) is increased or diminished, and to which a certain idea
(\textit{idea}) corresponds.

With the systematic-deductive presentation of the psychology of the feelings that now begins,
one must constantly compare the more purely descriptive presentation given in the appendix to
the third book, which is often very illuminating, although it does not always agree with the
presentation in the book itself. The descriptive appendix ends, e.g., with a ``general
definition of the feelings'', by which these are conceived as passive states and made one with
``confused ideas'' (\textit{confusæ ideæ}). Here, then, no account is taken of feelings, by
Def. 3, being able to be both active and passive, according as they have their origin in
man's own nature or are due to external causes. And no account is taken of its being expressly
shown in Propositions 58---59 how a feeling can get an active character.

\opage{59}Both in Def. 3 and in the appendix feeling is made one with idea (\textit{idea}).
Here, as with the concept of will at the end of the second book, an intellectualist conception
lies at the foundation. It is only in the part of the book that begins with Proposition 4 that
a more realistic and voluntarist line of thought asserts itself.

The concept of idea (\textit{idea}) has on the whole a fateful influence in \textsc{Spinoza}.
By the soul's being defined from the first as the idea corresponding to the body (\textit{idea
corporis}), the confusion of theory of knowledge and psychology shown in the foregoing is made
possible. And on account of the active character of ideas, which \textsc{Spinoza} rightly
stresses, an absolute identity between knowledge and will is maintained (II, 48---49),
although there is only partial identity between them, since there is spiritual self-activity
other than that which appears as knowledge. As for the psychology of the feelings in
particular, it is overlooked that there is a difference between feeling conditioned by
sensation and feeling conditioned by ideas (elementary and ideal feelings), --- and that the
very feelings that are most of all conditioned by ideas are just as little absolutely identical
with them as the expressions of the will are.

\subsection*{b. Propositions 1---3.}
\begin{center}\textit{Spiritual Activity and Passivity.}\end{center}

We are passive when we have only inadequate ideas (\textit{ideæ inadæqvatæ}), that is, such as
cannot be understood from our own nature alone but presuppose action from outside. But so as not
to be misunderstood, \textsc{Spinoza} enjoins very energetically and clearly the view of the
relation between soul and body that in his opinion follows from the relation between the
different attributes maintained in the first two books: III, 2 Schol. is grounded without more
ado on II, 7. A state or activity of the soul and of the body, e.g. a resolution and the bodily
state corresponding to it, are one and the same thing and assert themselves simultaneously
(\textit{simul natura}). By virtue of this both a psychological and a physiological explanation
must be demanded of the passive or active states men can be in, and of the changes they can
undergo.

By maintaining, in consequence of his doctrine of attributes, the simultaneity of the states of
the soul and the corresponding bodily states, \textsc{Spinoza}, without being aware of it,
comes indirectly to point to the only possible way in which the problem of the relation between
soul and body could be thought to receive a really experimental decision. For if it could be
established that a state of the soul, a sensation for example, set in only after the
corresponding bodily state had arisen, --- or that (in action) a movement was only begun in the
brain after a resolution had been taken, the popular doctrine of interaction would be proved.
We should then in the one case get a series of material states that at a certain point was
broken off and continued by a shorter or longer series of states of the soul, --- and in the
other case a series of states of the soul that at a certain point was broken off and continued
by a series of material states. Now it is not \opage{60}easy to see how one should come to
make such an experiment. It is another matter that interesting experiments have established the
continuous course of material processes (in as outside the organism) and their agreement with
what is now called the law of the conservation of energy, even where processes of the soul
could be ascertained at the same time, and that a proportionality between the organic and the
psychical has further been shown. But although \textsc{Spinoza} could not know on these points
what we now know or think we know, it is nevertheless precisely, as the scholium to III, 2 shows,
the continuity of material phenomena and the proportionality between the psychical and the
organic that he appeals to here. And the general doctrine of attributes to which he appeals in
the end has itself, as I have sought to show, arisen as a consequence of the method of the ``new
science''. He appeals expressly (in III, 2 Schol.) to the method of empirical science. No limit
has yet been shown, he says here, to what can be explained within the organism by the laws of
bodily nature, and to appeal to an intervention of the soul to explain bodily processes is only
to admit one's ignorance. This reasoning is the same as he applies elsewhere (in the concluding
remark to the first book) against the teleological explanation of nature and (in Ep. 75)
against the belief in miracles. Spiritualism, theology and belief in miracles rest for him on the
same principle. In this connection he regards his doctrine of attributes (the hypothesis of
identity) as a working hypothesis, as is more and more the case with modern adherents of this
doctrine too. At bottom no other working hypothesis can be applied in treating the relation
between soul and body than the hypothesis of attributes, which builds on the possibility of
inferring from the one circle of phenomena to the other, although a causal relation (which is
also a temporal relation) is not assumed. --- Expressions such as parallelism or duplicism, so
often used of the Spinozist hypothesis, are due to a misleading analogy\footnote{On the other
hand the expressions named (and the polemic they usually give rise to) fit \textsc{Chr.
Wolff}'s teaching. Cf. \textit{Den nyere Filosofis Historie\textsuperscript{2}}, I, p. 524
(Note 76).}. If one seeks an analogy that really fits, one can use \textsc{Fechner}'s image of
the relation between the convex and the concave side of one and the same curve, or
\textsc{Brunschwicg}'s analogy (which according to him is more than an analogy) with the relation
between algebra and geometry in analytic geometry.

\subsection*{c. Propositions 4---12.}
\begin{center}\textit{The Urge of Self-Preservation.}\end{center}

\textit{α.} As an introduction to the psychology of the feelings \textsc{Spinoza} puts forward
the teaching that every being has a tendency to preserve and maintain itself. What is to be able
to take away the existence of a thing must be something other than the thing itself. The ground
of this proposition is sought in every individual being (mode) being a special form (manner) of
the being and working of substance and therefore having its share in its own way in the ability
of substance to exist through itself and be understood through itself. The tendency to
self-preservation is not a subordinate property, but is one with the innermost nature of the
individual being. \opage{61}\textsc{Spinoza} refers here to the proposition set up earlier (I,
36), that there is nothing from whose nature effects do not spring, a proposition grounded
precisely in everything that exists expressing God's being in its definite way. It is the
positive side of the relation between substance and mode that \textsc{Spinoza} brings forward
here. The very positive (\textit{qvid positivum}) by which things are determined to act springs
from God (I, 26 Dem.). For the most part \textsc{Spinoza} stresses the negative side of that
relation, on the strength of the presupposition that every limitation, every closer
determination, is a negation (\textit{omnis determinatio negatio}, Ep. 50). We know individual
beings only from experience, but since they do after all exist, substance must show itself in
them, and therefore their power or striving (\textit{potentia sive conatus}) for
self-preservation is one with their given, actual nature (7 Dem.).

Nevertheless this concept appears rather unprepared in \textsc{Spinoza}. He does not, like
\textsc{Galileo}, \textsc{Gassendi}, \textsc{Hobbes} and later \textsc{Leibniz}, make any use
of this concept in natural philosophy. Such a use could have led him beyond the Cartesian theory
of motion, whose untenability he admitted (in his last letter to \textsc{Tschirnhausen}). In the
physical doctrines in the second book he gives a purely mechanical theory of motion in the
Cartesian spirit.

It is plainly by a purely psychological way that he was led to give the concept of
self-preservation a fundamental significance in the science of mind, and here he has
\textsc{Bruno}\footnote{Of interest for the relation between \textsc{Spinoza} and
\textsc{Bruno} is that the Short Treatise (I, 5) regards the individual thing's striving for
self-preservation as the expression of a special providence. A quite similar thought is found in
\textsc{Bruno}. (\textit{Den nyere Filosofis Historie\textsuperscript{2}}, I, p. 128).} and
\textsc{Hobbes}, and through them \textsc{Telesio}, as forerunners. Through \textsc{Telesio},
as \textsc{Dilthey}\footnote{\textit{Weltanschauung und Analyse des Menschen seit Renaissance
und Reformation}. (1914). p. 433. 463.} has shown, this whole line of thought leads back to
Stoicism.

\textit{β.} As regards man, the urge of self-preservation expresses itself both in so far as he
has clear and distinct ideas and in so far as he has obscure ideas. Therefore a widening is now
made of the concept of will set up in II, 48---49. The urge of self-preservation is called will
(\textit{voluntas}) in so far as it is referred to the soul alone, urge (\textit{appetitus}) in
so far as it is referred to both body and soul, and desire (\textit{cupiditas}) in so far as we
become aware of the urge [and its object]. What we hold good rests on the urge of
self-preservation, that is, on will (in the widest sense of the word). We therefore do not will
something because it is good, but we hold something good because we will it (9 Schol. and 39
Schol.). We can will what at bottom we do not will, and turn away from (\textit{nolle}) what at
bottom we will, as when, e.g., fear of a greater evil makes us will a lesser evil or give up what
would otherwise be a good (39 Schol.).

Without going again into the relation between the concept of will that appears here and the one
set up at the end of the second book, I shall only remark that even if, as was stated in II, 49,
understanding and will were absolutely one and the same, one could with the same right say that
we know the good because we will it, \opage{62}and that we will it because we know it.
% the Danish misprints hlev (blev, was)
When A = B, it holds after all both that A is B and that B is A. If we look at the development of
the life of the will, the first proposition can hold at more primitive stages, the latter at
higher stages: we begin by holding good what we will (have a liking and an urge for), and we can
end by willing what we have learned to regard as a good. The most fundamental proposition is and
remains, however, that our willing is the expression of our urge. This proposition has its
validity right up into the highest regions of the life of the mind. We will science and art only
because we have an urge for them. Only by this consideration does it become possible to build a
bridge between \textsc{Spinoza}'s two concepts of will. But \textsc{Spinoza} did not undertake
such a consideration, since the genetic point of view does not fully come into its own in him.

The significant difference between the two concepts of will is seen from IV, 7---8 and 14, where
it is enjoined that the knowledge of good and evil is a power in us only because it is a feeling
(\textit{affectus}), not because it is true knowledge. A feeling or urge is thus in play wherever
knowledge does an inner work (and is not merely itself a work).

\textit{γ.} The difference between urge (\textit{appetitus}) and desire (\textit{cupiditas})
consists in our being aware, in desire, of what we strive after, whereas the mere urge is blind
(III, 9 Schol.). To this difference, however, \textsc{Spinoza} attaches no decisive significance.
In the appendix to the third book (Ch. 1 Expl.) he says that whether a man is aware of his urge
or not, it is the same urge, and he can therefore use the word desire (\textit{cupiditas}) to
denote all striving (\textit{conatus}), urge (\textit{appetitus}), impulse (\textit{impetus})
and willing (\textit{volitiones}).

It is a large question whether \textsc{Spinoza} is right in this. The awareness of the aim of our
urge can become of decisive significance. When the urge is no longer blind, its direction can
become more definite; means can be sought that do not offer themselves immediately; attention can
be roused to obstacles. If only a single idea, that of the aim of the urge, is given, it can call
up other ideas, even very remote ones, and the action can thus be determined by a larger circle of
experiences. It is the introduction to what we, in purely psychological terms, call a higher
development of the life of the will, that the awareness of aims and means becomes clear. With
blindness easily go unruliness and unrest, whereas clarity about aims and means gives the
limitation that makes mastery possible. This development from obscurity to clarity does not go
on without danger. The immediate urge, the spontaneous craving of the reflex movement and of
instinct, can be weakened by reflection and speculation. But precisely this danger shows that the
relation between conscious and unconscious urge is not so indifferent as the Spinozist
proposition quoted states. And \textsc{Spinoza}'s own brilliant presentation of the significance
of knowledge for the life of feeling in the following sections of the third book must also hold
for the life of the will, since, as we shall see, all feelings of pleasure and pain are for
\textsc{Spinoza} symptoms of urge\footnote{On the narrower and the wider concept of will and on
the significance of degrees of consciousness for the life of the will see my paper
\textit{Begrebet Villie} (in the third series of Mindre Arbejder). --- Whereas
\textsc{Schopenhauer} expresses himself in a similar direction to \textsc{Spinoza},
\textsc{Fichte} assumes an altogether decisive difference between unconscious and conscious urge.
--- In the most recent times \textsc{Henri Bergson} has regarded the replacement of instinct by
reflection as a kind of fall. Cf. \textit{Henri Bergson's Filosofi}, p. 16 ff; 39 f.}.

\opage{63}In his interesting work ``Psychologie des emotionalen Denkens'' (1908) p.
566---574 \textsc{Heinrich Maier} denies that there are expressions of will without ideas. He
regards will as the innermost being of the life of the soul, but spontaneous and reflex movements
he does not take to be psychical at all, and in instinct, he claims, there is always an idea of
the aim present. % the Danish misprints forestilingsløse (without ideas)
Now it is indeed so that spontaneous and reflex movements can take place unconsciously. But
expressions of consciousness can be bound up with them, partly in that they are felt while they
are carried out, partly also in that pleasure is felt when they have free outlet and pain when
they are checked, partly even in that an unrest is felt, swinging between pleasure and pain,
before they can be carried out. In all this a tendency or urge makes itself known that is
intelligible only through the importance for self-preservation of the movements it leads to. It
is this that \textsc{Spinoza} has in view in his concept of urge (\textit{appetitus}). --- As for
the second point, it is a far-reaching anthropomorphism, that is, here, an interpretation from
the standpoint of the grown and reflecting man, when \textsc{Maier} reads an idea of purpose into
every instinctive action. He adds, it is true, that this idea of purpose appears only as a member
of a complex, is not analysed out of it. But precisely this difference can be of great
significance. A psychical element cannot without more ado be said to be the same before and after
it has been analysed out of the complex in which it first appeared. The undifferentiated
approaches the unconscious as regards the individual elements. There is in any case a whole series
of transitions. And it is always risky to jump over them without more ado and to read into the
still undifferentiated states what one has found in the more differentiated ones. The difference
is qualitative.

After this psychological excursus we return to \textsc{Spinoza}.

\textit{δ.} From the urge of self-preservation there follows (11---13) an urge to hold fast or
prefer what increases the power of action in spiritual and bodily respects. With regard to the
power of action or power of existence (\textit{potentia agendi --- vis existendi}), which is one
with perfection or reality (11 with Schol. and the closing section of the third book), a
transition can be experienced now to greater, now to less activity or perfection. In the first
case joy (pleasure) (\textit{lætitia}) arises, in the second case sorrow (pain) % the Danish doubles Tilfælde at the line turn
(\textit{tristitia}). In the appendix to the third book (Ch. 2 and 3) \textsc{Spinoza} stresses
very strongly that joy and sorrow consist in the very transitions (\textit{actus transeundi}) from
less to greater, or from greater to less perfection, --- not in the very perfection or
imperfection to which the transition is made. The decisive thing is thus the relation to the
preceding state. One must, after all, possess a certain perfection (power of action or of
existence) to be able to pass over to less perfection, and conversely a certain imperfection must
be present for a passage to greater perfection to be possible. What \textsc{Spinoza} enjoins here
is what in my Psychology I call ``the law of relativity in the domain of the feelings''. In the
closing section of the third book it is enjoined again that it is the transition itself that
\opage{64}matters. But here it is not the difference between transition and state that he wants to
impress. He maintains here, on the contrary (something in which he puts several modern
psychologists to shame), that no express comparison need take place between the preceding and the
following state; it is an altogether immediate relation between two states that conditions joy
and sorrow. \textsc{Spinoza}'s fineness as a psychological observer perhaps appears nowhere more
than in these remarks.

Beside joy and sorrow, desire is set as the third fundamental feeling (\textit{affectus
primarius}). From these three all other feelings can be derived (11 Schol.). Later on (57 Dem.)
joy and sorrow are said to be one with urge or desire in so far as this is furthered or checked by
external causes. And this is consistent within \textsc{Spinoza}'s whole philosophy, since man,
like every individual being, stands as a special form of substance, represents as it were a part
of its power: it then depends on what the relation is between the working determined from within
and the external circumstances, that is, between the individual mode and the other modes, which
like it are special forms of the power of substance.

When this relation is such that self-assertion is favoured by it, there is, according to
\textsc{Spinoza}, a tendency to dwell in the imagination on everything that can maintain and
develop it (12). And the same tendency shows itself also when a man feels joy in and love for
another being; then the ideas of what serves the preservation and development of this being will
be held fast (25). --- This is one side of the phenomenon (described by \textsc{Hume},
\textsc{Lotze} and others) that may be called the expansion of feeling. But \textsc{Spinoza}
speaks of the phenomenon only in so far as it is the expression of self-assertion. He shows (26
Schol.) how it can lead to overvaluation of oneself and of those one loves, and to undervaluation
of others. But experience shows that such expansion can also take place with feelings (such as
fear and pity) that do not further self-assertion. The relation is not so simple as
\textsc{Spinoza} thought. The individual elements of the life of the soul can, by virtue of a
psychical law of inertia, tend to spread at the expense of the life of the soul as a whole. Here
there is more possibility of disharmony than \textsc{Spinoza} assumed\footnote{In 13 Cor.
\textsc{Spinoza} says that the soul is averse to imagining what diminishes or checks its own and
the body's power. --- But owing to the tendency to expansion a strange urge can arise to dwell
precisely on what checks (to bury oneself in sorrow, for example).}.

Before we go over to \textsc{Spinoza}'s presentation of the special feelings and desires, the
remark presses forward that the unfitness of \textsc{Spinoza}'s ``geometrical method'' shows
itself nowhere so clearly as in the section that now follows. The different phenomena all come to
stand on a line with one another, although they fall into different groups, and although within
each individual group there is a difference between fundamental and derived phenomena. Yet
\textsc{Spinoza} occasionally (especially in the appendix) indicates certain points of division,
and on the strength of them the following survey can be given of the psychology of the feelings
in \textsc{Spinoza}:

\opage{65}I. Passive feelings (13---57).

\textit{α.} Influence of ideas on feeling (13---18).

\textit{β.} Feelings can be determined 1) by the relation to others (19---29) or 2) by the
relation to ourselves (30---35).

\textit{γ.} Feelings can have the character of desire (36---59).

II. Active feelings (58---59).

\subsection*{d. Propositions 13---57 (and the corresponding sections of the Appendix).}
\begin{center}\textit{Passive Feelings.}\end{center}

\textit{α.} Two special feelings \textsc{Spinoza} derives at once from the fundamental feelings
(which are, as we know, joy, sorrow and desire), namely love (\textit{amor}), which is joy
(\textit{lætitia}) joined with the idea of its cause, and hatred (\textit{odium}), which is sorrow
(\textit{tristitia}) joined with the idea of its cause. Both contain urge or desire, the former
to hold fast and preserve, the latter to remove and abolish. Already here the influence of ideas
on feeling shows itself clearly. It is by being joined with the idea of its cause that joy becomes
love, sorrow hatred. When something that would in itself awaken sorrow also has sides from which
it can give us joy, it will be able to call forth love and hatred at once; the mind will then
waver to and fro between the two feelings. This wavering of the mind (\textit{animi
fluctuatio}) corresponds in the domain of feeling to doubt (\textit{imaginationis fluctuatio}) in
the domain of ideas. And there is at bottom, says \textsc{Spinoza} (17 Schol.), only a difference
of degree between the two phenomena. --- By this last remark \textsc{Spinoza} comes to admit that
there are no ideas without effect on feeling. Besides, a certain discrepancy must be noted between
17 (the proposition and the demonstration) and the scholium belonging to it. In the proposition and
in the demonstration there is talk of the simultaneity of the two feelings, but in the scholium of
a wavering, that is, of an alternation\footnote{That the word wavering (\textit{fluctuatio}) must
be understood as an alternation is seen from a passage in the second book to which
\textsc{Spinoza} refers here. For it says there (II, 44 Schol.) that when an idea has earlier,
at different times, been joined with two different ideas, it will later call up now the one of
these, now the other, and this state is called the wavering of the imagination
(\textit{imaginationis fluctuatio}).}. It is naturally difficult for observation to distinguish
between a rapid transition and a simultaneity.

\textit{β.} 1) A whole series of feelings is determined by whether what we love or hate exists
or perishes, is furthered or checked. What is connected in an external way or by likeness with the
object of love or hatred will get a share in our feeling towards that object. By this way pity,
benevolence, indignation, envy, shame arise, --- all by the same principle, only in different
modifications and complications. \textsc{Spinoza} stresses in particular (32 Schol.) how pity,
envy and ambition spring from the same property of human nature, namely from the way in which a
feeling of joy or sorrow is altered through the influence of ideas of their causes or of things that
in one way or \opage{66}another stand in connection with these causes. --- Pity, however,
\textsc{Spinoza} explains not merely by the association of ideas --- namely by our getting a
certain feeling ourselves when we imagine a being like ourselves dominated by it; --- but also by a
reflex tendency to imitate the states and movements of others. ``He who flees because he sees
others flee, or fears because he sees others fear, or he who draws back his hand when he sees
another burn his hand, --- of him we say that he imitates the feelings of others'' (Appendix Ch.
33, Note). A starting point is indicated here for the explanation of sympathetic feelings, later
carried further by \textsc{Adam Smith}\footnote{See \textit{Den nyere Filosofis
Historie\textsuperscript{2}}, I, p. 450.}.

2) In the appendix it says, after the group of feelings just spoken of has been gone through:
``These are the feelings of joy and sorrow that are conditioned by the idea of an external thing as
immediate or mediate cause. I now go over to other feelings, which are determined by the idea of an
internal thing as cause.'' (Ch. 24 Expl.). This internal thing is ourselves. --- In the third book
itself this distinction is not made, since there love of oneself and love of another being are
spoken of together (see Proposition 25). If we follow the appendix, which is more systematic, we
now come to a series of feelings determined by the idea of ourselves and our lot. To this belong
self-contentment (\textit{acquiescentia in se ipso}), humility, repentance, pride, dejection,
shame. --- \textsc{Spinoza} thinks that humility and dejection are very rare, since human nature
has a spontaneous urge to dwell on or stress everything that can strengthen its power of action and
keep it up. Very often it turns out besides that those who seem to be in the highest degree
dejected and humble are exceedingly ambitious and envious. (Appendix Ch. 29, Expl.).

Through this whole teaching on how ideas of external or internal, immediate or mediate causes of
feelings can bring about changes of quality and direction in the life of feeling, \textsc{Spinoza}
has become the real founder of the theory of the displacement of motives, which was later carried
further by English psychology in the eighteenth and nineteenth centuries, and which is so important
for understanding the development of the soul. There is probably a historical connection between
\textsc{Spinoza} and the English psychologists, since \textsc{Hartley}, the first in their line,
shows traces of Spinozist influence\footnote{See \textit{Den nyere Filosofis
Historie\textsuperscript{2}}, I, p. 455. --- In the same work II, p. 372 I wrongly ascribed a very
clear presentation of the theory to \textsc{James Mill}; it is due to \textsc{Mackintosh}.}.

\textit{γ.} After the account of the two groups of feelings due to the relation to others and to
ourselves, it says in the appendix (Ch. 31, Expl.): ``With this I have closed the description of
the feelings that belong under joy and sorrow; I now go over to the feelings I count under desire
(\textit{cupiditas}).'' Nor is this division used in the third book itself. And besides,
\textsc{Spinoza} himself teaches (Proposition 57, Dem., see above p. 64) that joy and sorrow are
really one with the urge or desire of self-preservation, arising according as this is furthered
or checked.

The first example of a feeling of desire is longing (\textit{desiderium}), which (Appendix Ch. 32)
\opage{67}is described as urge or desire checked by the idea of something that excludes the
existence of its object. Longing is thus a desire that expresses itself under strong resistance. It
presupposes more composite and mutually conflicting connections of ideas than the simple desire
determined merely by a single idea. In the third book itself longing is called a sorrow determined
by the absence of a loved object (36 Schol.), and in the appendix too it is admitted that it forms
the opposite of the joy felt at the absence of a hated object; when it is nevertheless counted as
a feeling of desire, it is because of the predominance of the active element. --- There is here a
difference of degree, determined by the relation between the feeling and the element of will in
longing.

In Propositions 37---38 it is enjoined that a desire will stir the more strongly the greater the
decline in power, that is, the greater the sorrow that the desire replaces, and how the same holds
of joy that replaces sorrow or hatred. It also holds conversely that sorrow and hatred become
stronger when joy and love have gone before. --- \textsc{Spinoza} does not mention that there is a
limit here to heightening by contrast. For the decline can become so great that no power can be
released afterwards, and so no joy arise. In the opposite direction the limit does not announce
itself so easily; yet there may be cases where excessive joy and rapture use up all power to such a
degree that after their cessation there is only the possibility of torpor, not of suffering and
sorrow. \textsc{Spinoza} was too much of an optimist to go into this side of the matter.

Hatred is increased by mutual hatred, but can be abolished by love (Proposition 43). And even if
joy is felt that what we hate suffers or perishes, this joy will not be without an element of
sorrow when what perishes is a being like ourselves (47).

The more we think of the being we love or hate as isolated or ``free'', the greater will the love
or the hatred become, because then no other ideas and feelings can arise that could limit or check
the influence of the object on our feeling (49).

After putting forward some further propositions on individual differences and complex phenomena
within the life of feeling, \textsc{Spinoza} remarks (57 Schol.): ``All this concerns the feelings
that arise in man as a passive being (\textit{qvatenus patitur}). I have now left to add a few
things about the feelings man has as active (\textit{qvatenus agit}).'' % the Danish does not close the quote

\subsection*{e. Propositions 58---59.}
\begin{center}\textit{Active Feelings.}\end{center}

\textit{α.} By active feelings \textsc{Spinoza}, following his general definition of activity and
passivity (III, Def. 2), would have to understand all feelings whose whole cause (or at any rate
the greater part of whose conditions) lies in our own nature. But since it is a presupposition with
him that only as clearly thinking or understanding beings are we in states that have their ground
wholly in ourselves, only those feelings \opage{68}are to be called active for him that are
determined by clear thoughts. In the demonstration of 59 he expressly sets the power of
understanding (\textit{potentia intelligendi}) as one with the power of acting (\textit{potentia
agendi}), and can thereby appeal to III, 1. But in III, 9 he had maintained that the urge of
self-preservation expresses itself both in so far as we have obscure ideas and in so far as we have
clear ideas. That is, we have an urge to hold fast the stage at which we happen to stand, whether
it is rational or not. Even in the demonstration of 58 this is held fast; it is only in the
demonstration of 59 that the narrower definition of spiritual activity comes into force. ---
Nowhere does one see so clearly as here how intellectualism and voluntarism are refracted in
\textsc{Spinoza}.

In \textsc{Leibniz}' doctrine of monads, according to which the universe consists of psychical
beings at all possible stages of obscurity and clarity, and yet each unfolding itself by an inner
law, \textsc{Spinoza}'s concept of self-preservation has come more into its own. And later, in the
nineteenth century, the importance of man's living out of himself, whatever stage of the knowledge
of truth he may stand at, has been maintained from different sides. Experience had been gained of
the intricate and composite nature of the human life of the soul, and the belief of rationalism that
man's being was exhausted by ``reason'' was no longer held. Nature, originality, is wanted at any % the Danish has a stray stop (udtømt. ved)
price, because it is a necessary condition for getting further. And the possibility of
self-activity is found at every stage of development, since there can stir a striving to free
oneself from everything taken up from outside in an external and superficial way. Then, and only
then, is clarity won about how life stirs in individuals. Such a view came forward in Denmark in
\textsc{Poul Møller}'s and \textsc{Kierkegaard}'s concept of personal truth\footnote{\textit{Danske
Filosofer}, p. 123---127; 125.}. In this concept the proposition in \textsc{Spinoza} is taken in
earnest that man strives to assert himself not only when he is led by clear, but also when he is
led by obscure thoughts. \textsc{Spinoza} himself, after all, also held the conviction that if one
only worked consistently with the thoughts one has, the way would be paved for clarity and truth.
(See above p. 17).

\textit{β.} Not all feelings can get the active character. Sorrow, as the expression of the
checking and diminishing of the power of action, is excluded in advance. Only joy and desire can be
conditioned by clear knowledge. Clear knowledge is due to self-activity and causes the joy of
knowledge and thereby an urge to continued self-activity (58 Dem.), --- a love of spiritual freedom
(\textit{amor libertatis}), as \textsc{Spinoza} calls it in a later connection (V, 10 Schol.).

All working that springs from active feeling \textsc{Spinoza} counts under strength of mind
(\textit{fortitudo}), which is courage (\textit{animositas}) when, e.g. in self-control and
presence of mind, it aims at what is good for the individual himself, and is generosity
(\textit{generositas}) when, e.g. in mildness and modesty, it aims at what is good for others (59
Schol.). % the Danish misprints Sehol.

It is a consequence of the sharp opposition \textsc{Spinoza} thinks he finds between activity and
passivity, and therefore between active and passive feelings, that he denies that sorrow can be an
active feeling. Yet in sorrow too an urge stirs, \opage{69}namely to hold fast, in spite of
everything, to the value of what has been lost; or rather, sorrow itself is an expression of a
conflict between the urge to hold fast to something valuable and the experience of its loss.
Therefore man's innermost nature must be able to express itself in sorrow too. And in any case
sorrow can become an element in a more composite feeling that can have the character of activity,
even if each single element in it should be passive; the activity then shows itself in the power by
which the elements are united in spite of their different, perhaps opposite, character.
\textsc{Spinoza} too has seen this (58 Schol.). But he passes very lightly over the complex
feelings; he contents himself with remarking (59 Schol.) that exceedingly many different
compositions of the fundamental feelings are possible. He has not followed the path that
\textsc{Giordano Bruno} took in his enthusiastic work Dell' heroici furori. \textsc{Bruno} and
\textsc{Shakespeare} are the first who have been able to present complex feelings in their
peculiarity and in their value. % the Danish misprints BBUNO
\textsc{Spinoza} himself has described, in ``courage'' and ``generosity'', two complex feelings
(total feelings) of great value; but he would not admit that ``passive'' elements entered into
them\footnote{Cf. \textit{Den store Humor}, p. 39 f; 154---159.}. ---

\textsc{Spinoza}'s account of the psychology of the feelings is justly famous. He has striven to
treat this part of psychology with scientific objectivity. In the foregoing we have dwelt on the
inconsistencies that arise from the conflict of different points of view, and on the
intellectualism that asserts itself again and again in spite of significant beginnings in another
direction. But all the most important questions that the psychology of the feelings must raise come
up in him and have also been brought out in the foregoing. Thus the relation between feeling and
urge, the influence of idea on feeling, the displacement of motive, the effect of contrast as
distinct from comparison, expansion, etc.

The psychology of feeling stands in \textsc{Spinoza} as the introduction to an ethics. Man can
become unfree, overwhelmed by feelings called forth from without. The question then is how
spiritual freedom and strength become possible. Already in the third book several hints in this
direction are given. When hatred can be overcome by love (43. 47), --- when new connections of
ideas can tear the feeling away from its original object (48---49), --- and when there are even
feelings that are not of a passive nature but are expressions of strength of mind (58---59), then
there are possibilities of development in the direction of activity and freedom. At this point the
fourth book takes up the question.

\opage{70}\section*{Fourth Book.}
\subsection*{a. Preface, Definitions and Axioms.}

\textit{α.} In the preface to the fourth book the transition is made to ethics, which has given
the whole work its name. For an ethics, concepts of value of one kind or another are necessary. And
yet \textsc{Spinoza} begins by repeating his polemic, from the end of the first book, against
applying concepts of value to nature. But men and their actions too do after all belong to nature
--- according to \textsc{Spinoza}'s energetic statements at the end of the first book and in the
preface to the third book! An explanation is then required of how \textsc{Spinoza} himself can
nevertheless think himself entitled to undertake an ethical valuation of human qualities and
actions.

When the application of concepts of value in our view of nature is rejected by \textsc{Spinoza},
the reason given is that such concepts are formed by our comparing the phenomena with general
concepts or with purposes. We call things good or evil, perfect or imperfect, according as they
agree with the general types we have formed for ourselves, or with the purpose we think they are to
serve. But such comparison and its results have nothing whatever to do with things as they are
``in themselves''. The infinite being that we call God or Nature acts with the same necessity as
that by which it exists. There is no higher concept and no purpose by which it could be measured.

Nevertheless, \textsc{Spinoza} suddenly declares, we must hold fast to the concepts of perfection
and imperfection, good and evil, for we desire (\textit{cupimus}) to form an intuitive ideal of man
(\textit{exemplar vitæ humanæ, qvod intueamur}), in relation to which we can judge human actions!

If it is asked why we desire this, the answer is that our power of action is limited; it can
constantly be increased or diminished, and there is always --- as the only axiom of the fourth book
asserts --- something mightier than any individual being; --- therefore we naturally strive for the
greatest possible power of action, and with it for the greatest possible freedom (in the sense of
the ability to have the cause of one's states and actions in oneself). At its height this power or
freedom (as \textsc{Spinoza} shows in the fifth book) makes man one with substance (God or Nature),
the only thing that exists through itself and is understood through itself. \textsc{Spinoza} can so
far say that he still holds fast to the definition of perfection as one with reality given in the
second book (Def. 6). \opage{71}But I have already had occasion above (p. 33---35) to state that
this cosmological concept of perfection or reality is already colored by a determination of
value, is due to a valuation, an idealization. Only thereby does it become possible for the ethical
concept of perfection, too, to find a place within \textsc{Spinoza}'s system. According to
\textsc{Spinoza}, ethical striving becomes possible only in that a ray of substance vibrates in
every individual being, and in that this ray can reach ever higher degrees of clarity.

Properly speaking, then, \textsc{Spinoza} could very well have said that it is substance (God or
Nature) that is the highest ethical model. And thereby there arises an analogy between his
philosophy and his view of positive religion. The ideas of God of the positive religions have
significance for him only in the practical sphere. In religion God is ``a model of the true life''
(\textit{veræ vitæ exemplar}) (Tr. theol. polit. Chap. 14), while religion or theology cannot
teach us anything about God's being in itself. It is characteristic that the expression ``model''
(\textit{exemplar}) can be used in religion of God, in philosophy of the true man. The cosmological
view forms a transition between the two worlds of thought.

\textit{β.} But there still lies a difficulty in the reason \textsc{Spinoza} gives why concepts of
value cannot be applied to existence itself. The reason was to lie in their being modes of thought,
concepts we form by comparing things with one another (\textit{cogitandi modi seu notiones, qvas
formamus ex eo, qvod res ad invicem comparamus}). % the Danish misprints anvendec (applied)
(See also Ep. 21.) Now it holds of all our concepts, the purely theoretical ones too, that they are
formed by comparison. All science rests, according to \textsc{Spinoza}'s own theory of knowledge, on
comparison. The second kind of knowledge (\textit{ratio}), science proper as opposed to chance
experience, consists precisely in our not considering things in isolation but understanding their
``agreements, differences and oppositions''; only thereby does clear and distinct knowledge come
about (II, 29 Schol., see above p. 43). All things, not values alone, become objects of clear and
distinct concepts only through comparison. The very concept of truth or reality can be applied only
by means of a standard (\textit{norma veritatis}), which according to \textsc{Spinoza} is given in
the purely logical relation between ground and consequent (see above p. 6---7). Just as the
theoretical concepts have their norm in this relation, so for \textsc{Spinoza} the concepts of value
have their norm in the concept of the highest power or freedom. Considered as concepts of
comparison, then, the concepts of value stand on a line with all other concepts. The difference is
that the norm of value (the fundamental value) is more special and concrete than the norm of truth,
since it springs from the needs and the urge of human nature. In this special form the norm of
value cannot be transferred to existence. In this, according to \textsc{Spinoza}, philosophy
differs from theology. Since the theologians regard God as an ideal man, it is consistent of them to
transfer concepts of value to him and to attribute to him wishes, and pleasure and displeasure at
what happens in the world (Ep. 23 and 54).

Even from the concrete human point of view the concepts of value have, according to
\textsc{Spinoza}, only a limited validity, namely so long as man has not reached \opage{72}perfect
freedom but must still struggle to assert himself in his peculiar being. Only so long does he need
a model, and only so long has he any use for concepts of good and evil. Towards the end of the
fourth book \textsc{Spinoza} sets up the proposition that if men were born free, which presupposes
that they had only complete (adequate) ideas, they would form no concept of good and evil (68).

At the height of spiritual development we thus come ``beyond good and evil''. To that extent the
ethical, on \textsc{Spinoza}'s standpoint, can be said to be only a means and a preparation for
realizing the ideal of human nature, which for \textsc{Spinoza} consisted in the freedom that is one
with knowledge of our mind's connection with the whole of nature. In the treatise on the
improvement of the understanding, where \textsc{Spinoza} tells of the choice he made among the
various ends men can set themselves, he names moral philosophy as one of the means of reaching the
end that is for him the true one. This agrees entirely with the train of thought in the
Ethica\footnote{It is in reality also a misunderstood idealism that would set the ethical up as the
highest spiritual form of life. Cf. on this my \textit{Etik} IV, 2. 7 (4th ed. p. 88. 103).}. ---
It is not right to say with \textsc{Wundt} (Ethik, II, p. 113) that cosmology and theory of
knowledge were for \textsc{Spinoza} only aids and preparation for reaching an ethical view of the
world. For \textsc{Spinoza} the ethical is a stage of transition. On the other hand one can say that
all the thoughts he has laid down in his system signify for him the way to what was for him at once
the highest goal of thought and of the urge of life, even if on his view it would be meaningless to
regard it as the goal of all existence, since existence in itself is eternal and unchangeable. For
it is precisely in this eternity and unchangeableness that man must seek to gain a share.

There is one problem \textsc{Spinoza} does not touch on (except right at the beginning of the
treatise on the theory of knowledge), namely whether there is a norm of value that can and must be
acknowledged by all men, just as there is a norm of truth to which all must bow in the end. It is
at this point that the real sting, in ethical respects, shows itself more and more to lie.

\textit{γ.} The last difficulty as regards the position of the concept of value within
\textsc{Spinoza}'s system comes to be this: how there can be anything at all other than the eternal
and unchangeable. Why are there differences and degrees as regards ``reality or perfection'', so
that there can be talk of a transition from one state to another, as in the definition of joy and
sorrow in the third book, and so that there can arise the task of reaching from lower stages to
higher? --- In other words it is the differences of time and of quality that cause difficulties.
Why does anything at all exist other than substance, that which has its ground in itself and
therefore alone can be called perfectly free? Why does it cost so hard a labor to reach up to the
stage where neither impulses nor concepts of value are to hold?

\textsc{Spinoza}'s answer is given in the closing words of the first book: ``There has been no
lack of material for producing everything, from the highest down to the lowest degree of
perfection. The laws of nature were so comprehensive that they could contain everything that an
infinite thought could contain.'' But this answer would settle the matter only if
\opage{73}\textsc{Spinoza}'s fundamental concepts could comprise not only the constant but also the
changing, not only the completed but also the striving, not only the continuous but also the
discontinuous. It is the unchangeableness of existence in its innermost ground that here expresses
itself in its consequences and causes difficulties, both for ethics and for the theory of
knowledge. There betrays itself here a presupposition that goes far back in the history of thought,
and whose stubbornness is perhaps best seen from its still asserting itself in the founder of the
critical philosophy, in his dogmatic holding fast to the constancy of the ``Ding an
sich''\footnote{See on this \textit{Kontinuiteten i Kant's filosofiske Udviklingsgang} (Vid.
Selsk. Skr. 1893), p. 34. --- \textit{Den menneskelige Tanke}, p. 327 f.}.

It must at once be added, however, that while this difficulty comes out especially sharply in
\textsc{Spinoza}'s system, it makes itself felt in all science. Science can lead us from
differences of time and quality to general laws, which can perhaps even be formulated in the form
of mathematical equations; but it cannot derive from these laws and equations that just these or
those differences of quality and differences of time exist in the world. It is precisely a chief
point of interest in \textsc{Spinoza}'s system that it leads us to see so clearly this relation
between the fundamental concepts in which we think existence and the factually given things that
make up the content of existence. \textsc{Spinoza} was not himself conscious of this relation. His
philosophy is, as I have expressed it above, an attempt to think existence by means of the formal
categories alone.

\subsection*{b. Propositions 1---18.}
\begin{center}\textit{The Relative Strength of the Feelings.}\end{center}

\textit{α.} The strength of a feeling, its ability to maintain itself, is not determined by our own
urge of self-preservation alone, but depends on the relation of this urge to the external causes by
which we are affected. I am only a part of the whole of nature, and therefore undergo many changes
that cannot be explained by my own being alone. And still less can the whole order of nature be
explained from my being. But it is possible for a single feeling to gain the mastery over all the
other feelings, when the relation between inner and outer causes favors it (2---6). Here the
difference between true and false knowledge makes itself felt, though not directly. For, as shown
in II, 32---35, error is nothing positive, but is due to the absence of supplementing and
correcting ideas, and therefore true knowledge does not in and of itself annul what is positive in
the state to which the error belongs (1). I have the same sense image of the sun whether I think
its distance from me is great or small. When the true knowledge dawns on me, some of the ideas
attached to the sense image are changed, but not the image itself. If it is objected that the fear
of a danger falls away when it turns out that in reality there was no danger, \textsc{Spinoza}
replies that the fear also falls away when the news that there is no danger is false. It is thus
not the true as such (\textit{qvatenus verum}) that \opage{74}annuls the fear, but the strength
with which the ideas that exclude the imagined dangers come forward (1 Schol.).

\textit{β.} When \textsc{Spinoza} here speaks of the strength of ideas, in such a way that this
seems to be something different from the logical opposition (the relation of exclusion) between
them and other ideas, it must be assumed that this strength rests on the feelings attached to the
excluding ideas. For in 7 it is maintained that a feeling is not displaced except by another and
stronger feeling. The proof is conducted from the law of inertia. Neither a spiritual nor a bodily
state is annulled except by means of another and stronger state.

The important proposition \textsc{Spinoza} sets up here occurs also in \textsc{Plato}, when he
(as in the Phaedo and in the Republic) sets the urge for knowledge as a counterweight to the urge
for sensual enjoyment or for honor --- the same application of the proposition that we shall see
\textsc{Spinoza} make in the fifth book. \textsc{Spinoza} may have come to the proposition by way
of his own experience. He tells at the beginning of the treatise on the improvement of the
understanding how in knowledge he found a good that could put all other goods in the shade. He
experienced that in his hours of reflection (\textit{modo possem serio deliberare}) he was freed
from everything that could otherwise ensnare him. Possibly, however, he is influenced here by
\textsc{Bacon}, who says that just as in hunting one uses a falcon to catch other birds, so one
feeling can be used to check or regulate other feelings (De Augmentis Scientiarum. VII, 3). In
several features \textsc{Spinoza}'s ethics recalls \textsc{Bacon}'s account. Thus \textsc{Bacon}
divides ethics into the doctrine of the model (\textit{exemplar}) and the doctrine of the work of
bringing about agreement with the model (\textit{cultura animi}) (De Augm. VII, 1). And this latter
doctrine is divided again into the doctrine of the feelings (\textit{de affectibus}) and the
doctrine of the means of regulating them (\textit{de remediis}) (De Augm. VII, 3). Both in scheme
and in expression something related appears in \textsc{Spinoza}. --- After \textsc{Spinoza}'s
time, too, this proposition, important both psychologically and pedagogically, has been
maintained. Thus by \textsc{Rousseau}, who in Emile (Livre IV) expresses it thus: On n'a de prise
sur les passions que par des passions. \textsc{Kant} maintains it (in the third section of the
essay on negative magnitudes and in the introduction to the ``Tugendlehre''\footnote{Note the
distinction, fateful for \textsc{Kant}, between freedom as an abstract ability (\textit{facultas})
and as actual strength (\textit{robur}).}, without heeding the contradiction with his concept of
freedom into which he thereby falls), and \textsc{Schiller} applies it in the eighth letter on
aesthetic education. In my Religionsfilosofi (§ 116) I have shown what significance it has for the
religious problem.

The proposition needs a corrective, however. For besides through the appearance of another,
stronger feeling, a feeling, urge or desire can also disappear for lack of nourishment from the
conditions under which life is lived, just as an organ can disappear for lack of use and --- in
consequence --- of nourishment. And when a feeling is displaced by a new feeling, this is probably
in fact also because the latter lays claim to all spiritual and bodily energy.

From Proposition 7 it now follows that if the knowledge of good and evil is to become a power in
\opage{75}the mind, it must appear as a feeling (urge or desire) (14). But according to III, 9
Schol. it is an urge or desire that leads us to regard something as good or evil. The knowledge of
good and evil is thus nothing but consciousness of joy or sorrow, that is, of the way in which our
power of action is increased or diminished (8). What matters is that this knowledge can maintain
itself even though the present affects us more strongly than the absent and the future, and even
though the difference between what in our eyes is contingent, possible and necessary also exerts a
great influence on us (9---13). Where the effect of what affects us at the moment stands in sharp
opposition to the knowledge of good and evil we have so far gained, a conflict arises within us.
The Preacher is right, then, that he who increases his knowledge increases his pain. This is no
reason, however, to praise ignorance; but it contains a summons to investigate the conditions of
spiritual strength and spiritual impotence (17 Schol.).

\textit{γ.} In Proposition 18 it is maintained that the desire arising from joy is, other things
being equal, stronger than that arising from sorrow. The reason lies in sorrow (pain) being a
symptom of a decline in power of action.

In the scholium to 18 \textsc{Spinoza} himself comes to admit the unsuitability of the
``geometrical'' method of exposition, especially because it leads to great prolixity and hinders a
general view. Therefore, before he shows in detail how man can develop from being the passive prey
of commotions of the mind to possessing self-control and freedom, he will give a survey of what
reason shows to be the condition for reaching this goal.

What matters is to act according to the laws of one's own nature and thereby reach greater
perfection, a goal that must be sought for its own sake, not as a means to something else. Those
who entirely lose the urge of self-preservation and take their own lives are --- perhaps without
noticing it themselves --- overpowered by external causes (\textit{causæ latentes externæ}. 20
Schol.).

Among the means of self-preservation, cooperation with other men is of especially great
importance. When several unite their forces for common purposes, they become as it were one soul
and one body. Men who are guided by reason will therefore not desire anything for themselves that
they do not also desire for others; they will be just, faithful and honorable. --- One sees
already here the difference between \textsc{Spinoza}'s and \textsc{Hobbes}' ethics. For
\textsc{Spinoza} social ethics is a positive and direct continuation of individual ethics, while
for \textsc{Hobbes} it stands predominantly in a negative or indirect relation to it. In a letter
(Ep. 50) \textsc{Spinoza} himself has spoken of this difference.

\subsection*{c. Propositions 19---28.}
\begin{center}\textit{The Ethical Character and the Highest Good.}\end{center}

\textit{α.} The urge of self-preservation decides what is regarded as good and evil, and therefore
everyone, according to the laws of his own nature, yearns for what he regards as good and abhors
what he regards as evil. That urge is one with man's being. And the knowledge \opage{76}of good
and evil is, as we have seen, nothing but consciousness of joy and sorrow, i.e. of the increase
and diminution of the ability to assert ourselves. --- One sees here how knowledge, feeling and
will in the end become one for \textsc{Spinoza}, yet in such a way that it is a willing and acting
that is our innermost being and that becomes conscious of itself in knowledge and feeling. Here
lies the possibility of a richer and deeper psychology than the marked intellectualistic tendency
that \textsc{Spinoza} brought with him to his philosophizing could allow to develop.

But when man acts blindly or is guided by incomplete (inadequate) ideas, he does not, according to
\textsc{Spinoza}, act from himself; he is then not active, not free. There is therefore nothing of
which we can know with certainty that it is good, except knowledge. \textsc{Spinoza} speaks here,
as one sees from the autobiographical fragment in the treatise on the improvement of the
understanding, from his own experience. External things, he says in the appendix to the fourth
book, are good for us only in so far as they are conditions for a share in that life of the mind
which consists in knowledge (\textit{ut mentis vita fruamur, qvæ intelligentia definitur}. IV App.
Chap. 5).

\textit{β.} \textsc{Spinoza} has, as mentioned earlier, not noticed any difficulty in the relation
between the two starting points that condition his psychological and ethical train of thought,
--- the one according to which the urge of self-preservation, --- the other according to which the
urge for knowledge constitutes man's being. % the Danish misprints hyilket
The opposition is there, but it need not have been insuperable. The urge of self-preservation can
lead to regarding knowledge as a good, first perhaps as a derived good, a means, later, perhaps
after a whole psychological process, as an end, a good in itself, so that it can be said (29
Dem.) that it is not for the sake of any other end that we strive after knowledge. When
\textsc{Spinoza} grounds this proposition on our not striving to preserve ourselves for the sake of
any end at all (25), the presupposition underlying it is that knowledge and self-preservation are
one. And yet experience shows, as \textsc{Spinoza}'s own experience had also shown him, that they
do not simply coincide. Through the third book's demonstration that displacement of motive is
possible, \textsc{Spinoza} had provided himself with a basis for understanding how what first has
value as a means can later become an end in and for itself. It is strange that he does not make
use of this theory here in the fourth book. His sound eye as a psychologist fails him here, and the
reason must be sought in his being unable, given the fundamental concepts of his philosophy, to
attribute any decisive importance to becoming, transition and development. He regards the urge of
self-preservation and the urge for knowledge as identical, instead of holding to this, that the
latter can develop out of the former according to natural laws and perhaps, at its height, gather
all strength of mind into itself. % the Danish misprints Aaandskraft
--- In the fifth book this peculiarity of his train of thought will come out even more sharply.

If we look back at the relation between knowledge and will in \textsc{Spinoza}, we find three
points of view that conflict in him: 1) man's being is thought (III, 1---3, cf. II, 40); 2) thought
and urge are identical (IV, 25---26, cf. II, 49); 3) thought is dependent on urge (III, 9 Schol.).
A genetic psychology would try to \opage{77}show how a development can take place from 3) through
2) to 1). But the dogmatic and the genetic points of view conflict in \textsc{Spinoza} in a way of
which he himself was not conscious.

\textit{γ.} True knowledge, which for \textsc{Spinoza} stands as the highest good, would be a
one-sided determination in an ethics meant to hold not merely for men of science (or rather for
men of science in so far as they are merely men of science). The life of the mind (\textit{mentis
vita}) of which \textsc{Spinoza} speaks can appear to us in more forms than the intellectual. The
life of mood and of imagination are forms of the life of the mind that, just as well as the pure
life of thought, can satisfy the spiritual urge, and can do so without breaking the norms of the
life of thought. If we want to adapt \textsc{Spinoza}'s ethics to our own use, we must therefore
put spiritual culture in general in the place of knowledge or understanding. New problems do indeed
arise thereby, since conflicts can arise between the different spheres of culture; but this is
precisely a testimony to the wealth the human life of the mind can possess\footnote{Cf. my
\textit{Etik}, Chap. 27---35 and my essay \textit{Aandelig Kultur} (Mindre Skrifter, Tredie
Række).}.

\subsection*{d. Propositions 29---37 (and the Appendix Chap. 9---29).}
\begin{center}\textit{The Significance of Social Life.}\end{center}

\textit{α.} So long as men's feelings are essentially determined from without, it is a matter of
chance whether they agree, and then it is also a matter of chance whether the individual man agrees
with himself. There is soon a possibility of opposition and strife, for what different men strive
for may be something that only one can possess. When Peter and Paul love one and the same thing,
they agree, to be sure, about its value, but in reality they stand opposed to each other; for if
Peter possesses the thing he will feel joy, but Paul will feel sorrow, and vice
versa\footnote{Already \textsc{Aristotle} remarks that for all to will one thing is not enough for
concord, when each wills the possession for himself alone (\textit{Eth. Nic.} IX, 6). ---
\textsc{Kant} has a similar reflection (\textit{Kritik der praktischen Vernunft} § 4 Anm.) and
cites the words of Francis I: ``What my brother Charles wants, I want too!'' --- For they both
wanted Milan.} (34 Schol.). If there is to be concord and a real life together, the feeling must
be attached to something that can be a common good, so that the striving of the one does not
hinder that of the other, but they can help and support each other. And this will be the case the
more men are guided by active feelings, that is, by such as are conditioned by knowledge and
reason. Reason is a spiritual good, in the acquiring of which the one does not stand in the other's
way. Here it is no accident when men agree. Without the ability to rejoice in spiritual goods man
would not be man. And a good such as knowledge the individual will not keep for himself, but will
strive to communicate to all. As regards material goods too, it will be of importance that forces
are united; one can then better secure oneself against danger and obtain what is required to
sustain life.

In the sphere of material as well as of spiritual culture, then, common \opage{78}striving is
possible when man is guided by reason, --- by what constitutes man's real being.

\textit{β.} \textsc{Spinoza} is not so optimistic, however, as to overlook the great hindrances to
such a community. Art and vigilance are required to lead men in this direction. An education for
freedom is required. And there is no way in which a man can better show his skill and his talent
than by influencing men so that they come to live under the authority of thought itself
(Appendix, Chap. 9).

Men are inconstant. Few are those who let themselves be guided by reason. Most are envious and
vengeful. But one nevertheless goes the wrong way when one scolds men for their faults instead of
guiding them to the acquisition of valuable traits of character, and when one breaks their courage
instead of strengthening it. One should bear men's insults with a calm mind and work for everything
that can further concord and friendship, and that is one with working for what justice, equity and
honor require.

The state is a support for this striving to subordinate what is changing and conflicting to common
fixed rules. Its working rests on two psychological facts. The first is that a feeling is overcome
only by another feeling, the second that fear of suffering an evil oneself can restrain one from
inflicting an evil on others. The state relies on these facts with its threats of punishment,
which can accomplish what reason cannot. Only in the state, therefore, not in the pure state of
nature, can there be talk of guilt and merit, of wrong and right (37 Schol. 2). Beyond what can be
compelled by threats the state has no authority, and therefore no right, and so there remains a
sphere in which man, while living in the state, keeps the power and the right to follow his urge
of self-preservation. It is of great importance that the state know its limits. For one cannot
rely on a concord that is brought about only by instilling fear; it is not by arms but by love and
magnanimity that souls can be conquered. (Appendix, Chap. 16).

\textsc{Spinoza} gives a closer development of the relation between the individual and the state
in the theological-political treatise, Chap. 16.

\subsection*{e. Propositions 38---58.}
\begin{center}\textit{A Scale of Value for Joys and Sorrows.}\end{center}

\textit{α.} As \textsc{Spinoza} passes on to a valuation of the different feelings according to the
standard he has gained in the foregoing, he begins with a biological and sociological
consideration. Good is everything that makes the body fit to act in as many-sided a way as
possible, and that preserves the constant relation between rest and motion within the body. Good,
further, is everything that leads to concord in society and thereby makes it possible for men to
live according to true knowledge.

Then the various feelings are gone through. The main proposition is that joy \opage{79}is always
good in and of itself, sorrow (pain) in and of itself evil\footnote{On the justification of this
proposition see my \textit{Etik} VII, 1, where some historical notes on it are also given.}. For,
it is said in the appendix (Chap. 31), the greater the joy we feel, the greater the perfection to
which we pass, and the more we partake of the divine nature. But it is joy as a total feeling
(\textit{hilaritas}) that is meant, --- joy as embracing man's whole being. In this sense, now, joy
is easier to conceive than to find in experience. For mostly there is a single special feeling that
expresses itself in excessive degree and makes the mind fasten so on a single object that nothing
else exists for it at all (44 Schol.). As for the joy that embraces only single parts of man's
being, that is, partial joy, ``titillation'' (\textit{titillatio}) (``pleasure'' in the narrowest
sense of the word), it can be carried so far that it is at the expense of the perfection of the
whole individual. --- As a partial feeling, sorrow or pain can be good. For it indicates that there
is still life, --- that the injured part is not rotten. Although repentance and shame, for example,
are in themselves sorrow, he who is in such sorrow is yet more perfect than the shameless man who
no longer has any urge for honor. (58 Schol.). A feeling of which sorrow is an element cannot be
called good in and of itself. This holds of fear and hope, both of which presuppose pain (expected
to come or expected to be removed), and it holds of feelings that presuppose a preceding hope and
fear, such as the feeling of security and the joy over sorrow overcome (\textit{securitas et
gaudium}) (47 Schol.).

By total feeling \textsc{Spinoza}, as one sees, means feelings that do concern man's whole being but
are simple, homogeneous and uncompounded, --- not such feelings as those in which elements and
aftereffects of several mutually different, perhaps conflicting, feelings make themselves felt. To
feelings of this latter kind he would in any case attribute a lesser value\footnote{On the
significance of such total feeling see \textit{Den store Humor} Chap. 1. Cf. above p. 55 and 68
f.}.

\textit{β.} Hatred of men can never be good. The same holds of envy, mockery, contempt, revenge.
Hatred is not rooted out by hatred, but by magnanimity. Those who are conquered in this way yield
with joy, not from impotence, but precisely because their spiritual power is increased. Feelings
such as repentance, pity and humility are not good in themselves, but, men being as they are, the
feelings named can be of use and ease the transition to wisdom. Despondency, which is inclined to
throw itself away (\textit{abjectio}), is not good in itself either, but is yet easier to cure than
self-overestimation and pride (\textit{superbia}), to which, for that matter, it may very well be
related (57 Schol.), since it is above all sorrow at one's own impotence, by comparison with the
power and fortune of others, which in turn means that it easily passes into envy. --- These last
remarks recall \textsc{Nietzsche}'s characterization of the unfortunate and their ``slave
revolt''\footnote{Cf. \textit{Moderne Filosofer} p. 140.}. But there is a difference in principle
in the standard: \textsc{Nietzsche}'s standard is power in and of itself, \textsc{Spinoza}'s
\opage{80}is taken from the spiritual goods that can become accessible to all, and from which pride
and humility can each in its own way exclude us.

\textit{γ.} At the close of this scale of value of joys and sorrows \textsc{Spinoza} reminds us that
the whole valuation he has undertaken here holds only from the point of view that human striving
provides. The laws of nature, he adds, concern the general order of existence, of which man is only
a part. It has not been his task to report on men's faults and unreasonable conduct, but to show
the nature and properties of things. It is still his intention to consider human feelings and
their properties in the same way as other natural phenomena. And even if human feelings do not
exactly bear witness to men's cleverness and wisdom, they still bear witness to the power and
wisdom of nature, no less than other phenomena we delight in investigating (57 Schol.). For
\textsc{Spinoza} there is no discord between cosmology and ethics.

\subsection*{f. Propositions 59---66.}
\begin{center}\textit{A Scale of Value for Desires.}\end{center}

\textit{α.} Although \textsc{Spinoza} has seen clearly that there is something artificial in
distinguishing between feeling and desire, since desire is determined by the idea of the cause of
the feeling (III, 12---15; 57 Dem.), he nevertheless says (IV, 58 Schol.) that now, having spoken of
the value of joys and sorrows, he will go on to give a valuation of the desires.

The actions that spring from passive feelings can also be performed from true knowledge. For to act
according to true knowledge (\textit{ratio}) is nothing but to act according to the innermost being
of our nature. When sorrow, and sometimes joy too, is evil, it is because they are passive
feelings, expressions of a checked power of action. If the check falls away, so that self-assertion
and transition to greater power of action become possible, the actions that were previously due to
passive feeling are performed in a far more perfect way. As an example the difference between
blind revenge and deliberate punishment is cited (59).

It is of decisive importance whether the desire springs from only a single part of our being, or
whether it takes account of all sides of our nature. Further --- whether it springs from a cause
that acts only in the present moment, or whether it takes account of the future. From the
standpoint of true knowledge the whole difference between past, present and future falls away. The
difference between the stages of knowledge rests essentially on this, that differences and
relations of time play a part at the lowest stages but fall away at the higher stages, where one
sees things ``under the aspect of eternity''. (Cf. above p. 47---54). The knowledge of good and evil
does not belong to the highest stage of knowledge, precisely because it presupposes the influence
of differences of time. For we do not look equally clearly on past, present and future with regard
to what is to be preferred. Therefore the knowledge of good and evil that we have, even if it is
correct, is only abstract and general, and cannot gain in the particular case the influence that in
and of itself it lies in its nature to exert. Knowledge of evil, now, is always an incomplete
(inadequate) knowledge, since, \opage{81}being conditioned by a transition to lesser power of
action, it is not explained by our own being, but is understood only through our dependence on what
is outside us. If we had only perfect knowledge, we would form no concept of evil. He who is guided
by his own being, i.e. by clear thoughts independently formed, always does what he wills, namely
what is for him the most excellent thing in life. He can in truth be called a free being (60---66).

\textit{β.} If we now look back on the whole development of imperfect and perfect action that
\textsc{Spinoza} gives in the fourth book, it turns out that what is irrational in ethical respects
is due, on his view, partly to differences in the degree of strength of the feelings (IV, 1---18),
partly to differences between inner and outer causes (19---28), partly to differences with respect
to time (59---66). In the end we come back to the fact that there are differences in existence at
all, which brings it about that our thinking comes to consist in comparison (see the preface to the
fourth book), not in beholding the inner bond between things. So long as we do not see everything
as one, we act more or less blindly. The ideal (\textit{exemplar}) is the free man, to whose
characterization \textsc{Spinoza} now passes.

Although in \textsc{Spinoza}, then, a sharp opposition makes itself felt between man as experience
shows him to us, tossed about in blindness and in illusions, and the ideal man, the model for our
thinking and acting, he nevertheless feels a deep aversion to those whom he calls now the
superstitious, now the melancholics, now the satirists, who are better at castigating the vices
than at showing the way to true freedom, --- who teach men rather to shun evil than to love good,
--- and who work by instilling fear instead of leading men forward through true knowledge. (63
Schol. and Appendix Chap. 13 and 31, cf. already 35 Schol. and 45 Schol.).

\subsection*{g. Propositions 67---73.}
\begin{center}\textit{The Free Man.}\end{center}

\textit{α.} The free man denotes the highest unfolding of human nature, a life according to its own
laws. As is clearly expressed in the Tractatus politicus (II, 7. 11), freedom is a virtue or a
perfection; therefore man can be called free only when he follows true knowledge, exists and acts
from the inner laws of human nature, from causes that are understood from man himself. Freedom in
this sense does not annul necessity but presupposes it.

The free man thinks not of death but of life. He yearns for what is good for its own sake, not from
fear of death. His wisdom is a meditation on life (\textit{meditatio vitæ}). The difference between
good and evil does not exist for him; this is the right thought in the old story of the prohibition
against eating of the tree of the knowledge of good and evil. His strength of mind shows itself as
much in averting dangers as in overcoming them. Since he acts from true knowledge, he does not
practice deceit. If true knowledge could permit deceit as a general rule, the whole of society would
become impossible. The free man lives in the conviction \opage{82}that everything follows from the
necessity of the divine nature, and therefore, when something appears to him as evil, horrible or
unjust, he will assume that the reason is obscure and fragmentary knowledge, and he will strive to
understand things as they are in themselves, without letting himself be influenced by hatred and
scorn, by anger and arrogance.

This characterization is, as \textsc{Spinoza} himself remarks, a further working out of what was
already described in the third book (III, 59 Schol.) as strength of mind (\textit{fortitudo}) in its
two main forms, as courage (\textit{animositas}) and as generosity (\textit{generositas}). That man
rises above passive feelings, such as pity and sorrow, therefore by no means leads to
self-absorption or cold indifference. At the height of spiritual development the opposition between
one's own welfare and that of others no longer makes itself felt, because there is the strength and
the mind to work for both, and --- we may add in \textsc{Spinoza}'s spirit --- because in the full
development of one's own personality the best contribution is made to the true interest of human
society. The state's own end is freedom. (Tract. theol. polit. Chap. 20). % the Danish reads polt.

\textit{β.} In \textsc{Spinoza}'s ``free man'' there appears the most thoroughly clarified form of
the thought that, as already mentioned above (p. 61), carries the whole ethics of the Renaissance.
This thought developed under the influence of Stoicism. But there is a significant difference
between ancient Stoicism\footnote{On Stoicism cf. my essay \textit{Hedenske Sandhedssøgere}
(Mindre Arbejder. Første Række) and the essays on \textit{Epiktet og Marcus Aurelius} (Mindre
Arbejder. Tredie Række).} and its successors in the time of the Renaissance. While self-assertion
was with the Stoics a refuge, an expression of resignation, in the thinkers of the Renaissance
(from \textsc{Telesio} through \textsc{Bruno}, \textsc{Descartes} and \textsc{Hobbes} to
\textsc{Spinoza}) it acquired an active and positive character; it sprang from trust in life and
enthusiasm for its great laws, and it was nourished by the widening of the understanding of nature
and of the view of the world that the new research brought.

The proposition (68) that if men were born free they would form no concepts of good and evil is of
great importance. For it follows from it that the burden of proof rests on him who considers it
necessary to pass moral judgment on human actions. The ethical problem lies not merely in the
question whether and how ethical judgments can be grounded, but first and foremost in the question
whether ethical judgments are necessary at all. That a man passes judgments of good and evil is for
\textsc{Spinoza} a sign that he does not live his life out of inner strength (with courage), and
that he is not able either (with generosity) to help others to live their lives out of themselves.
As it is put in a statement already quoted above: ``There is no way in which a man can better show
his skill and his talent than by influencing others so that they can live under the authority of
thought itself''. (Appendix. Chap. 9).

This thought contains the most intimate connection between individual and social ethics that is
possible at all\footnote{Cf. on this my \textit{Etik}, III, 10 and XIII, 6.}. By giving in to the
inclination to moralize one easily comes to put obstacles in the way of what should be precisely
\opage{83}the goal of all morality: the development of independent personalities. That our actions
are judged --- by ourselves or by others --- is ethically justified only when it is a link in an
education for freedom. In \textsc{Rousseau}'s pedagogical gospel it was therefore also a main
proposition that education must first and foremost proceed indirectly, or, as he puts it, be
``negative''. It is the principle of personality itself that requires this.

In his grounding of the claim that the free man does not act deceitfully (72), \textsc{Spinoza} uses
a line of thought that later plays a great part in \textsc{Kant}, in fixing the general moral
principle. \textsc{Kant} maintains that it could not be set up as a general principle that one may
appropriate goods entrusted to one, or give a promise with the intention of not keeping
it\footnote{The line of thought is old and occurs already in \textsc{Plato} (Republic I, p. 362
C).}. \textsc{Kant} did not see that underlying this line of thought is the presupposition that
society is to continue, and that it can continue only through mutual trust\footnote{Cf. \textit{Den
nyere Filosofi's Historie\textsuperscript{2}}, p. 72---84.}. \textsc{Spinoza}, on the other hand,
starts from society as a historical fact whose conditions are to be found. The reason by which the
free man is guided in \textsc{Spinoza} includes also the conditions of the continuance of society,
and declares it absurd that society could continue without mutual trust (72 Schol.). In spite of
all \textsc{Spinoza}'s speculation, it is always to be felt in his ethics, as in the other parts of
his philosophy, that experience is for him the first stage of knowledge.

\opage{84}\section*{Fifth Book.}
\subsection*{a. Preface and Axioms.}

\textit{α.} When \textsc{Spinoza} begins the preface to the fifth book by saying that he now
``passes on to the second part of the ethics'', he no doubt regards the fourth book as the first
part, since ethics in the narrower sense begins with it. But when he says that this second part is
to treat of ``the way that leads to freedom'', its content does not stand in the same relation to
the fourth book as freedom to bondage. Already in the fourth book he has shown the possibility of a
development from dependence on feelings and desires aroused from without to spiritual freedom, and
the book reached its height in the portrayal of the free man. The fourth book oddly enough bears
the title ``Of Human Bondage'', but this title does not fit its content. Indeed, already at the end
of the third book (Propositions 58---59) the possibility is shown of a transition from passive to
active feelings, that is, from bondage to freedom.

The fourth book ought to have ended with Proposition 66, and Propositions 67---73 (the portrayal of
the free man) ought to have formed one book together with the first twenty propositions of the
fifth book, which deal with the remedies against dependence on the feelings (\textit{de remediis
affectuum}), that is, with the work of becoming a free man. The remaining part of the fifth book
develops how he who has reached true knowledge is independent of time and stands as an eternal form
of God or substance. Already in the fourth book (74 Schol. and Appendix Chap. 4---5) he had come
to speak of intuitive knowledge as the highest goal and as what determines the value of everything
else.

\textsc{Joachim}\footnote{\textit{The Ethics of Spinoza}, p. 268.} determines the relation between
the fourth and the fifth book thus: the fourth book describes man's knowledge and life at a stage
where perfect ideas do guide his conduct but do not fill his whole soul. % the Danish marks the note 2) and sets it as 1)
This view, however, decidedly conflicts with the portrayal of the free man at the end of the fourth
book. What is added in the fifth book is the consequence that the free man is raised above time,
not only thinks under the aspect of eternity but is himself eternal. At the end of 20 Schol.
\textsc{Spinoza} says that, having concluded the consideration of everything that concerns the
present life, he will pass on to what concerns the continuance of the soul apart from the body
(\textit{sine relatione ad corpus}). Apart from the first twenty propositions, it is a
religious-mystical rather than an ethical mode of consideration that is characteristic of the fifth
book.

\opage{85}\textit{β.} The preface is essentially a critique of \textsc{Descartes}' psychology, called
forth by \textsc{Descartes}' ascribing to man without more ado, like the Stoics, the ability to
master the feelings. As \textsc{Brochard} remarks in his comparison between \textsc{Descartes}' and
\textsc{Spinoza}'s psychology\footnote{\textit{Le traité des passions de Descartes et l'Éthique de
Spinoza} (in his book Études de philosophie ancienne et de philosophie moderne, p. 330).},
\textsc{Spinoza} dwells only on the points on which he disagrees with \textsc{Descartes}, and not on
the agreements; thus he does not mention the weight that \textsc{Descartes} (like the Stoics before
him) laid on understanding the nature of the feelings, and in general on the influence of ideas on
the life of feeling.

\textsc{Spinoza}'s critique concerns four points:

1. \textsc{Descartes} gives no explanation of how soul and body (thought and a fragment of
extension) can be united. He assumes such an opposition between them that he must appeal to God,
the cause of the whole universe, to be able to solve the riddle. But that is no scientific solution,
no explanation by a special cause (\textit{per proximam causam})\footnote{\textit{Proxima causa} in
\textsc{Spinoza} corresponds to \textit{vera causa} in \textsc{Kepler}.}.

2. It becomes unintelligible how it can be decided whether the soul or the ``animal spirits'' are
the stronger. Which of them is able at the decisive moment to give the pineal gland the decisive
push? How can the two forces that are to meet here be compared? Willing and motion are not
commensurable. % the Danish has a stray stop: eller. (or)

3. Anatomy shows that the pineal gland does not have its place in the middle of the brain, as
\textsc{Descartes} thought. It cannot therefore be moved as easily and in as many ways as he
supposed. Moreover, not all nerves lead out to the cavities of the brain\footnote{These and similar
anatomical objections had been made by \textsc{Thomas Bartholin} (1651) and \textsc{Steno} (1665).
See on this \textit{Den nyere Filosofis Historie\textsuperscript{2}}, I, p. 517. ---
\textsc{Steno}'s critique \textsc{Spinoza} may have known from his personal acquaintance with him in
the years 1661---1663. \textsc{Steno}'s treatise on the anatomy of the brain has been translated
into Danish by Dr. \textsc{Maar}.}.

4. A purely psychological explanation of the mastery over the feelings must be required. Since all
spiritual power is conditioned by understanding alone, understanding must be made the basis where
the talk is of remedies against bondage to the feelings. But instead of giving a purely
psychological explanation of how understanding can develop and gain power, \textsc{Descartes}
describes a purely imaginary struggle between soul and brain. The right procedure would have been
to describe the phenomena of the soul in their mutual connection and then to consult anatomy and
physiology. On \textsc{Spinoza}'s view every struggle between motives in the soul has its
simultaneous expression in the relation between different brain processes. There can then be no
question of the soul's being able to take a decision while at the same time the brain could
initiate a movement in a quite different direction. % the Danish ends the paragraph with a comma

As regards the second objection, it rests on the reasons and points of view that led
\textsc{Spinoza} to his doctrine of the attributes. (See above p. 38---45). But he overlooks that he
himself assumes a proportionality between psychical and bodily phenomena, so that he comes to face
the difficulty of showing a common measure for the two groups of phenomena. There is, to be sure,
the difference that in \textsc{Descartes} \opage{86}the one kind of phenomena was at a certain point
to form the continuation of the other (the movement of the pineal gland to succeed the decision, or
sensation and feeling to succeed the movement of the pineal gland) and to form a series within
which continuity was broken at a single point, while in \textsc{Spinoza} two series, each
continuous in itself, run proportionally to each other, so that man's being shows itself precisely
in this proportionality. But the difficulty then becomes that of finding a member of the one series
that corresponds to certain members of the other series, a difficulty that becomes the greater the
more weight is laid on the continuity within each series\footnote{I have touched on this point in
\textit{Filosofiske Problemer} (Universitetsprogram 1902), p. 23---25.}. To remedy this difficulty a
concept of psychical work must be set up that can correspond to the concept of material
work\footnote{Cf. on this \textit{Den menneskelige Tanke}, p. 4---6.}, and \textsc{Spinoza} did not
reach this, though he so often speaks of power of action in the spiritual as well as in the
material sphere. It is his Cartesian philosophy of nature that hinders him in this.
\textsc{Leibniz} was the first to begin to see more clearly here\footnote{See \textit{Den nyere
Filosofis Historie\textsuperscript{2}}, I, p. 357---363.}.

\textit{γ.} No new definitions are set up in the fifth book, only two axioms. --- The first lays
down that when two different activities (\textit{actiones}) arise in the same person
(\textit{subjectum}), changes must take place in both of them, or at least in one, until the
opposition is removed. It is evidently a proposition of experience that is to be set up here. But
it is clear that this proposition rests on the presupposition that it is not indifferent to the two
activities that they take place in the same subject. What is presupposed is a fundamental tendency
to work together the different elements and processes that may occur within the same whole. By
bringing out this presupposition \textsc{Spinoza} would have arrived at a concept of work that could
have been fruitful for his philosophy. The axiom is used in Proposition 7, where there is talk of
the opposition between such feelings as are attached to thoughts of laws, that is, of ``what holds
alike in the part and in the whole'', and such as are attached to a single, isolated thing. The
conclusion is drawn from Ax. 1 that feelings of the latter kind must in the long run yield to
feelings of the former kind. What \textsc{Spinoza} concludes here is no doubt also correct,
although it can take a very long time before an impatient feeling, concentrated on a single thing,
yields to feelings determined by relations that assert themselves again and again. The life of
feeling, after all, develops more slowly than the life of ideas, single outbursts excepted. But even
if \textsc{Spinoza} is right in the special application of the axiom, he is hardly right in setting
up the content of the axiom as the only possibility. It may be that there is a swinging between the
two tendencies, or that a state of tension continues to exist between them, or that the one gains
sole mastery, or that a total feeling arises whose timbre differs according to the mutual relation
of the opposites\footnote{Cf. \textit{Den store Humor}, p. 39, 73, 88 ff.}. It is dogmatic to assume
so simple a mechanical leveling as \textsc{Spinoza} (and some modern \opage{87}psychologists with
him) imagines. More psychical oppositions can persist side by side than one usually
imagines\footnote{Cf. \textsc{August Wimmer}: \textit{Psykogene Sindssygdomme} (Jubilæumsskrift for
St. Hans Hospital, p. 107), where this is grounded by way of pathological observation.}.

The second axiom --- that the power of the effect is determined by the power of the cause ---
brings nothing new. It is accordingly grounded by a reference to III, 7, which in turn is grounded by
I, 29. 36.

\subsection*{b. Propositions 1---20.}
\begin{center}\textit{Remedies against the Feelings.}\end{center}

According to \textsc{Spinoza} it is not really against the feelings themselves that remedies are
needed. What matters is to annul the passive character some feelings have, by converting them into
active form. Through a purely psychological development \textsc{Spinoza} shows the possibility of
this, which had, for that matter, already been hinted at in the third and fourth books.

According to \textsc{Spinoza}'s doctrine of the attributes and the identity hypothesis formulated in
it, this possibility can and must be investigated both from the psychical and from the physical
side, and it is indifferent from which side we begin. In the second, third and fourth books
\textsc{Spinoza} began his proofs from the physical side, though he constantly (see e.g. III, 12
Dem.) holds fast to the two-sided view. In the first part of the fifth book (1---20) he begins with
the psychical side, evidently because, as regards the question treated here, this side is the
better known. Later (from Proposition 21 on) he starts from the physical side, since here he thinks
he has clear and simple premises to start from. It is not self-contradictory, as has sometimes been
thought, but fully consistent, when \textsc{Spinoza} thus sees the matter now from the one side, now
from the other; and every scientific treatment of the questions at issue here in reality proceeds
in this way, whether one is conscious of it or not.

The ways \textsc{Spinoza} points out to mastery over, or metamorphosis of, the passive feelings are
three.

\textit{α.} What matters is to bring the feelings into new connections of ideas, in that the
thought of their external cause is replaced by other thoughts (2). Of special importance here are
the thoughts of the common properties of things and of their general laws, --- of what holds
equally in the particular and in the whole and is always present (7). By being attached to such
thoughts the feeling no longer has a single and vanishing ground, but is conditioned by true
knowledge (\textit{ratio}) and can be understood as a link in the great connection of things.
Obscure thoughts are replaced by clear thoughts, and the feeling now acquires another character.
Thus, for example, the lust for power or the pride that would force others to live as we see fit
passes over into the generosity that strives to develop others to live according to clear
understanding. It is one and the same urge that underlies the passive and the active states, but a
metamorphosis has taken place in the transition to active form by means of clear knowledge
\opage{88}of the circumstances (4 Schol.). Sorrow ceases to be sorrow when its cause is clearly
known (15 Dem.; 18 Schol.): for it is replaced by the joy of knowledge. A feeling transformed in
this way will mean a greater power in the mind than a feeling attached to a single object. A firm
connection of thoughts also gives connection in the life of feeling. Greater power is required to
check feelings that are ordered and connected according to clear knowledge than to check
indefinite and changing feelings (10 Schol.). The thought of the eternal and unchangeable can in
the end limit the dominion of the passive feelings in the highest degree (20 Schol.).

For \textsc{Spinoza} the transition to clear knowledge is a special case of displacement of motive.
In and of itself displacement of motive can also lead to one passive feeling's being replaced by
another passive feeling. Thus joy in a thing, or longing for it, can lead to envy of the one who
possesses this thing (III. 32). But \textsc{Spinoza} is undoubtedly right that a displacement is
possible from joy in a thing to understanding of this joy and its cause, so that both the feeling
and its cause (object) are thereby ordered in as links in a whole connection of nature, whereby the
being ensnared and captivated by the single thing ceases. Clear understanding is \textsc{Spinoza}'s
great universal remedy, and in this fifth book he can lay the main weight on this remedy all the
more since here --- after the fluctuations shown above in the third and fourth books --- he
decidedly returns to his intellectualistic view, according to which ``the power of the mind
consists in knowledge alone'' (\textit{mentis potentia sola intelligentia definitur}) (see the end
of the preface and 20 Schol.). Only through clear understanding does the soul become wholly and
fully itself.

\textit{β.} But the development takes time. True knowledge is reached only step by step. What
matters is to use the intervals between strong commotions of the mind for reflection.
\textsc{Spinoza} speaks here from his own experience. In the period of ferment he describes at the
beginning of the treatise on the improvement of the understanding, he felt freer when only he could
give himself up to serious reflection (\textit{modo possem serio deliberare}!). And the intervals in
which this happened gradually became more frequent and longer. Therefore he now exhorts us (V, 10
and 20 Schol.) to use the intervals to impress on ourselves certain principles of life
(\textit{vitæ dogmata}) that may be needed when passive feelings again threaten to overwhelm us.
Such principles of life are, for example, that hatred is not overcome by hatred but by magnanimity,
and that men's conduct, like everything in nature, is subject to necessary laws. By impressing such
thoughts on oneself one can make headway, even if it does not happen without fluctuations
(\textit{sine animi fluctuatione}). On the other hand it is an unfortunate path one takes when, for
example, one consoles oneself for not being rich by reflecting on how often wealth is misused and
how vicious the rich often are. One thereby only causes new sorrow in oneself. No, the guiding
motive must be the urge for spiritual freedom (\textit{amor libertatis}). One is then led to
understand the origin of valuable traits of character, and the mind will be filled with the joy of
knowledge.

\textit{γ.} This joy of knowledge reaches its height when the understanding of our feelings and
their causes goes back to the ultimate foundation of our knowledge, \opage{89}the thought of God
(substance, nature). In this foundation the joy of knowledge has its ultimate cause, and since all
joy becomes love when it is combined with the idea of its cause (III, 53 and V, 15 Dem.), love of
God will arise, --- a love that does not demand love in return and gives no occasion for jealousy
or envy, but that can occupy the mind in the highest degree as long as life lasts (14---20).

\subsection*{c. Propositions 21---40.}
\begin{center}\textit{The Eternity of the Mind.}\end{center}

\textit{α.} In this section \textsc{Spinoza} again takes his starting point in the physical
(physiological) side of man's being. He distinguishes between the body in its special or actual
existence, with a definite time of existence and duration in time, and the nature of the body as
contained in the eternal order of nature, which is expressed in the constant relation between rest
and motion in the universe; for him the universe is a single great individual in which no change
takes place, however great the changes that go on in its single parts. (See above p. 28---30; 43).
In one of his letters he says that our bodies were not created but existed before our birth,
though in other forms (\textit{qvamvis alio modo formata}) (Ep. 4, cf. the preface to the second
part of the Short Treatise and L. M.'s preface to \textsc{Spinoza}'s exposition of
\textsc{Descartes}' philosophy). In so far as the nature of the single body is contained in the
material order of nature, and through it in the attribute of extension, it is an eternal mode, an
object of knowledge ``under the aspect of eternity''. In a corresponding way the soul too is
contained in the spiritual order of nature, and through it in the attribute of thought, as an
eternal mode, an object of knowledge ``under the aspect of eternity''. Only so long as the definite
existence in time lasts does the soul have sensation and memory. There is in God (substance,
nature) a thought that corresponds to the existence of the body in all its forms (before and after
birth and after death)\footnote{``Our body presented before birth another relation between rest and
motion [than after], and will likewise after death subsist under another relation between rest and
motion [than before death]. But nonetheless there was then [i.e. before birth], and there will be
afterwards [after death], just as well as now [in the lifetime], a thought corresponding to the
body in substance as thinking, only not the same one, since the body will be placed in another
relation between motion and rest''. \textit{Short Treatise}. II. Preface (Note 10).}. Therefore the
soul, no more than the body, becomes nothing at death (21---23).

This is grounded by means of II, 8 Cor., where already (see above p. 41 f) a distinction was made
between a double form of existence for all single things (\textit{res singulares}), --- the one in so
far as things are contained in God's attributes, the other in so far as they also have existence in
time. This proposition \textsc{Spinoza} now takes up, not only in the proof of 21 and 23, but more
definitely in 29 Schol. (where reference is made back to II, 45, which in turn is grounded on II, 8
Cor.). The first-named form of existence is eternal, just as the relation between ground and
consequent is. For there has never been a time when the ground (the premises) was, but the
consequent (the conclusion) was not. When two chords \opage{90}intersect within a circle, the
rectangles formed by the segments cut off are equal. That is an eternal truth. In each single case
only certain definite chords are drawn; but the others exist nonetheless as contained in the
concept of the circle (II, 8 Cor.). It is here again \textsc{Spinoza}'s identification of ground and
cause that makes itself felt. As determined in time, man is a link in a causal series in which
every link is at once the effect of the preceding and the cause of the following; but in his
eternal being he is an ungenerated, unchangeable and imperishable mode of substance. The opposition
between these two forms of existence would stand out very glaringly, were it not the case for
\textsc{Spinoza} that the difference of time is altogether inessential to the causal relation, while
what is decisive is the timeless relation between ground and consequent.

\textsc{Spinoza} recalls \textsc{Plato} here. The Neoplatonic train of thought had indeed also
influenced him through Jewish and scholastic theology. Just as for \textsc{Plato} the world of
ideas expresses the content of existence as eternal, while the things observation shows us are
finite, imperfect and time-bound forms of this eternal content, so for \textsc{Spinoza} the eternal
relation between ground and consequent is the true reality, while the phenomena given in experience
are time-bound formations (\textit{affectiones}) of this eternal content. The difference from
\textsc{Plato} consists, as already emphasized above (see p. 48 f), in \textsc{Spinoza}'s putting
(as \textsc{Bacon} and \textsc{Bruno} had already done) the concept of law in place of the concept of
idea.

\textit{β.} On the question of the relation between the part of man's nature that belongs to time
and the part that belongs to eternity, \textsc{Spinoza} relies (24---31) on his theory of knowledge.
The more we understand single things, the more we understand God, who is the ultimate
presupposition of their being and their continuance. We then build not on loose and chance
observations, but on the lawful connection between single things, that is, on the foundation of
scientific understanding, in which knowledge of the individual and insight into the general laws
become one. (See above p. 51---54). Thereby not only is the innermost urge of the human mind
satisfied, but the knowledge by which man comes to partake of God's being is possible only by means
of what is eternal in man's being itself. The more man is absorbed in perfect knowledge, the more
he can know that he is in God and is understood from God's own being, which is indeed the eternal
relation between ground and consequent (30). To a knowledge of the kind named there corresponds,
within the attribute of extension, not the body's time-bound existence, but the body's being as an
eternal mode, seen under the aspect of eternity (29 and 31).

For \textsc{Spinoza} the ideal of knowledge coincides with the ideal of existence, that is, with
substance. But the highest kind of knowledge can be developed in very different degrees (26), and
therefore the relation between the eternal and the perishable part of the soul can also be very
different (38).

When intellectualism reaches its height in the interpretation of the highest human knowledge as
God's indwelling in man, there nevertheless contributes to this a mystical tendency that never
disappeared in \textsc{Spinoza}, even after he had broken \opage{91}with his theological past, but
that comes out more strongly in the fifth book than in the earlier books. It is an immediate feeling
of unity that presses forward. His religious urge will not acknowledge any difference between
thought itself and the highest, concluding object of thought. They are to become wholly one. Man's
thinking of God is God's thinking of man. --- If we compare the fifth book with the Short Treatise,
mysticism has acquired a more intellectualistic stamp in the course of \textsc{Spinoza}'s
development. But his intellectualism shows itself not in passive contemplations but in the activity
of thought. It is by becoming wholly self-active in his knowledge, by living and acting from
himself, that man becomes one with substance, becomes as it were one of the rays that go forth from
it as from a source of light. God is not an object for thought, so that thought should stand in a
relation of dependence on God. In the activity of thought itself God stirs. (Cf. above p. 7---9).
\textsc{Spinoza}'s mysticism is here a consequence of his original presupposition about the
relation between knowledge and its content.

\textsc{Spinoza} does not investigate whether there might not be other forms of psychical energy,
and thereby of spiritual culture, than the purely intellectual. When \textsc{Goethe} found the
foundation for his belief in immortality in the thought that our spirit is of an ever-active
nature, and that if he was restlessly active until his death, nature was obliged to provide a new
form of existence for him when the present one could no longer hold out, --- he was thinking not of
the energy of thought alone, but also of other forms in which the human spirit can act. The thinker
and the poet agree, however, that only those men are immortal in whom a great energy has developed:
``Wir sind nicht auf gleiche Weise unsterblich, und um sich künftig als grosse Entelechie zu
manifestiren, muss man auch eine sein''. (Gespräche mit \textsc{Eckermann}. 1 Sept. 1829).
\textsc{Spinoza} says the same in his language, when he makes immortality conditional on the
attainment of the highest form of knowledge and its dominion in the soul, and when he acknowledges
different degrees in these respects (38---39). Starting from the highest knowledge as the
expression of eternity, he could not arrive at immortality for all\footnote{Cf. \textit{Den nyere
Filosofis Historie\textsuperscript{2}}, I, p. 552 (Note 67), where a series of authors from antiquity
and modern times is named in whom the thought appears that only those are immortal who have reached
the highest spiritual development. To the names given there must be added \textsc{Renouvier}, and
perhaps \textsc{Sibbern}. --- This assumption occurs among theologians too. Thus \textsc{Otto
Møller} interprets ``eternal damnation'' as annihilation (Brevveksling med \textsc{Rørdam}, II, p.
321).}.

\textsc{Spinoza} nevertheless thinks that all men are conscious of immortality, only they confuse it
with duration or length of time and believe that after death there can be talk of memory and of
pictorial ideas (34 Schol.). Strictly speaking, one can speak neither of a beginning nor of a ceasing
nor of a continuation of the soul's true being (31 Schol. and 34 Schol.). The immortality
\textsc{Spinoza} speaks of lies not in the continuation of existence in time, but is in every moment
and has always been. It comes to our consciousness when our activity of thought stirs in a
decidedly active way. ``We feel and experience'', says \textsc{Spinoza}, ``that we are eternal, for
\opage{92}we feel no less what we conceive in understanding (\textit{res, qvas mens intelligendo
concipit}) than what we remember. The eyes of the soul are the proofs of thought themselves''. The
external difference between thought and object falls away in true knowledge. Eternity stirs in
thought itself, not merely in its results.

There are two questions that press forward in face of this whole train of thought in
\textsc{Spinoza}. The one is whether a man's being, if it is to be wholly understood within the
connection that knowledge points out, must not also include its relation to (its being conditioned
by, and its intervention in) the definite circumstances of time and place under which it exists.
This relation stamps and conditions its innermost nature and cannot be indifferent to the full
understanding --- even from the standpoint of eternity, in so far as it is possible for us to take
that standpoint, which according to \textsc{Spinoza} himself is possible only very rarely. This
first question concerns the understanding of existence in time. The second question concerns its
value. For according to \textsc{Spinoza} a long theoretical and practical course of development
(see above p. 89---90) is required to reach the stage where eternal existence first becomes
self-evident. But then eternity becomes dependent on time, not merely the other way round.
\textsc{Spinoza} himself would surely have been very reluctant to admit that the whole spiritual
work he has described so thoroughly should be an altogether indifferent episode. When the real
validity of the form of time for our knowledge is not acknowledged, the reality of spiritual work
cannot be maintained either. Indeed, the very concept of value presupposes the reality of time,
since it presupposes urge and striving.

This presupposition, now, \textsc{Spinoza} would not acknowledge. He demands, after all, that the
perfection of things be judged by their ``nature and power'', and not by our urge. But the concept
of value has not therefore disappeared from \textsc{Spinoza}'s philosophy; it has, as shown above
(p. 21 f. 33), only come to rest in the fundamental thought of the innermost law of existence. And
it is in relation to cosmological value that all other value in the end becomes vanishing.

\textit{γ.} Man comes to partake of this value through the joy of knowledge. Since God, as the
ultimate ground of the possibility of knowledge, is the cause of the joy of knowledge, there arises
an intellectual love of God (\textit{amor intellectualis dei}). It is not directed toward any object
in time, but toward an eternal being that only thought grasps (32 Cor.). Properly speaking there
can be no talk of any object of this love. For just as God is and acts in human knowledge itself at
its highest stage, so the intellectual love of God is a part of God's love of himself, as the one
who acts in man's being when this has unfolded into an eternal mode of thought. The intellectual
feeling is therefore an active feeling in the highest degree (36). In relation to it every other
feeling pales.

From \textsc{Spinoza}'s own psychology there follows a difficulty with regard to the concept of
``intellectual love'', if this is to be a feeling that arises and persists according to
psychological laws. It is to be eternal and unchangeable like substance, in whose being man comes to
partake through the highest degree of knowledge. But according to \textsc{Spinoza}'s psychology (see
above p. 63 f) every feeling presupposes a change of state \opage{93}and consists in a transition
--- something he himself recalls in this connection (33 Schol.), in the same moment as he declares
that intellectual love is not joy (\textit{lætitia}), which consists in the transition to
perfection, but blessedness (\textit{beatitudo}), which consists in the possession of perfection.

At another point too, at the height of his mysticism, \textsc{Spinoza} is compelled to alter
psychological propositions he had set up earlier. In the axiom at the beginning of the fourth book
it was maintained that there was not in existence (\textit{in rerum natura}) any individual being
that could not come to face something stronger. But now he teaches (37) that there is nothing in
nature (\textit{in natura}) that is stronger than intellectual love and could annul it. He has felt
the contradiction himself, but maintains (37 Schol.) that the axiom in the fourth book concerns only
individual beings as existing at a definite time and place, not eternal modes that are one with
substance. But --- a psychological development to clear knowledge and to the joy of knowledge does
after all always take place under definite conditions of place and time. And \textsc{Spinoza}
himself teaches that intellectual love can only gradually permeate the whole soul, and that a great
work is needed to bring it about that ``the greatest part of the soul'' is ruled by it (38---39).
Intellectual love thus does appear in the psychological arena, and wherein lies the certainty that
it will always conquer? One simply cannot escape psychology. It is omnipresent in the sphere of
spiritual life, right up into the highest intellectual, ethical and religious regions.

Later thinkers too have sought to free themselves from psychology as regards these ideal regions.
\textsc{Kant} rejected every motive of feeling in the ethical sphere, and when he must nevertheless
admit that the moral law can become a power within us only through the respect and reverence it
inspires, he does not shrink from the consequence of declaring that respect and reverence are not
feelings at all, since they have nothing to do with pleasure and pain, indeed cannot be explained
by way of experience at all\footnote{See \textit{Den nyere Filosofi's Historie\textsuperscript{2}},
II, p. 83 f. --- \textsc{Kant}'s ``respect'' belongs to what \textsc{Spinoza} calls ``active
feelings''.}. A distinction like the one \textsc{Spinoza} makes between blessedness and joy is made
by \textsc{Carlyle} between blessedness and happiness\footnote{\textit{Sartor Resartus}, Book II,
Chap. 9: Man can do without happiness, and instead hereof find blessedness. --- Cf. my remarks on
the terminology, \textit{Etik} III, 11 and VII, 1.}.

It would be wrong to see a mere dispute about words here, when eminent thinkers refuse to bring
their highest ideals under psychological illumination. When continuity seems to them broken, it is
because psychology has not succeeded in showing lawful and motivated transitions from the lowest to
the highest forms of spiritual life. \textsc{Spinoza} himself (in the third and fourth books) has
given valuable contributions in this direction, but in spite of all that has been accomplished
since, there are still always tasks here for descriptive and analytical psychology. In distinctions
such as those \textsc{Spinoza}, \textsc{Kant} and \textsc{Carlyle} maintain, there are contained
tasks for psychology that still demand to be treated.

\opage{94}\subsection*{d. Propositions 41---42.}
\begin{center}\textit{Ethics Is Independent of Speculation and Dogma.}\end{center}

\textsc{Spinoza}'s last word, however, is not speculation and mysticism. Proposition 41, the last
but one in the whole work, runs thus: ``Even if we did not know that our mind is eternal, we would
still regard the sense of duty, religion, and everything that in the fourth book is reckoned to
courage and generosity, as the highest''. And in his proof of this proposition he points out that
in grounding the value of these traits of character he has made no use of his doctrine of
immortality. Most men, he adds, think otherwise, to be sure; they would not esteem those traits of
character if there were nothing to hope or to fear after death. In one of his letters (Ep. 43) he
maintains firmly that whether the ethical propositions (\textit{documenta moralia}) are God's
commandments or not, they are equally divine and salutary. And already in the Short Treatise (II,
26, 4) he mocks at the fact that many who are regarded as great theologians say that if there were
no eternal life they would seek only their own advantage. It is as if a fish were to say that,
since no other life follows its life in the water, it would leave its own element and go up on
land!

The last proposition of the Ethica (42) states that blessedness (\textit{beatitudo}) is not the
reward of valuable traits of character but one with them. We do not master our passive feelings in
order to attain blessedness; but conversely --- we can master the passive feelings because we
possess blessedness. For the highest life of thought, and the joy bound up with it, are able to
secure us against becoming the slave of passive feelings. It shows itself again in this closing
consideration how deeply \textsc{Spinoza} is convinced of the power of thought to harmonize life.

\opage{95}\section*{Concluding Characterization.}

\textit{a.} The innermost nerve of \textsc{Spinoza}'s thinking is his conviction of the rationality
of existence. The primal phenomenon\footnote{On this concept, which comes from \textsc{Goethe}, and
its significance for every attempt at a view of the world, see \textit{Den menneskelige Tanke}, p.
311 f.} on which his system is built was this, that the things or events in the world could be
conceived as links in logical and mathematical series, in which every link is related to the others
as ground to consequent or as consequent to ground. By two ways, according to the conjecture put
forward in the present treatise, the conviction of the reality of this primal phenomenon became
fixed in him. In the concept of God to which he was led from his childhood he found, when it was
thought through consistently, the principle of a rational ground for everything expressed. And he
found the same principle as the presupposition of the new natural science's striving for a
mathematical formulation of physical laws. The concept of God and the concept of nature met for him
in the concept of substance, defined as that whose concept does not in turn presuppose any other
concept. The series of grounds and consequents (causes and effects) must, if existence is absolutely
rational, rest on a principle that in turn presupposes nothing, but ``has its ground in itself'' or
is ``cause of itself''.

It was moreover his conviction that all spiritual life not only presupposes thought, but at its
height is a life of thought. All spiritual needs are then satisfied when perfect knowledge is
reached. All culture is for him in the end intellectual. All value is in the end contained in the
joy of knowledge.

All science and all philosophy work, and must work, in the direction \textsc{Spinoza} points out. It
is always a matter of finding the inner connection between the different phenomena of experience
(modes) and between the different sides or properties of existence (attributes). To that extent
\textsc{Hegel} was right when he said: ``Entweder hast du den Spinozismus, oder du hast keine
Philosophie''.

\textit{b.} About 60 years ago (1857) there appeared here at home a thorough and acute monograph on
\textsc{Spinoza} by \textsc{Hans Bröchner}. In much my view of the great thinker agrees with
\textsc{Bröchner}'s. It was indeed \textsc{Bröchner} who first really taught me to value
\textsc{Spinoza} and his thinking, even if \opage{96}many years passed before I came to the view I
have put forward in the present work. And the memory of my teacher and friend, who in character
and conduct of life had so much in common with \textsc{Spinoza}, has accompanied me during my
renewed immersion in the subject.

When \textsc{Bröchner}'s work appeared, \textsc{van Vloten} had not yet made the important find that
brought ``the Short Treatise'' and the original text of the letters to light, and thereby not only
corrected various biographical particulars but also threw light on \textsc{Spinoza}'s course of
development. Not only in that I, like the many authors who in the last 50 years have occupied
themselves with \textsc{Spinoza}, have been able to use these documents, but also on several single
points my view differs from \textsc{Bröchner}'s, and for that matter from that of the Spinoza
scholars I know. I have especially sought to understand his ideas in connection with fundamental
thoughts present in the natural science of the time, and have thereby come to stress the realistic
character of his philosophy, which was wholly misjudged in the Romantic period, however much people
enthused over \textsc{Spinoza}. It is this realistic character of his thinking that has aroused so
much interest in him in the last part of the nineteenth century. \textsc{Spinoza} had the quiet
conviction that in order to hold fast to the reality of spiritual life in the world it is by no
means necessary to set limits in advance to the work of empirical science, either in the sphere of
natural science or in the historical sphere. \textsc{Bröchner} conceives \textsc{Spinoza} ``as the
representative of the starting point of modern idealism'', as the initiator of a movement that was
to find its conclusion in \textsc{Hegel}. On my view, romantic idealism became possible precisely
only through a break with the direction \textsc{Spinoza} had taken. \textsc{Spinoza} is no
idealist, if that involves a denial in principle of the points of view and methods of empirical
science.

But the most essential difference between my view of \textsc{Spinoza} and \textsc{Bröchner}'s will
appear from some remarks with which I shall close this treatise.

\textit{c.} Can we really, as \textsc{Spinoza} presupposes, conclude our science with the thought of
something that is understood through itself and exists through itself? --- Thinking consists in
synthesis and comparison, the former as the presupposition of the latter. But every synthesis
brings with it a limitation that can always be annulled by new experience, and every comparison
presupposes a point of view that can itself become the object of comparison with other points of
view. We can end with a question, but not with a thought that concludes once and for all. What for % the Danish doubles en (en en Gang)
\textsc{Spinoza} was the primal phenomenon, that which had its ground in itself, was the lawfulness
in the world. But this lawfulness is a great hypothesis, which science in its constant work seeks to
find confirmed through ever new application. And (as \textsc{Spinoza} himself has admitted) the
single phenomena and events cannot be derived from the very concept of lawfulness, even if they can
be understood (and can only be understood) by applying this concept.

In his precritical period \textsc{Kant} pursued a line of thought that recalls \textsc{Spinoza}'s.
He inferred from the lawful connection in nature an absolute \opage{97}being that lay at the ground
of everything. In this period it was for him ``the only possible ground of proof'' for the
objective validity of the concept of God. But the transition to the critical philosophy was made
when \textsc{Kant} no longer asked about the objective ground of the world, but about the
possibility of an objective concept of the world, and when he found this possibility in the very
fundamental form of human thinking, whose validity rests on this, that without strict logical
connection of observations experience in the strict sense is not possible\footnote{Cf. Om
\textit{Kontinuiteten i Kant's filosofiske Udviklingsgang} (Vid. Selsk. Skr. 1893), p. 14---15.}.
The active character of thinking, which \textsc{Spinoza} had so strongly emphasized, was now made
the basis of the whole consideration. The concept of synthesis replaced the concept of substance,
and while \textsc{Spinoza}'s philosophy was, or meant to be, first and foremost cosmology
(metaphysics), a doctrine of the innermost nature of existence, the problem of knowledge now pressed
forward ahead of all other philosophical problems, and a scientific cosmology came to stand as a
limiting problem. To \textsc{Spinoza}'s substance (``that which is in itself'') there corresponds
in \textsc{Kant}'s philosophy ``the thing in itself'', of which nothing can be known (not even that
it ``is understood through itself'').

A closer discussion of the problem of knowledge led in turn to distinguishing between different
groups of categories, corresponding to \textsc{Spinoza}'s distinction between substance, attributes
and infinite modes of different degrees, but without the immediate metaphysical significance he % the Danish misprints Belydning (Betydning)
attributed to them. A distinction is made between ground and cause, that is, between formal and
real categories, and the concepts of value appear in their difference from both these groups. From
this it follows in turn that all spiritual culture cannot, as \textsc{Spinoza} at bottom
presupposes, be one with intellectual culture.

But the ideal conception of the work of thought that \textsc{Spinoza} maintained, and the inward joy
of knowledge that for him shed luster over life, cannot be lacking in any true inquirer. Every % the Danish misprints kam (ham)
observation and every thought that gives an insight into the great connection that all
understanding of existence presupposes lets us feel the significance of the words of
\textsc{Goethe}, that faithful devotee of Spinoza: ``Ort für Ort sind wir im Inneren'', although we
cannot, as \textsc{Spinoza} thought, unveil this inner once and for all.

% ===== p. [99], Indhold =====
\clearpage
\begin{center}{\large CONTENTS}\end{center}
\medskip
\hfill{\scriptsize Page}\par
\tocline{Introduction}{1---13}
\tocline{First Book}{14---35}
\tocline{Second Book}{36---57}
\tocline{Third Book}{58---69}
\tocline{Fourth Book}{70---83}
\tocline{Fifth Book}{84---94}
\tocline{Concluding Characterization}{95---97}

\end{document}
